% Traduction et production d'articles de journaux scolaires et de textes de presse — Allemand 1ère A — Cours signé OnBuch+
% Programme officiel MINESEC. Compiler avec : tectonic ALA20-traduction-et-production-d-articles-de-journaux-scolaires-et.tex
\newcommand{\DOCMATIERE}{Allemand}
\newcommand{\DOCNIVEAU}{1ère A}
\newcommand{\DOCMODULE}{Chapitre 4 : Média et communication}
\newcommand{\DOCLECON}{ALA20}
\newcommand{\DOCTITRE}{Traduction et production d'articles de journaux scolaires et de textes de presse}
% OnBuch+ — préambule « cours signés » ENRICHI (compilé avec tectonic / XeLaTeX)
% Système visuel « Pop » d'OnBuch+ : fond crème (jamais blanc), bordures encre
% épaisses, ombres « dures » décalées, titres Archivo Black en capitales.
% Version MATHÉMATIQUES : graphes (pgfplots/tikz), boîtes pédagogiques
% (définition, propriété, méthode, attention, le savais-tu, exercice résolu,
% exercice, TP, bilan), page de garde et signature OnBuch+ ancrée en bas.
%
% Macros attendues dans main.tex AVANT \input{preamble} :
%   \DOCMATIERE  \DOCNIVEAU  \DOCMODULE  \DOCLECON (numéro)  \DOCTITRE
\documentclass[11pt]{article}

\usepackage{fontspec}
\usepackage[ngerman,french]{babel} % français = langue principale ; l'allemand via \all{..} / textebox
\usepackage[a4paper,margin=2cm,top=3.4cm,bottom=2.8cm,headheight=1.4cm,footskip=1.3cm]{geometry}
\usepackage{amsmath,amssymb,amsfonts}
\usepackage{mathtools} % \xmapsto, \coloneqq... souvent utilisés par les modèles
\usepackage{cancel} % simplifications barrées (bilans de forces)
% \abs{z} (module, valeur absolue) et \norm{u} (norme), utilisés par les modèles
\providecommand{\abs}{}\let\abs\relax\DeclarePairedDelimiter{\abs}{\lvert}{\rvert}
\providecommand{\norm}{}\let\norm\relax\DeclarePairedDelimiter{\norm}{\lVert}{\rVert}
\usepackage[table]{xcolor} % [table] : fournit aussi \rowcolors (alternance de lignes)
\usepackage{pagecolor}
\usepackage{tikz}
\usetikzlibrary{arrows.meta,positioning,calc,shapes.geometric,shapes.misc,shapes.arrows,decorations.pathmorphing,decorations.pathreplacing,decorations.markings,patterns,shadows,mindmap,trees,fit,backgrounds,matrix,intersections,angles,quotes}
\usepackage{pgfplots}
\usepackage[version=4]{mhchem} % équations nucléaires \ce{^{235}_{92}U}
\usepackage[european,siunitx]{circuitikz} % schémas électriques
\pgfplotsset{compat=1.18}
\usepgfplotslibrary{fillbetween,groupplots} % aires entre courbes (calcul intégral)
\usepackage{enumitem}
\usepackage{fancyhdr}
% [most] charge toutes les bibliothèques tcolorbox (listings, minted, raster,
% documentation...) et gonfle inutilement la mémoire TeX sur un document long
% (~300 boîtes) : on ne charge que ce qui est réellement utilisé.
\usepackage[skins,breakable]{tcolorbox}
\usepackage{graphicx}
\usepackage{colortbl}
\usepackage{booktabs}
\usepackage{tabularx}
\usepackage{diagbox} % case d'en-tête diagonale des tableaux à double entrée
% Colonnes extensibles centrée / gauche / droite (C, L, R), souvent utilisées par les modèles
\newcolumntype{C}{>{\centering\arraybackslash}X}
\newcolumntype{L}{>{\raggedright\arraybackslash}X}
\newcolumntype{R}{>{\raggedleft\arraybackslash}X}
\usepackage{array}
\usepackage{multicol}
\usepackage{multirow}
\usepackage{siunitx}
% parse-numbers=false : accepte aussi \SI{2,5 \times 10^{-3}}{...}
\sisetup{locale=FR,output-decimal-marker={,},group-minimum-digits=4,parse-numbers=false}
\DeclareSIUnit\ohm{\ensuremath{\symup{\Omega}}} % U+2126 absent de la police
\DeclareSIUnit\division{div} % réglages d'oscilloscope (ms/div, V/div)
\DeclareSIUnit\tr{tr} % tours (vitesse de rotation, ex. \SI{1500}{\tr\per\minute})
\DeclareSIUnit\molar{M} % concentration molaire (ex. \SI{3}{\molar}, \SI{1}{\micro\molar})
\DeclareSIUnit\an{an} % année (le modèle confond parfois avec \year, un compteur TeX) — ex. \SI{5}{\kilo\meter\per\an}
\DeclareSIUnit\FCFA{FCFA} % franc CFA (ex. \SI{500}{\FCFA\per\kilo\gram})
\DeclareSIUnit\degreeBrix{\textdegree Brix} % teneur en sucre d'un sirop/jus (réfractométrie)
\DeclareSIUnit\month{mois} % le modèle utilise parfois \month, un compteur TeX, comme unité de durée
\usepackage{qrcode}
% \degree : commande fréquemment utilisée par les modèles pour un angle ou une
% température en dehors de \SI{}{} (non fournie par les paquets chargés).
\providecommand{\degree}{^{\circ}}
% \qed (fin de démonstration), utilisé par les modèles sans amsthm
\providecommand{\qed}{\hfill$\square$}
% \barre : double barre des valeurs interdites dans un tableau de variations (tkz-tab)
\providecommand{\barre}{\ensuremath{\|}}
% environnement proof (amsthm non chargé) utilisé par les modèles
\ifdefined\proof\else\newenvironment{proof}[1][Démonstration]{\par\noindent\textit{#1.}\ }{\qed\par}\fi
\usepackage{lastpage}
\usepackage{hyperref}
\hypersetup{hidelinks}

% ============================================================
% Palette « Pop » OnBuch+ — mêmes hex que la classe P (plus_ui.dart)
% ============================================================
\definecolor{poporange}{HTML}{F59321}
\definecolor{popdark}{HTML}{E07A0C}
\definecolor{popcream}{HTML}{FFF6E8}
\definecolor{popink}{HTML}{1C1714}
\definecolor{poppurple}{HTML}{7A5AE0}
\definecolor{popgreen}{HTML}{2E9E6A}
\definecolor{poppink}{HTML}{E0457B}
\definecolor{popblue}{HTML}{2F7DE1}
\definecolor{popgold}{HTML}{C98A12}
\definecolor{popmuted}{HTML}{8A7F76}
\definecolor{popline}{HTML}{EADFD2}
\definecolor{popgray}{HTML}{8A7F76}
\colorlet{popgrey}{popgray}
% Alias tolérants (le modèle nomme parfois les couleurs différemment)
\colorlet{popwhite}{white}
\colorlet{popblack}{popink}
\colorlet{popred}{poppink}
\colorlet{popyellow}{popgold}
% Teintes claires (fonds de tableaux/encadrés) — le modèle nomme parfois
% les couleurs avec un suffixe L (ex. popblueL) sans les avoir définies.
\colorlet{poporangeL}{poporange!20}
\colorlet{popdarkL}{popdark!20}
\colorlet{poppurpleL}{poppurple!20}
\colorlet{popgreenL}{popgreen!20}
\colorlet{poppinkL}{poppink!20}
\colorlet{popblueL}{popblue!20}
\colorlet{popgoldL}{popgold!20}
\colorlet{popredL}{popred!20}
\colorlet{popgrayL}{popgray!20}
% Teintes claires pour fonds de tableaux / zones
\colorlet{popblueL}{popblue!10!white}
\colorlet{poporangeL}{poporange!14!white}
\colorlet{popgreenL}{popgreen!12!white}
\colorlet{poppurpleL}{poppurple!12!white}
\colorlet{poppinkL}{poppink!12!white}
\colorlet{popgoldL}{popgold!14!white}

\pagecolor{popcream}

% ============================================================
% Typographie
% ============================================================
\defaultfontfeatures{Ligatures=TeX}
\setmainfont{PlusJakartaSans-Medium.ttf}[Path=../../../fonts/,BoldFont=PlusJakartaSans-Bold.ttf]
% Symboles, API et emojis absents de la police principale : repli sur DejaVu Sans (caractères actifs)
\newfontface\symfont{DejaVuSans.ttf}[Path=../../../fonts/]
\makeatletter
\newcommand\symactive[1]{\catcode#1=13 \begingroup\lccode`\~=#1 \lowercase{\endgroup\def~}{{\symfont\char#1}}}
\newcommand\symblank[1]{\catcode#1=13 \begingroup\lccode`\~=#1 \lowercase{\endgroup\def~}{}}
\newcommand\symhyph[1]{\catcode#1=13 \begingroup\lccode`\~=#1 \lowercase{\endgroup\def~}{\mbox{-}}}
\newcount\symc
\newcommand\symrange[2]{\symc=#1 \@whilenum\symc<#2 \do{\expandafter\symactive\expandafter{\the\symc}\advance\symc by 1 }}
\newcommand\symblankrange[2]{\symc=#1 \@whilenum\symc<#2 \do{\expandafter\symblank\expandafter{\the\symc}\advance\symc by 1 }}
\makeatother
\symrange{"0250}{"02FF}\symactive{"2713}\symactive{"2714}\symactive{"2717}\symactive{"2718}\symactive{"2610}\symactive{"2611}\symactive{"2612}
\symhyph{"2011}\symhyph{"2E1A}\symblankrange{"1F300}{"1FAFF}\symblank{"FE0F}
% Maths en sans-serif (Fira Math) avec chiffres et lettres droites de Plus
% Jakarta, pour que les formules se fondent dans le texte.
\usepackage[math-style=ISO,bold-style=ISO]{unicode-math}
\setmathfont{FiraMath-Regular.otf}[Path=../../../fonts/]
\setmathfont{PlusJakartaSans-Medium.ttf}[Path=../../../fonts/,range={up/num,up/{Latin,latin},\mathup}]
\usepackage{icomma} % virgule décimale sans espace en mode math
% Caractères mathématiques Unicode absents des polices de texte (Plus Jakarta,
% Archivo Black) : rendus par leur équivalent mathématique, en texte comme en
% formule — sinon ils s'affichent en carré vide (ex. « Arithmétique dans ℤ »).
\usepackage{newunicodechar}
\newunicodechar{ℝ}{\ensuremath{\mathbb{R}}}
\newunicodechar{ℤ}{\ensuremath{\mathbb{Z}}}
\newunicodechar{ℂ}{\ensuremath{\mathbb{C}}}
\newunicodechar{ℕ}{\ensuremath{\mathbb{N}}}
\newunicodechar{ℚ}{\ensuremath{\mathbb{Q}}}
\newunicodechar{𝔻}{\ensuremath{\mathbb{D}}}
\newunicodechar{α}{\ensuremath{\alpha}}
\newunicodechar{β}{\ensuremath{\beta}}
\newunicodechar{γ}{\ensuremath{\gamma}}
\newunicodechar{δ}{\ensuremath{\delta}}
\newunicodechar{ε}{\ensuremath{\varepsilon}}
\newunicodechar{θ}{\ensuremath{\theta}}
\newunicodechar{λ}{\ensuremath{\lambda}}
\newunicodechar{μ}{\ensuremath{\mu}}
\newunicodechar{σ}{\ensuremath{\sigma}}
\newunicodechar{τ}{\ensuremath{\tau}}
\newunicodechar{φ}{\ensuremath{\varphi}}
\newunicodechar{ψ}{\ensuremath{\psi}}
\newunicodechar{ω}{\ensuremath{\omega}}
\newunicodechar{Σ}{\ensuremath{\Sigma}}
% majuscules produites par \MakeUppercase dans les titres (α-aminés -> Α)
\newunicodechar{Α}{\ensuremath{\alpha}}
\newunicodechar{Β}{\ensuremath{\beta}}
\newunicodechar{Γ}{\ensuremath{\Gamma}}
\newunicodechar{Δ}{\ensuremath{\Delta}}
\newunicodechar{Ε}{\ensuremath{\varepsilon}}
\newunicodechar{Θ}{\ensuremath{\Theta}}
\newunicodechar{Λ}{\ensuremath{\Lambda}}
\newunicodechar{Μ}{\ensuremath{\mu}}
\newunicodechar{Τ}{\ensuremath{\tau}}
\newunicodechar{Φ}{\ensuremath{\Phi}}
\newunicodechar{Ψ}{\ensuremath{\Psi}}
\newunicodechar{Ω}{\ensuremath{\Omega}}
\newunicodechar{↦}{\ensuremath{\mapsto}}
\newunicodechar{∈}{\ensuremath{\in}}
\newunicodechar{∉}{\ensuremath{\notin}}
\newunicodechar{∧}{\ensuremath{\wedge}}
\newunicodechar{⇔}{\ensuremath{\Leftrightarrow}}
\newunicodechar{⟺}{\ensuremath{\Longleftrightarrow}}
\newunicodechar{⇒}{\ensuremath{\Rightarrow}}
\newunicodechar{⟹}{\ensuremath{\Longrightarrow}}
\newunicodechar{∘}{\ensuremath{\circ}}
\newunicodechar{∪}{\ensuremath{\cup}}
\newunicodechar{∩}{\ensuremath{\cap}}
\newunicodechar{≡}{\ensuremath{\equiv}}
\newunicodechar{⊂}{\ensuremath{\subset}}
\newunicodechar{⊥}{\ensuremath{\perp}}
\newunicodechar{∃}{\ensuremath{\exists}}
\newunicodechar{∀}{\ensuremath{\forall}}
\newunicodechar{∛}{\ensuremath{\sqrt[3]{\,}}}
\newunicodechar{∜}{\ensuremath{\sqrt[4]{\,}}}
\newunicodechar{⃗}{\ensuremath{{}^{\to}}}
\newunicodechar{ⁿ}{\ifmmode^{n}\else\textsuperscript{\ensuremath{n}}\fi}
\newunicodechar{ˣ}{\ifmmode^{x}\else\textsuperscript{\ensuremath{x}}\fi}
\newunicodechar{ᵇ}{\ifmmode^{b}\else\textsuperscript{\ensuremath{b}}\fi}
\newunicodechar{ʳ}{\ifmmode^{r}\else\textsuperscript{\ensuremath{r}}\fi}
\newunicodechar{ᵘ}{\ifmmode^{u}\else\textsuperscript{\ensuremath{u}}\fi}
\newunicodechar{ʸ}{\ifmmode^{y}\else\textsuperscript{\ensuremath{y}}\fi}
\newunicodechar{ᵃ}{\ifmmode^{a}\else\textsuperscript{\ensuremath{a}}\fi}
\newunicodechar{ᵉ}{\ifmmode^{e}\else\textsuperscript{\ensuremath{e}}\fi}
\newunicodechar{ᵏ}{\ifmmode^{k}\else\textsuperscript{\ensuremath{k}}\fi}
\newunicodechar{ᵖ}{\ifmmode^{p}\else\textsuperscript{\ensuremath{p}}\fi}
\newunicodechar{ᴺ}{\ifmmode^{N}\else\textsuperscript{\ensuremath{N}}\fi}
\newunicodechar{ᶜ}{\ifmmode^{c}\else\textsuperscript{\ensuremath{c}}\fi}
\newunicodechar{ʰ}{\ifmmode^{h}\else\textsuperscript{\ensuremath{h}}\fi}
\newunicodechar{ᵠ}{\ifmmode^{q}\else\textsuperscript{\ensuremath{q}}\fi}
\newunicodechar{ₙ}{\ifmmode_{n}\else\textsubscript{\ensuremath{n}}\fi}
\newunicodechar{ₐ}{\ifmmode_{a}\else\textsubscript{\ensuremath{a}}\fi}
\newunicodechar{ᵢ}{\ifmmode_{i}\else\textsubscript{\ensuremath{i}}\fi}
\newunicodechar{ᵣ}{\ifmmode_{r}\else\textsubscript{\ensuremath{r}}\fi}
\newunicodechar{ₖ}{\ifmmode_{k}\else\textsubscript{\ensuremath{k}}\fi}
\newunicodechar{ᵦ}{\ifmmode_{\beta}\else\textsubscript{\ensuremath{\beta}}\fi}
\newunicodechar{ₓ}{\ifmmode_{x}\else\textsubscript{\ensuremath{x}}\fi}
\newfontfamily\popdisplay{ArchivoBlack-Regular.ttf}[Path=../../../fonts/]
\newfontfamily\poplogo{SpaceGrotesk-Bold.ttf}[Path=../../../fonts/]
\newfontfamily\popsemi{PlusJakartaSans-SemiBold.ttf}[Path=../../../fonts/]
\newfontfamily\popxbold{PlusJakartaSans-ExtraBold.ttf}[Path=../../../fonts/]
\newfontfamily\popxboldls{PlusJakartaSans-ExtraBold.ttf}[Path=../../../fonts/,LetterSpace=3.0]

\setlist[enumerate]{leftmargin=1.7em,itemsep=5pt,topsep=4pt}
\setlist[itemize]{leftmargin=1.7em,itemsep=4pt,topsep=4pt,label={\color{poporange}\textbullet}}
\setlength{\parindent}{0pt}
\setlength{\parskip}{5pt}
\renewcommand{\arraystretch}{1.25}

% pgfplots : style Pop par défaut
\pgfplotsset{
  every axis/.append style={axis line style={popink,line width=1pt},tick style={popink},
    grid style={popline,line width=0.5pt},label style={font=\small},tick label style={font=\footnotesize}},
  popcurve/.style={poporange,line width=1.8pt,no markers,smooth},
}
% Même style disponible dans un \draw TikZ simple (hors pgfplots)
\tikzset{popcurve/.style={poporange,line width=1.8pt}}
\tikzset{popgrid/.style={popline,very thin}}
\colorlet{popcurve}{poporange} % le nom du style est parfois utilisé comme couleur

% ============================================================
% Pill / squiggle
% ============================================================
\newcommand{\poppill}[3]{%
  \tikz[baseline=-0.55ex]\node[draw=popink,line width=1.2pt,rounded corners=2.6pt,%
    fill=#2,inner xsep=6.5pt,inner ysep=3pt]{%
    \popxboldls\fontsize{7}{7}\selectfont\color{#3}\MakeUppercase{#1}};%
}
\newcommand{\popsquiggle}[1]{%
  \tikz[baseline=-0.4ex]{
    \draw[#1,line width=2.6pt,line cap=round]
      plot [smooth] coordinates {(0,0.09) (0.34,0.01) (0.68,0.21) (1.02,0.01) (1.36,0.10)};
  }%
}
% Mot-clé surligné (marqueur orange sous le texte)
\newcommand{\cle}[1]{{\popxbold\color{popdark}#1}}

% ============================================================
% En-tête / pied
% ============================================================
\pagestyle{fancy}
\fancyhf{}
\renewcommand{\headrulewidth}{0pt}
\renewcommand{\footrulewidth}{0pt}
\lhead{\raisebox{-0.15ex}{\poplogo\fontsize{13}{13}\selectfont\color{popink}OnBuch\color{poporange}+}}
\rhead{\poppill{\DOCMATIERE\ \textbullet\ \DOCNIVEAU\ \textbullet\ Leçon \DOCLECON}{white}{popink}}
\lfoot{\popxbold\fontsize{7.6}{7.6}\selectfont\color{popmuted}\MakeUppercase{Cours signé OnBuch+}\ \textbullet\ {\popsemi\DOCTITRE}}
\rfoot{\tikz[baseline=-0.6ex]\node[rounded corners=8.5pt,fill=popink,minimum height=17pt,minimum width=17pt,inner xsep=5pt,inner sep=0]{\popxbold\fontsize{8}{8}\selectfont\color{popcream}\thepage\,/\,\pageref*{LastPage}};}

% ============================================================
% Titres
% ============================================================
\newcommand{\chaptitre}[2]{%
  \begin{tcolorbox}[enhanced,colback=poporange,colframe=popink,boxrule=2.6pt,arc=11pt,
      drop shadow=popink,left=15pt,right=15pt,top=12pt,bottom=11pt,before skip=3pt,after skip=18pt]
    \poppill{#2}{popink}{popcream}\par\vspace{9pt}
    {\raggedright\hyphenpenalty=10000\exhyphenpenalty=10000\popdisplay\fontsize{22}{24}\selectfont\color{popcream}\MakeUppercase{#1}\par}
  \end{tcolorbox}
}
% Section numérotée : pastille orange avec le numéro + titre Archivo
\newcounter{popsec}
\newcommand{\coursec}[1]{%
  \stepcounter{popsec}\setcounter{popsub}{0}%
  \par\vspace{16pt}\noindent
  \begin{minipage}{\linewidth}
  \tikz[baseline=-0.7ex]\node[draw=popink,line width=1.6pt,fill=poporange,rounded corners=4pt,minimum size=22pt,inner sep=0]{\popdisplay\fontsize{12}{12}\selectfont\color{popcream}\thepopsec};%
  \hspace{8pt}{\popdisplay\fontsize{15}{17}\selectfont\color{popink}\MakeUppercase{#1}}\par
  \vspace{4pt}\noindent{\color{poporange}\rule{36pt}{3.6pt}}
  \end{minipage}\par\nopagebreak\vspace{8pt}
}
\newcounter{popsub}
\newcommand{\courssub}[1]{%
  \stepcounter{popsub}%
  \par\vspace{9pt}\noindent{\popxbold\fontsize{11.5}{13}\selectfont\color{popink}{\color{poporange}\thepopsec.\thepopsub}\hspace{6pt}#1}\par\nopagebreak\vspace{4pt}
}

% ============================================================
% Boîtes pédagogiques — même recette : fond blanc, bordure épaisse colorée,
% ombre dure encre, badge plein. Titre optionnel [#1].
% ============================================================
\tcbset{popbase/.style={breakable,enhanced,colback=white,boxrule=2.3pt,arc=9pt,
  shadow={3pt}{-3pt}{0pt}{popink},left=14pt,right=14pt,top=11pt,bottom=11pt,before skip=11pt,after skip=15pt,
  fontupper=\popsemi\fontsize{10.6}{14.5}\selectfont\color{popink},
  fontlower=\popsemi\fontsize{10.6}{14.5}\selectfont\color{popink}}}
% Coche dessinée (le glyphe ✓ n'existe pas dans les polices du document)
\newcommand{\popcheck}[1]{\tikz[baseline=-0.5ex]\draw[#1,line width=1.6pt,line cap=round,line join=round](-0.13,0.0)--(-0.04,-0.09)--(0.13,0.1);}
\providecommand{\coche}{\popcheck{popgreen}}
\providecommand{\resultat}[1]{\par\smallskip\noindent\fcolorbox{popgreen}{popgreenL}{\parbox{\dimexpr\linewidth-2\fboxsep-2\fboxrule\relax}{\textbf{Résultat.} #1}}\par\smallskip}
\newcommand{\popboxhead}[3]{\poppill{#1}{#2}{white}\ifx\relax#3\relax\else\hspace{6pt}{\popxbold\fontsize{10.6}{12}\selectfont\color{popink}#3}\fi\par\vspace{7pt}}

\newtcolorbox{aretenir}[1][]{popbase,colframe=popgreen,before upper={\popboxhead{À retenir}{popgreen}{#1}}}
% Texte en allemand (document de lecture, dialogue, extrait) : cadre
% d'étude. Un titre optionnel (source, auteur) s'affiche dans l'en-tête.
\newtcolorbox{textebox}[1][]{popbase,colframe=popblue,before upper={\popboxhead{Text}{popblue}{#1}\begin{otherlanguage}{ngerman}},after upper={\end{otherlanguage}}}
% Marques ✓ / ✗ (forme correcte / fautive) souvent utilisées par les modèles
\providecommand{\cross}{\textcolor{poppink}{\ensuremath{\times}}}
\providecommand{\xmark}{\cross}\providecommand{\crossmark}{\cross}\providecommand{\cmark}{\ensuremath{\checkmark}}
% Ligne de réponse / justification à compléter (inventée par les modèles)
\providecommand{\justification}[1]{\par\noindent\textit{Justification~:} \dotfill\par}
% \textipa (package tipa non chargé) : la police ne couvre pas l'API -> simple passage
\providecommand{\textipa}[1]{#1}
% Soulignement (ulem non chargé) : \uline, \ul
\providecommand{\uline}[1]{\underline{#1}}\providecommand{\ul}[1]{\underline{#1}}
% \glossaire{\item ...} inventée par les modèles : liste de glossaire
\providecommand{\glossaire}[1]{\begin{itemize}#1\end{itemize}}
% \alert (beamer) inventée par les modèles : texte mis en évidence
\providecommand{\alert}[1]{\textbf{\textcolor{poppink}{#1}}}
% variantes ASCII d'ordinaux écrites en macro
\providecommand{\iere}{\textsuperscript{re}}\providecommand{\ieme}{\textsuperscript{e}}\providecommand{\ere}{\textsuperscript{re}}\providecommand{\eme}{\textsuperscript{e}}\providecommand{\er}{\textsuperscript{er}}\providecommand{\ie}{\textsuperscript{e}}
% Ordinaux français écrits en macro par les modèles : 1\ière, 3\ième
\newcommand{\ière}{\textsuperscript{re}}\newcommand{\ième}{\textsuperscript{e}}
% \exemple (inventée par les modèles) : étiquette « Exemple : »
\providecommand{\exemple}{\textbf{Exemple~:}\ }
% \square utilisable aussi hors mode math (cases à cocher)
\AtBeginDocument{\let\squareMath\square\renewcommand{\square}{\ensuremath{\squareMath}}}
% Mot ou phrase en allemand à l'intérieur d'un paragraphe français
\newcommand{\all}[1]{\ifmmode\text{\textit{#1}}\else\begingroup\selectlanguage{ngerman}\itshape #1\endgroup\fi}
\let\esp\all % alias toléré
\newtcolorbox{exemplebox}[1][]{popbase,colframe=poporange,before upper={\popboxhead{Exemple}{poporange}{#1}}}
\newtcolorbox{definition}[1][]{popbase,colframe=popblue,before upper={\popboxhead{Définition}{popblue}{#1}}}
\newtcolorbox{propriete}[1][]{popbase,colframe=poppurple,before upper={\popboxhead{Propriété}{poppurple}{#1}}}
\newtcolorbox{methode}[1][]{popbase,colframe=popgold,before upper={\popboxhead{Méthode}{popgold}{#1}}}
\newtcolorbox{attention}[1][]{popbase,colframe=poppink,before upper={\popboxhead{Attention}{poppink}{#1}}}
\newtcolorbox{experience}[1][]{popbase,colframe=popblue,colback=popblueL,before upper={\popboxhead{Expérience}{popblue}{#1}}}
% « Le savais-tu ? » : carte violette pleine (contexte, culture, Cameroun)
\newtcolorbox{savaistu}[1][]{popbase,colback=poppurple,colframe=popink,
  fontupper=\popsemi\fontsize{10.4}{14.2}\selectfont\color{white},
  before upper={\poppill{Le savais-tu\,?}{white}{poppurple}\ifx\relax#1\relax\else\hspace{6pt}{\popxbold\color{white}#1}\fi\par\vspace{7pt}}}
% Exercice résolu : énoncé en haut, \tcblower puis la solution sur fond crème
\newcounter{popexo}\newcounter{popexr}
\newtcolorbox{exoresolu}[1][]{popbase,colframe=poporange,colbacklower=popcream,
  segmentation style={popink,line width=1.2pt,dashed},
  before upper={\stepcounter{popexr}\popboxhead{Exercice résolu \thepopexr}{poporange}{#1}},
  before lower={\poppill{Solution}{popink}{popcream}\par\vspace{6pt}}}
% Exercice (non résolu en place) — la correction est dans la partie Corrigés
\newtcolorbox{exercice}[1][]{popbase,colframe=popink,
  before upper={\stepcounter{popexo}\popboxhead{Exercice \thepopexo}{popink}{#1}}}
% Corrigé
\newtcolorbox{corrige}[1][]{popbase,colframe=popgreen,colback=popgreenL,
  before upper={\popboxhead{Corrigé}{popgreen}{#1}}}
% Objectifs / prérequis / bilan
\newtcolorbox{objectifs}{popbase,colframe=popink,colback=poporangeL,
  before upper={\popboxhead{Objectifs de la leçon}{popink}{}}}
\newtcolorbox{prerequis}{popbase,colframe=popink,colback=popblueL,
  before upper={\popboxhead{Prérequis}{popblue}{}}}
\newtcolorbox{bilan}{popbase,colframe=popink,colback=white,boxrule=2.8pt,
  before upper={\popboxhead{Fiche bilan}{poporange}{L'essentiel à connaître pour l'examen}}}
% Figure encadrée avec légende
\newtcolorbox{popfigure}[1][]{enhanced,breakable=false,colback=white,colframe=popline,boxrule=1.4pt,arc=8pt,
  left=10pt,right=10pt,top=10pt,bottom=8pt,before skip=10pt,after skip=12pt,halign=center,
  lower separated=false,halign lower=center,
  fontlower=\popsemi\fontsize{9}{11.5}\selectfont\color{popmuted},
  #1}
\newcommand{\legende}[1]{\par\vspace{6pt}{\popsemi\fontsize{9}{11.5}\selectfont\color{popmuted}#1}}

% ============================================================
% Signature OnBuch+
% ============================================================
\newcommand{\poplogoinline}[1]{{\poplogo\fontsize{#1}{#1}\selectfont\color{popink}OnBuch\color{poporange}+}}
\newcommand{\signatureonbuch}{%
  \begin{tcolorbox}[enhanced,colback=popink,colframe=popink,boxrule=2.2pt,arc=10pt,
      drop shadow=poporange,left=14pt,right=14pt,top=10pt,bottom=10pt,before skip=0pt,after skip=0pt]
    \tikz[baseline=-0.7ex]\node[circle,fill=popgreen,draw=popcream,line width=1.3pt,minimum size=22pt,inner sep=0]{\popcheck{white}};%
    \hspace{9pt}\begin{minipage}[c]{0.62\linewidth}
      {\popxbold\fontsize{12}{13}\selectfont\color{popcream}Signé\ \textbullet\ {\poplogo OnBuch\color{poporange}+}}\par\vspace{2pt}
      {\raggedright\popsemi\fontsize{8}{10.5}\selectfont\color{popline}\MakeUppercase{Équipe pédagogique OnBuch+}\\\MakeUppercase{Conforme au programme officiel MINESEC}\par}
    \end{minipage}\hfill
    {\poplogo\fontsize{22}{22}\selectfont\color{popcream}OnBuch\color{poporange}+}
  \end{tcolorbox}%
}

% ============================================================
% Page de garde
% #1 titre · #2 matière · #3 niveau · #4 module · #5 n° leçon · #6 durée ·
% #7 description / objectifs (texte ou itemize)
% ============================================================
\newcommand{\pagegarde}[7]{%
  \thispagestyle{empty}
  \newgeometry{a4paper,margin=1.7cm,top=1.6cm,bottom=1.4cm}%
  \begin{tikzpicture}[remember picture,overlay]
    \fill[poporange!18] ([xshift=-3.2cm,yshift=-3.6cm]current page.north east) circle (4.6cm);
    \fill[poppurple!14] ([xshift=2.4cm,yshift=7.6cm]current page.south west) circle (3.8cm);
  \end{tikzpicture}%
  {\poplogo\fontsize{30}{30}\selectfont\color{popink}OnBuch\color{poporange}+}\hfill
  \raisebox{0.6ex}{\poppill{Cours signé \textbullet\ Premium}{popink}{popcream}}\par
  \vspace{1pt}\popsquiggle{poppurple}\par
  \vspace{8pt}
  \begin{tcolorbox}[enhanced,colback=poporange,colframe=popink,boxrule=2.8pt,arc=13pt,
      drop shadow=popink,left=18pt,right=18pt,top=12pt,bottom=13pt,after skip=10pt]
    \poppill{#2 \textbullet\ #3}{popink}{popcream}\hspace{5pt}\poppill{#4}{white}{popink}\par\vspace{8pt}
    {\popdisplay\fontsize{14}{14}\selectfont\color{popink}LEÇON #5}\par\vspace{4pt}
    {\raggedright\hyphenpenalty=10000\exhyphenpenalty=10000\popdisplay\fontsize{25}{27}\selectfont\color{popcream}\MakeUppercase{#1}\par}
    \vspace{8pt}
    {\popxbold\fontsize{9.5}{9.5}\selectfont\color{popink}\MakeUppercase{Durée conseillée : #6}}
  \end{tcolorbox}
  \begin{tcolorbox}[enhanced,colback=white,colframe=popink,boxrule=2.2pt,arc=10pt,
      drop shadow=popink,left=16pt,right=16pt,top=10pt,bottom=10pt,after skip=8pt]
    \poppill{Dans cette leçon}{poppurple}{white}\par\vspace{5pt}
    \popsemi\fontsize{9.3}{12.6}\selectfont\color{popink}#7
  \end{tcolorbox}
  \noindent\begin{minipage}[t]{0.487\linewidth}
    \begin{tcolorbox}[enhanced,colback=popgreen,colframe=popink,boxrule=2.2pt,arc=10pt,
        drop shadow=popink,left=11pt,right=11pt,top=10pt,bottom=10pt,height=3.1cm,valign=center]
      \tcbox[colback=white,colframe=popink,boxrule=1.3pt,arc=5pt,boxsep=0pt,nobeforeafter,
        left=3pt,right=3pt,top=3pt,bottom=3pt,valign=center]{\qrcode[height=1.9cm]{https://play.google.com/store/apps/details?id=com.luvvix.onbuch&hl=fr}}%
      \hspace{8pt}\begin{minipage}[c]{0.5\linewidth}
        {\popdisplay\fontsize{11}{12.5}\selectfont\color{white}\MakeUppercase{Télécharge OnBuch}}\par\vspace{4pt}
        {\raggedright\popsemi\fontsize{8.4}{11}\selectfont\color{white}Gratuit \textbullet\ Google Play\\Tuteur IA, annales, concours, résultats\par}
      \end{minipage}
    \end{tcolorbox}
  \end{minipage}\hfill
  \begin{minipage}[t]{0.487\linewidth}
    \begin{tcolorbox}[enhanced,colback=poppurple,colframe=popink,boxrule=2.2pt,arc=10pt,
        drop shadow=popink,left=11pt,right=11pt,top=10pt,bottom=10pt,height=3.1cm,valign=center]
      \tikz[baseline=-0.7ex]\node[circle,fill=white,draw=popink,line width=1.3pt,minimum size=1.6cm,inner sep=0]{\popdisplay\fontsize{24}{24}\selectfont\color{poppurple}+};%
      \hspace{8pt}\begin{minipage}[c]{0.6\linewidth}
        {\popdisplay\fontsize{11}{12.5}\selectfont\color{white}\MakeUppercase{Passe à OnBuch+}}\par\vspace{4pt}
        {\raggedright\popsemi\fontsize{8.4}{11}\selectfont\color{white}Tous les cours signés, Tuteur IA illimité, fiches PDF — onglet OnBuch+ de l'app\par}
      \end{minipage}
    \end{tcolorbox}
  \end{minipage}
  \vspace{8pt}
  \popcontenu
  \vfill
  \signatureonbuch
  \restoregeometry
  \newpage
}

% Rangée « Ce cours contient » (page de garde)
\newcommand{\poptile}[3]{% #1 couleur #2 gros texte #3 légende
  \begin{tcolorbox}[enhanced,colback=#1,colframe=popink,boxrule=2pt,arc=9pt,shadow={2.5pt}{-2.5pt}{0pt}{popink},
      left=6pt,right=6pt,top=9pt,bottom=9pt,halign=center,valign=center,height=1.75cm,width=\linewidth,nobeforeafter]
    {\popdisplay\fontsize{12.5}{13}\selectfont\color{white}\MakeUppercase{#2}}\par\vspace{3pt}
    {\centering\popsemi\fontsize{7.8}{9.5}\selectfont\color{white}#3\par}
  \end{tcolorbox}}
\newcommand{\popcontenu}{%
  \par\noindent{\popxboldls\fontsize{7.5}{7.5}\selectfont\color{popmuted}CE COURS SIGNÉ CONTIENT}\par\vspace{7pt}
  \noindent\begin{minipage}[t]{0.235\linewidth}\poptile{popblue}{Cours}{complet, illustré, pas à pas}\end{minipage}\hfill
  \begin{minipage}[t]{0.235\linewidth}\poptile{poporange}{Méthodes}{et exercices résolus}\end{minipage}\hfill
  \begin{minipage}[t]{0.235\linewidth}\poptile{poppink}{Exercices}{type examen + corrigés}\end{minipage}\hfill
  \begin{minipage}[t]{0.235\linewidth}\poptile{popgreen}{Fiche bilan}{l'essentiel à retenir}\end{minipage}\par}

% Sommaire
\newtcolorbox{sommaire}{enhanced,colback=white,colframe=popink,boxrule=2.2pt,arc=10pt,
  drop shadow=popink,left=18pt,right=18pt,top=13pt,bottom=13pt,before skip=6pt,after skip=20pt,
  before upper={\poppill{Sommaire}{poppurple}{white}\par\vspace{9pt}\popsemi\fontsize{10.6}{14.5}\selectfont\color{popink}}}

% Signature de fin de cours — poussée tout en bas de la dernière page
\newcommand{\findecours}{%
  \par\vfill\nopagebreak
  \begin{center}\popsquiggle{poporange}\end{center}\vspace{4pt}
  \signatureonbuch
}
\AtBeginDocument{\def\llbracket{\mathopen{[\mkern-3.5mu[}}\def\rrbracket{\mathclose{]\mkern-3.5mu]}}} % intervalles d'entiers (glyphe ⟦ absent de la police)
% Glyphes absents de Fira Math : redéfinis à partir de glyphes disponibles
\AtBeginDocument{%
  \def\setminus{\mathbin{\backslash}}\let\smallsetminus\setminus
  \def\vdots{\mathord{\vbox{\baselineskip3.2pt\lineskiplimit0pt\kern4pt\hbox{.}\hbox{.}\hbox{.}}}}%
  \def\ddots{\mathinner{\mkern1mu\raise6.5pt\hbox{.}\mkern2mu\raise3.6pt\hbox{.}\mkern2mu\raise0.7pt\hbox{.}\mkern1mu}}%
  \def\triangle{\mathord{\scalebox{0.9}{$\symup{\Delta}$}}}%
  \def\nabla{\mathord{\raisebox{\depth}{\scalebox{1}[-1]{$\symup{\Delta}$}}}}%
  \def\checkmark{\popcheck{popdark}}%
  \def\star{\mathbin{\ast}}\def\bigstar{\mathord{\ast}}%
  \def\diamond{\mathbin{\scalebox{0.6}{$\Diamond$}}}%
  \def\bigcup{\mathop{\mathchoice{\raisebox{-0.3em}{\scalebox{1.7}{$\cup$}}}{\scalebox{1.25}{$\cup$}}{\cup}{\cup}}\displaylimits}%
  \def\bigcap{\mathop{\mathchoice{\raisebox{-0.3em}{\scalebox{1.7}{$\cap$}}}{\scalebox{1.25}{$\cap$}}{\cap}{\cap}}\displaylimits}%
  \def\lesseqqgtr{\mathrel{\lessgtr}}\def\bigoplus{\mathop{\oplus}\displaylimits}%
}
\providecommand{\ssi}{si et seulement si} % abréviation fréquente
\AtBeginDocument{\providecommand{\wideparen}{\overparen}} % arc de cercle

\begin{document}

\pagegarde{Traduction et production d'articles de journaux scolaires et de textes de presse}{Allemand}{1ère A}{Chapitre 4 : Média et communication}{ALA20}{2 h}{Cette leçon mixte (texte + traduction) guide l'élève de la lecture commentée d'un article de journal scolaire allemand vers la rédaction autonome d'articles, d'interviews et d'annonces, tout en consolidant la grammaire (Perfekt, Präteritum, ordre des mots) et en pratiquant le thème et la version de courts textes journalistiques.
\vspace{4pt}
\begin{multicols}{2}\small
\begin{itemize}
\item À la fin de cette leçon, je sais identifier la structure et le vocabulaire spécifiques d'un article, d'une interview et d'une annonce de presse en allemand.
\item Je sais analyser un texte de presse allemand (compréhension globale et détaillée) et en extraire les idées principales.
\item Je sais réemployer le Perfekt, le Präteritum et l'ordre des mots (V2, position du verbe en proposition subordonnée) dans mes propres productions.
\item Je sais traduire du français vers l'allemand (thème) et de l'allemand vers le français (version) des phrases et un court texte journalistiques en respectant le registre.
\item Je sais rédiger un article de journal scolaire, une interview et une annonce en allemand en respectant les conventions de mise en page.
\item Je sais relire et corriger ma production (orthographe, majuscules, cas, ordre des mots) à l'aide d'une grille d'auto‑évaluation.
\end{itemize}\end{multicols}}

\begin{sommaire}
\begin{multicols}{2}
\begin{enumerate}
\item 1. Texte support \& compréhension
\item 2. Lexique thématique – presse \& journal scolaire
\item 3. Grammaire – Révision Perfekt, Präteritum, ordre des mots
\item 4. Types de textes journalistiques – Article, Interview, Annonce
\item 5. Traduction – Thème (FR→DE) \& Version (DE→FR)
\item 6. Production guidée \& relecture
\item Activité d'intégration
\item Méthodes et exercices résolus
\item Exercices
\item Corrigés des exercices
\item Fiche bilan
\end{enumerate}
\end{multicols}
\end{sommaire}

\begin{objectifs}
\begin{itemize}
\item À la fin de cette leçon, je sais identifier la structure et le vocabulaire spécifiques d'un article, d'une interview et d'une annonce de presse en allemand.
\item Je sais analyser un texte de presse allemand (compréhension globale et détaillée) et en extraire les idées principales.
\item Je sais réemployer le Perfekt, le Präteritum et l'ordre des mots (V2, position du verbe en proposition subordonnée) dans mes propres productions.
\item Je sais traduire du français vers l'allemand (thème) et de l'allemand vers le français (version) des phrases et un court texte journalistiques en respectant le registre.
\item Je sais rédiger un article de journal scolaire, une interview et une annonce en allemand en respectant les conventions de mise en page.
\item Je sais relire et corriger ma production (orthographe, majuscules, cas, ordre des mots) à l'aide d'une grille d'auto‑évaluation.
\item Je sais mobiliser le lexique thématique (médias, rédaction, publication) dans des tâches de communication authentiques.
\end{itemize}
\end{objectifs}

\begin{prerequis}
\begin{itemize}
\item Connaissance de base du Perfekt et du Präteritum (formation auxiliaire + participe passé / terminaisons faibles et fortes).
\item Maîtrise de l'ordre des mots en phrase principale (V2) et en subordonnée (verbe en fin).
\item Vocabulaire de la vie scolaire (Klasse, Lehrer, Projekt, Veranstaltung) et notions de base sur les médias (Zeitung, Artikel, Redaktion).
\item Capacité à rédiger un paragraphe cohérent en allemand (niveau A2/B1).
\item Habitude de l'auto‑correction à l'aide de listes de contrôle.
\end{itemize}
\end{prerequis}

\newpage

\coursec{Texte support \& compréhension}

\courssub{Situation d'accroche et problématique}

Imaginez la scène : au lycée de Yaoundé, les élèves de Première préparent le prochain numéro du journal scolaire « La Voix du Lycée ». Leur professeur d'allemand leur propose un défi : publier une page entièrement en allemand. En feuilletant des journaux germanophones comme la \emph{Berliner Zeitung} ou la \emph{Süddeutsche Zeitung}, ils découvrent que la presse allemande suit des codes très précis : un titre accrocheur, un \all{Lead} (chapeau) qui résume l'essentiel, un corps d'article structuré en paragraphes, des interviews et des annonces. En comparant une brève dépêche de la \emph{Deutsche Welle} avec un article camerounais, ils perçoivent immédiatement deux difficultés majeures : le choix du temps (Perfekt ou Präteritum ?) et la place du verbe (la fameuse règle V2 : le verbe conjugué toujours en deuxième position).

\textbf{Problématique :} Comment lire, comprendre puis produire un article de presse scolaire en allemand, en respectant sa structure et ses marques linguistiques ?

Pour répondre à cette question, nous allons d'abord lire un article authentique paru dans une \all{Schulzeitung} allemande, l'analyser de près, puis en extraire les outils de langue qui nous serviront à rédiger et à traduire nos propres textes.

\begin{definition}[La Schulzeitung]
La \cle{Schulzeitung} (die, \emph{-en}) est le journal scolaire : un périodique rédigé et publié par les élèves, sous la supervision d'un professeur (der \all{Betreuer}). On y trouve des articles (der \all{Artikel}, \emph{-}), des reportages (der \all{Bericht}, \emph{-e}), des interviews (das \all{Interview}, \emph{-s}) et des annonces (die \all{Ankündigung}, \emph{-en}). Au Cameroun, l'équivalent serait « La Voix du Lycée ».
\end{definition}

\courssub{Lecture du texte support}

Lisez d'abord le texte en silence, sans dictionnaire, pour saisir le sens général. Ne vous arrêtez pas sur chaque mot inconnu : le contexte aide souvent à deviner.

\begin{textebox}[Unser Schulfest – Ein Rückblick (Schulzeitung „Die Schülermaus“, München)]
\textbf{Unser Schulfest – Ein Rückblick}\par
\emph{von Lena Hoffmann, Klasse 10b}\par
\medskip
\textbf{[1]} Am Samstag, den 15. Juni, hat unsere Schule ihr großes Sommerfest gefeiert. Über 500 Gäste – Eltern, Lehrer, Schüler und Freunde – sind gekommen.\par
\textbf{[2]} Zunächst eröffnete der Schulleiter, Herr Dr. Weber, das Fest um 10 Uhr mit einer kurzen Rede. Er bedankte sich bei allen Helfern und wünschte den Gästen einen schönen Tag.\par
\textbf{[3]} Danach gab es ein reichhaltiges Programm. Die Theatergruppe hat ein Stück von Brecht gespielt, und die Schulband sang deutsche und internationale Lieder. Außerdem konnten die Besucher an vielen Ständen Kuchen, Würstchen und Getränke kaufen. Die Klasse 10b verkaufte selbst gebackenen Kuchen – ein voller Erfolg!\par
\textbf{[4]} Am Nachmittag fand das traditionelle Fußballspiel Lehrer gegen Schüler statt. Die Schüler haben mit 3:1 gewonnen. Alle Zuschauer lachten und feierten laut.\par
\textbf{[5]} Schließlich endete das Fest am Abend mit einem großen Feuerwerk. Der Gewinn von über 800 Euro geht an ein Projekt für Kinder in Kamerun.\par
\textbf{[6]} Das Schulfest war ein unvergesslicher Tag. Wir danken allen Helfern und freuen uns schon auf das nächste Jahr!
\end{textebox}

\textbf{Glossaire personnel (mots difficiles) :}

\begin{center}
\begin{tabularx}{\linewidth}{|X|X|}\hline
\rowcolor{popblueL} \textbf{Deutsch} & \textbf{Français} \\ \hline
der Rückblick, -e & le retour en arrière, le bilan \\ \hline
der Gast, "e & l'invité \\ \hline
eröffnen (hat eröffnet) & ouvrir, inaugurer \\ \hline
die Rede, -n & le discours \\ \hline
sich bedanken bei + Dat. & remercier (quelqu'un) \\ \hline
der Helfer, - & l'aide, le bénévole \\ \hline
reichhaltig & riche, copieux \\ \hline
der Stand, "e & le stand \\ \hline
selbst gebacken & fait maison (cuit soi-même) \\ \hline
der Erfolg, -e & le succès \\ \hline
stattfinden (fand statt, hat stattgefunden) & avoir lieu \\ \hline
der Zuschauer, - & le spectateur \\ \hline
das Feuerwerk, -e & le feu d'artifice \\ \hline
der Gewinn, -e & le gain, le bénéfice \\ \hline
unvergesslich & inoubliable \\ \hline
\end{tabularx}
\end{center}

\begin{attention}[Bericht ou Artikel ?]
Ne confondez pas \all{der Bericht} (le reportage, le compte rendu) et \all{der Artikel} (l'article, terme générique). Le \all{Bericht} est plus narratif : il raconte des faits dans l'ordre chronologique, souvent au Präteritum. L'\all{Artikel} peut aussi donner une opinion ou analyser. Notre texte est un \all{Rückblick} (bilan), une forme de \all{Bericht}. Autre piège : \all{die Rede} signifie « le discours », pas « la règle » !
\end{attention}

\courssub{La structure de l'article : Titel, Lead, Hauptteil, Schluss}

Tout article de presse allemand suit une architecture fixe. Observons notre texte :

\begin{itemize}
\item \textbf{der Titel} (le titre) : « Unser Schulfest – Ein Rückblick ». Court, accrocheur, il annonce le sujet.
\item \textbf{der Lead / die Einleitung} (le chapeau, ligne 1) : il répond aux questions essentielles — \all{Wer? Was? Wann? Wo?} (qui ? quoi ? quand ? où ?). Ici : l'école (wer), la fête d'été (was), samedi 15 juin (wann), à l'école (wo).
\item \textbf{der Hauptteil} (le corps, lignes 2 à 5) : trois à quatre paragraphes qui détaillent les événements dans l'ordre chronologique.
\item \textbf{der Schluss} (la conclusion, ligne 6) : un bilan, un remerciement, une ouverture vers l'avenir.
\end{itemize}

\begin{popfigure}
\begin{center}
\begin{tikzpicture}[node distance=1.2cm,
box/.style={draw=popblue, fill=popblueL, rounded corners, thick, align=center, minimum width=6cm, minimum height=0.9cm}]
\node[box] (titel) {\textbf{Titel} (accrocheur)};
\node[box, below of=titel] (lead) {\textbf{Lead (Einleitung)} : Wer? Was? Wann? Wo?};
\node[box, below of=lead] (body) {\textbf{Hauptteil} (3--4 Absätze, chronologique)};
\node[box, below of=body] (schluss) {\textbf{Schluss} (bilan + ouverture)};
\draw[->, thick, poporange] (titel) -- (lead);
\draw[->, thick, poporange] (lead) -- (body);
\draw[->, thick, poporange] (body) -- (schluss);
\end{tikzpicture}
\end{center}
\legende{La structure en quatre parties d'un article de presse allemand.}
\end{popfigure}

\begin{aretenir}[Les marqueurs de structure]
Les connecteurs logiques guident le lecteur dans le corps de l'article. Repérez-les dans le texte :
\begin{itemize}
\item \all{Zunächst} (d'abord) — ligne 2 : début du récit ;
\item \all{Danach} (ensuite) — ligne 3 : enchaînement ;
\item \all{Außerdem} (de plus) — ligne 3 : ajout d'information ;
\item \all{Schließlich} (finalement) — ligne 5 : dernier événement.
\end{itemize}
Attention à la syntaxe : après ces connecteurs placés en tête de phrase, le verbe conjugué vient \emph{immédiatement} (règle V2) : \all{Zunächst \textbf{eröffnete} der Schulleiter das Fest.} et non *\all{Zunächst der Schulleiter eröffnete...}.
\end{aretenir}

\begin{methode}[La lecture active en quatre étapes]
\begin{enumerate}
\item \textbf{Première lecture globale} : lire sans dictionnaire, identifier le type de texte et le sujet.
\item \textbf{Surligner les connecteurs} (\all{zunächst, danach, außerdem, schließlich}) : ils révèlent le plan du texte.
\item \textbf{Noter les temps verbaux} : entourer les verbes au Perfekt et au Präteritum, se demander pourquoi l'auteur a choisi l'un ou l'autre.
\item \textbf{Poser des questions marginales} : dans la marge, écrire une question par paragraphe (Qui fait quoi ? Pourquoi ?). Cela prépare la compréhension détaillée et la traduction.
\end{enumerate}
\end{methode}

\courssub{Questions de compréhension}

Répondez en allemand, par des phrases complètes.

\begin{exercice}[Leseverstehen — Compréhension de l'écrit]
\begin{enumerate}
\item (Globale) \all{Worum geht es in diesem Text?} — De quoi parle ce texte ?
\item (Détaillée) \all{Wann hat das Schulfest stattgefunden und wie viele Gäste sind gekommen?}
\item (Détaillée) \all{Was hat die Klasse 10b verkauft?}
\item (Détaillée) \all{Wer hat das Fußballspiel gewonnen? Mit welchem Ergebnis?}
\item (Inférence) \all{Warum war das Fest auch ein soziales Ereignis?} — Pourquoi la fête a-t-elle aussi une dimension solidaire ?
\item (Inférence) \all{Wie findet die Autorin das Fest? Zitieren Sie zwei Wörter aus dem Text.}
\end{enumerate}
\end{exercice}

\begin{corrige}[Exercice « Leseverstehen »]
\begin{enumerate}
\item \all{Der Text berichtet über das Sommerfest einer Schule in München.} (Le texte rend compte de la fête d'été d'une école à Munich.)
\item \all{Das Fest hat am Samstag, den 15. Juni, stattgefunden. Über 500 Gäste sind gekommen.}
\item \all{Die Klasse 10b hat selbst gebackenen Kuchen verkauft.}
\item \all{Die Schüler haben mit 3:1 gegen die Lehrer gewonnen.}
\item \all{Der Gewinn von über 800 Euro geht an ein Projekt für Kinder in Kamerun. Das Fest hilft also auch anderen Menschen.}
\item \all{Sie findet das Fest sehr positiv: „ein voller Erfolg“ und „ein unvergesslicher Tag“.}
\end{enumerate}
\end{corrige}

\courssub{Repérage des temps verbaux : Perfekt et Präteritum}

Relisons le texte en entourant les verbes. On constate que l'auteure utilise \emph{deux temps du passé} :

\begin{exemplebox}[Deux temps pour raconter le passé]
\all{Die Schüler \textbf{organisierten} ein Fest.} (Präteritum) → Les élèves ont organisé une fête.\par
\all{Die Schüler \textbf{haben} ein Fest \textbf{organisiert}.} (Perfekt) → Les élèves ont organisé une fête.\par
\medskip
Les deux phrases se traduisent de la même façon en français, mais l'allemand distingue deux registres : le Präteritum est le temps du récit écrit (presse, roman, compte rendu), le Perfekt est le temps de l'oral et du style plus léger.
\end{exemplebox}

\begin{propriete}[Emplois du Perfekt et du Präteritum dans la presse]
\begin{itemize}
\item \textbf{Präteritum} : temps privilégié du récit écrit et journalistique. Obligatoire à l'écrit pour \all{sein} (\all{war}), \all{haben} (\all{hatte}) et les verbes modaux. Ex. ligne 2 : \all{eröffnete}, \all{bedankte}, \all{wünschte} ; ligne 4 : \all{fand ... statt}, \all{lachten}, \all{feierten}.
\item \textbf{Perfekt} : fréquent dans les articles de journal scolaire (style proche de l'oral) et pour les faits récents. Ex. ligne 1 : \all{hat ... gefeiert}, \all{sind gekommen} ; ligne 3 : \all{hat ... gespielt}.
\item \textbf{Règle pratique} : dans un même article, les deux temps peuvent coexister sans différence de sens. Les francophones doivent surtout éviter de traduire le Präteritum par un imparfait : \all{eröffnete} = « a ouvert » (passé composé), pas « ouvrait ».
\end{itemize}
\end{propriete}

\begin{center}
\begin{tabularx}{\linewidth}{|X|X|X|}\hline
\rowcolor{popblueL} \textbf{Infinitif} & \textbf{Präteritum (texte)} & \textbf{Perfekt (texte)} \\ \hline
feiern (fêter) & feierte & hat gefeiert \\ \hline
kommen (venir) & kam & ist gekommen \\ \hline
eröffnen (inaugurer) & eröffnete & hat eröffnet \\ \hline
stattfinden (avoir lieu) & fand statt & hat stattgefunden \\ \hline
gewinnen (gagner) & gewann & hat gewonnen \\ \hline
sein (être) & war & ist gewesen \\ \hline
geben (il y a) & gab & hat gegeben \\ \hline
\end{tabularx}
\end{center}

\begin{attention}[Pièges fréquents des francophones]
\begin{itemize}
\item \textbf{Auxiliaire} : les verbes de mouvement ou de changement prennent \all{sein} au Perfekt : \all{Über 500 Gäste \textbf{sind} gekommen} (et non *\all{haben gekommen}).
\item \textbf{Place du participe} : au Perfekt, le participe passé va \emph{à la fin} de la phrase : \all{Die Schule \textbf{hat} ihr Fest \textbf{gefeiert}.} Ne dites jamais *\all{Die Schule hat gefeiert ihr Fest}.
\item \textbf{Verbe à particule} : \all{stattfinden} se sépare au Präteritum : \all{Das Spiel \textbf{fand} am Nachmittag \textbf{statt}.}
\item \textbf{Majuscules} : tous les noms prennent la majuscule : \all{das Fest, die Gäste, der Gewinn}.
\end{itemize}
\end{attention}

\begin{exoresolu}[Du Präteritum au Perfekt]
\textbf{Énoncé.} Transformez au Perfekt la phrase extraite du texte : \all{Zunächst eröffnete der Schulleiter das Fest um 10 Uhr.}
\tcblower
\textbf{Solution détaillée.}
\begin{enumerate}
\item Verbe : \all{eröffnen} → participe passé : \all{eröffnet} (verbe faible, préfixe \all{er-} inséparable, donc \emph{pas} de \all{ge-}).
\item Auxiliaire : \all{eröffnen} n'est pas un verbe de mouvement → \all{haben}, conjugué : \all{hat}.
\item Ordre des mots : le connecteur \all{Zunächst} occupe la position 1, l'auxiliaire \all{hat} la position 2 (règle V2), le participe va en fin de phrase.
\end{enumerate}
\textbf{Réponse :} \all{Zunächst \textbf{hat} der Schulleiter das Fest um 10 Uhr \textbf{eröffnet}.} → D'abord, le directeur a inauguré la fête à 10 heures.
\end{exoresolu}

\begin{savaistu}[Journaux scolaires et Länder]
En Allemagne, l'école relève des Länder (régions) : chaque Land possède sa propre loi scolaire, le \all{Schulgesetz}, qui encadre notamment la publication des journaux scolaires — liberté de ton des élèves, mais responsabilité du professeur-conseil. Certaines \all{Schulzeitungen} allemandes existent depuis plus de cent ans ! Et au Cameroun ? Les clubs de journalisme des lycées, comme « La Voix du Lycée », jouent le même rôle de formation à l'expression et à la citoyenneté.
\end{savaistu}

\begin{experience}[À vous de jouer : lecture à voix haute]
Par groupes de quatre : chaque élève lit un paragraphe du texte à voix haute (attention à la prononciation : \all{Schulfest} se prononce « choul-fest », \all{Zuschauer} « tsou-chaou-eur »). Les trois autres notent les connecteurs et les temps verbaux entendus. Puis le groupe rédige ensemble un glossaire personnel de cinq mots nouveaux, avec article et pluriel pour chaque nom.
\end{experience}

\begin{aretenir}[Bilan de la section]
Un article de presse allemand se construit en quatre parties : \all{Titel → Lead → Hauptteil → Schluss}. Le récit au passé alterne Präteritum (style écrit) et Perfekt (style oral, presse scolaire). Les connecteurs (\all{zunächst, danach, außerdem, schließlich}) structurent le texte et imposent le verbe en deuxième position. Ces outils seront la base de nos propres productions et traductions dans les sections suivantes.
\end{aretenir}

\coursec{Lexique thématique – presse \ \& journal scolaire}

Dans la section précédente, nous avons lu et commenté un article de journal scolaire. Vous avez remarqué que ce texte contenait beaucoup de mots nouveaux : \all{die Redaktion}, \all{der Artikel}, \all{die Schlagzeile}... Avant de pouvoir rédiger ou traduire vous-mêmes des textes de presse, il faut maîtriser ce vocabulaire. C'est l'objectif de cette section : construire, champ par champ, votre lexique de la presse et du journal scolaire, avec l'article, le pluriel et un exemple pour chaque mot.

\courssub{Observer d'abord : le lexique en contexte}

Lisons d'abord un court texte de présentation d'un journal scolaire camerounais. Soulignez mentalement tous les mots qui appartiennent au monde de la presse.

\begin{textebox}[Unser Schülerjournal „Der Pfeil“]
An unserem Gymnasium in Yaoundé gibt es seit drei Jahren ein Schülerjournal. Es heißt „Der Pfeil“. Die Redaktion trifft sich jeden Montag nach dem Unterricht. Unsere Chefredakteurin, Marie-Claire, verteilt die Aufgaben: Wer schreibt einen Artikel über das Fußballturnier? Wer macht ein Interview mit dem neuen Deutschlehrer? Wer schreibt die Ankündigung für die Theateraufführung?

Jeder Artikel hat eine Schlagzeile und mehrere Absätze. Unter jedem Foto steht eine Bildunterschrift. Bevor wir das Layout fertig machen, lesen wir alle Texte noch einmal – das Korrekturlesen ist sehr wichtig! Dann veröffentlichen wir die Druckfassung und die Online-Ausgabe.

Unsere Texte sollen sachlich und objektiv sein, aber auch lebendig. Ein Kommentar darf meinungsstark sein, ein Bericht nicht. Am Ende des Schuljahres fassen wir die besten Artikel in einer Sonderausgabe zusammen.
\end{textebox}

\textbf{Glossaire :} \all{die Aufgabe, -n} la tâche ; \all{das Fußballturnier, -e} le tournoi de football ; \all{die Theateraufführung, -en} la représentation théâtrale ; \all{fertig machen} terminer (particule séparable : \all{macht ... fertig}) ; \all{die Sonderausgabe, -n} l'édition spéciale ; \all{der Bericht, -e} le compte rendu.

\textbf{Questions de compréhension :}
\begin{enumerate}
\item Wann trifft sich die Redaktion? (Quand la rédaction se réunit-elle ?)
\item Wer ist die Chefredakteurin? (Qui est la rédactrice en chef ?)
\item Was steht unter jedem Foto? (Qu'y a-t-il sous chaque photo ?)
\item Welche zwei Versionen veröffentlicht die Redaktion? (Quelles deux versions la rédaction publie-t-elle ?)
\end{enumerate}

\textit{Réponses attendues : 1. jeden Montag nach dem Unterricht ; 2. Marie-Claire ; 3. eine Bildunterschrift ; 4. die Druckfassung und die Online-Ausgabe.}

\courssub{Champ lexical 1 : la rédaction}

Le premier champ lexical concerne les personnes et les activités de la rédaction (\all{die Redaktion}).

\begin{definition}[Le vocabulaire de la rédaction]
\begin{itemize}
\item \all{die Redaktion, -en} : la rédaction (l'équipe et le lieu).
\item \all{der Chefredakteur, -e / die Chefredakteurin, -nen} : le/la rédacteur(-trice) en chef.
\item \all{das Ressort, -s} : la rubrique, le service (sport, culture...).
\item \all{das Layout, -s} : la mise en page.
\item \all{das Korrekturlesen} : la relecture (correction des fautes).
\end{itemize}
\end{definition}

\begin{exemplebox}[La rédaction en action]
\all{Die Redaktion trifft sich montags.} « La rédaction se réunit le lundi. »

\all{Der Chefredakteur verteilt die Ressorts.} « Le rédacteur en chef répartit les rubriques. »

\all{Vor der Druckfassung machen wir das Korrekturlesen.} « Avant la version imprimée, nous faisons la relecture. »
\end{exemplebox}

Remarquez la formation des mots composés, très fréquente en allemand : \all{Chef} (chef) + \all{Redakteur} (rédacteur) donne \all{der Chefredakteur} ; \all{Korrektur} (correction) + \all{lesen} (lire) donne \all{das Korrekturlesen}. Le genre du mot composé est toujours celui du \textbf{dernier} élément.

\courssub{Champ lexical 2 : les genres journalistiques}

\begin{definition}[Les genres journalistiques]
\begin{itemize}
\item \all{der Artikel, -} : l'article (texte informatif).
\item \all{das Interview, -s} : l'interview, l'entretien (questions/réponses).
\item \all{die Ankündigung, -en} : l'annonce (événement à venir).
\item \all{der Kommentar, -e} : le commentaire (opinion personnelle signée).
\item \all{die Reportage, -n} : le reportage (enquête sur le terrain).
\item \all{die Kolumne, -n} : la chronique (rubrique régulière et personnelle).
\end{itemize}
\end{definition}

\begin{definition}[L'annonce : die Ankündigung]
Une \all{Ankündigung} est une annonce courte : elle donne l'événement, la date, le lieu, parfois le prix. Le style est télégraphique : phrases brèves, informations essentielles d'abord.

\all{Theateraufführung: „Wilhelm Tell“, Freitag, 18 Uhr, Aula. Eintritt frei!} « Représentation théâtrale : “Guillaume Tell”, vendredi, 18 h, grande salle. Entrée libre ! »
\end{definition}

\begin{attention}[Orthographe : das Interview]
Le mot \all{Interview} s'écrit en allemand avec \textbf{un seul « v »} : \all{das Interview}. N'ajoutez jamais de « w » comme en français (« interview »). Autre piège : on dit \all{ein Interview führen} (mener un entretien) et non « machen ».

\all{Das Interview wurde geführt.} « L'interview a été menée. »
\end{attention}

\courssub{Champ lexical 3 : les verbes de la presse (avec Präteritum)}

Le programme de Première exige la maîtrise du Perfekt \textbf{et} du Prétérit. Voici les formes complètes pour les verbes forts et irréguliers.

\begin{definition}[Les verbes fréquents — formes conjuguées]
\begin{center}
\begin{tabular}{|l|l|l|l|l|}
\hline
\rowcolor{popblueL} Infinitif & Sens & Präsens (er/sie/es) & Präteritum (er/sie/es) & Perfekt (hat/ist + Partizip II) \\
\hline
\all{veröffentlichen} & publier & \all{veröffentlicht} & \all{veröffentlichte} & \all{hat veröffentlicht} \\
\hline
\all{interviewen} & interviewer & \all{interviewt} & \all{interviewte} & \all{hat interviewt} \\
\hline
\all{ankündigen} & annoncer & \all{kündigt an} & \all{kündigte an} & \all{hat angekündigt} \\
\hline
\all{berichten} & rendre compte & \all{berichtet} & \all{berichtete} & \all{hat berichtet} \\
\hline
\all{zusammenfassen} & résumer & \all{fasst zusammen} & \all{fasste zusammen} & \all{hat zusammengefasst} \\
\hline
\all{treffen} & toucher/atteindre & \all{trifft} & \all{traf} & \all{hat getroffen} \\
\hline
\end{tabular}
\end{center}
\end{definition}

\begin{exemplebox}[Les verbes en phrases]
\all{Wir veröffentlichen das Journal zweimal im Jahr.} « Nous publions le journal deux fois par an. »

\all{Die Reporterin interviewt den Bürgermeister von Douala.} « La journaliste interviewe le maire de Douala. »

\all{Der Schulleiter kündigt eine Deutschwoche an.} « Le proviseur annonce une semaine de l'allemand. » (particule \all{an-} rejetée en fin de phrase)

\all{Die Zeitung berichtete gestern über das Spiel der Löwen Indomptables.} « Le journal a rendu compte hier du match des Lions Indomptables. » (Prétérit)

\all{Fasse den Artikel in drei Sätzen zusammen!} « Résume l'article en trois phrases ! » (Impératif)
\end{exemplebox}

\begin{attention}[Verbes à particule séparable vs inséparable]
\all{ankündigen} et \all{zusammenfassen} sont \textbf{séparables} : au Präsens, la particule va à la fin (\all{Er kündigt das Fest \textbf{an}.}) ; au Prétérit, la particule reste à la fin (\all{Er kündigte das Fest \textbf{an}.}) ; au Perfekt, le \all{ge-} s'insère \textbf{entre} la particule et le radical (\all{hat an\textbf{ge}kündigt}, \all{hat zusammen\textbf{ge}fasst}).

\all{veröffentlichen} et \all{berichten} ont un préfixe \textbf{inséparable} (\all{ver-}, \all{be-}) : pas de séparation, \textbf{pas de \all{ge-}} au Perfekt (\all{hat veröffentlicht}, \all{hat berichtet}). Le Prétérit se forme régulièrement : \all{veröffentlichte}, \all{berichtete}.
\end{attention}

\courssub{Champ lexical 4 : les objets du journal et les adjectifs de style}

\begin{definition}[Les éléments d'une page de journal]
\begin{itemize}
\item \all{die Schlagzeile, -n} : le titre, le gros titre.
\item \all{der Lead, -s} : le chapeau (première phrase qui résume l'article).
\item \all{der Absatz, -¨e} : le paragraphe (pluriel : \all{die Absätze}).
\item \all{die Bildunterschrift, -en} : la légende (sous une photo).
\item \all{die Druckfassung, -en} : la version imprimée.
\item \all{die Online-Ausgabe, -n} : l'édition en ligne.
\end{itemize}
\end{definition}

\begin{definition}[Les adjectifs de style]
\begin{itemize}
\item \all{sachlich} : factuel, neutre.
\item \all{lebendig} : vivant, animé.
\item \all{prägnant} : concis, percutant.
\item \all{objektiv} : objectif.
\item \all{meinungsstark} : qui a des opinions tranchées.
\end{itemize}

\all{Ein Bericht muss sachlich und objektiv sein.} « Un compte rendu doit être factuel et objectif. »

\all{Ihre Kolumne ist lebendig und meinungsstark.} « Sa chronique est vivante et pleine d'opinions. »
\end{definition}

\courssub{Tableau récapitulatif et carte mentale}

\begin{center}
\begin{tabularx}{\linewidth}{|X|X|X|}
\hline
\rowcolor{popblueL} Deutsch & Français & Champ \\
\hline
die Redaktion, -en & la rédaction & Rédaction \\
\hline
der Chefredakteur, -e & le rédacteur en chef & Rédaction \\
\hline
das Ressort, -s & la rubrique & Rédaction \\
\hline
das Layout, -s & la mise en page & Rédaction \\
\hline
das Korrekturlesen & la relecture & Rédaction \\
\hline
der Artikel, - & l'article & Genres \\
\hline
das Interview, -s & l'interview & Genres \\
\hline
die Ankündigung, -en & l'annonce & Genres \\
\hline
der Kommentar, -e & le commentaire & Genres \\
\hline
die Reportage, -n & le reportage & Genres \\
\hline
die Kolumne, -n & la chronique & Genres \\
\hline
die Schlagzeile, -n & le gros titre & Objets \\
\hline
der Lead, -s & le chapeau & Objets \\
\hline
der Absatz, -¨e & le paragraphe & Objets \\
\hline
die Bildunterschrift, -en & la légende & Objets \\
\hline
die Druckfassung, -en & la version imprimée & Objets \\
\hline
die Online-Ausgabe, -n & l'édition en ligne & Objets \\
\hline
\end{tabularx}
\end{center}

\begin{popfigure}
\begin{tikzpicture}
\matrix (m) [matrix of nodes, nodes={anchor=west, minimum width=3.5cm, font=\small}, column sep=0.5cm, row sep=0.2cm] {
Deutsch & Français \\
\hline
Redaktion & Rédaction \\
Chefredakteur & Rédacteur en chef \\
Ressort & Rubrique \\
Layout & Mise en page \\
Korrekturlesen & Relecture \\
};
\draw[popblue, line width=1pt] (m-1-1.north west) rectangle (m-6-2.south east);
\foreach \i in {2,...,6} \draw[popline] (m-\i-1.north west) -- (m-\i-2.north east);
\draw[popline] (m-1-1.south west) -- (m-1-2.south east);
\end{tikzpicture}
\legende{Le champ lexical de la rédaction : cinq mots-clés à connaître par cœur.}
\end{popfigure}

\begin{popfigure}
\begin{tikzpicture}[mindmap, grow cyclic, every node/.style={concept, align=center, font=\small, text=white}, root concept/.append style={concept color=popblue, font=\normalsize\bfseries}, level 1/.append style={sibling angle=90, level distance=3.5cm}]
\node {Die Presse} [counterclockwise from=45]
  child[concept color=poporange] { node {Redaktion} }
  child[concept color=popgreen] { node {Genres} }
  child[concept color=poppurple] { node {Verben} }
  child[concept color=poppink] { node {Stil} };
\end{tikzpicture}
\legende{Carte mentale : les quatre champs lexicaux de la presse à compléter avec vos mots.}
\end{popfigure}

\courssub{Expressions idiomatiques}

\begin{definition}[Deux expressions à connaître]
\begin{itemize}
\item \all{ins Schwarze treffen} : littéralement « toucher dans le noir », c'est-à-dire \textbf{viser juste}, réussir parfaitement. \all{Die Schlagzeile trifft ins Schwarze.} « Le titre est parfaitement trouvé. »
\item \all{den Nagel auf den Kopf treffen} : littéralement « frapper le clou sur la tête », c'est-à-dire \textbf{mettre le doigt dessus}, dire exactement la vérité. \all{Mit ihrem Kommentar trifft sie den Nagel auf den Kopf.} « Avec son commentaire, elle met le doigt dessus. »
\end{itemize}
Attention : \all{treffen} est un verbe fort : \all{trifft, traf, hat getroffen}.
\end{definition}

\begin{savaistu}[D'où vient le mot „Zeitung“ ?]
Le mot \all{die Zeitung} (le journal) vient de \all{die Zeit} (le temps) + le suffixe \all{-ung} (qui exprime une action/résultat) : littéralement « ce qui donne le temps », c'est-à-dire ce qui informe sur l'actualité du moment. Le même suffixe se retrouve dans \all{die Ankündigung} (de \all{ankündigen}), \all{die Bildunterschrift} et \all{die Druckfassung}. Observer les familles de mots aide à deviner le sens des mots inconnus !
\end{savaistu}

\begin{methode}[Créer une fiche de révision par champ]
Pour mémoriser durablement ce lexique, fabriquez vos propres fiches :
\begin{enumerate}
\item Choisissez un champ, par exemple « Genres journalistiques ».
\item Notez 5 mots avec l'\textbf{article} et le \textbf{pluriel} : \all{der Artikel, -} ; \all{das Interview, -s}...
\item Écrivez pour chaque mot une \textbf{phrase exemple} complète : \all{Der Artikel hat drei Absätze.}
\item Ajoutez la traduction française au verso de la fiche.
\item Testez-vous en cachant une colonne, puis réutilisez les mots dans une production écrite.
\end{enumerate}
\end{methode}

\begin{exoresolu}[Complétez avec le mot qui convient]
Complétez : \all{Redaktion, Schlagzeile, Interview, Ankündigung, Korrekturlesen}.

a) Die \_\_\_ schreibt einen Artikel über das Schulfest.

b) Das \_\_\_ mit dem Fußballtrainer erscheint morgen.

c) Jeder Artikel braucht eine gute \_\_\_.

d) Die \_\_\_ informiert über das Konzert am Samstag.

e) Vor der Druckfassung ist das \_\_\_ sehr wichtig.
\tcblower
\textbf{Solution détaillée :}

a) \all{Redaktion} — c'est l'équipe qui écrit les articles. « La rédaction écrit un article sur la fête de l'école. »

b) \all{Interview} — un entretien avec un entraîneur (un seul « v » !). « L'interview avec l'entraîneur de football paraît demain. »

c) \all{Schlagzeile} — tout article a besoin d'un titre. « Chaque article a besoin d'un bon titre. »

d) \all{Ankündigung} — on informe sur un événement à venir. « L'annonce informe sur le concert de samedi. »

e) \all{Korrekturlesen} — la relecture précède l'impression. « Avant la version imprimée, la relecture est très importante. »
\end{exoresolu}

\begin{exercice}[À vous de jouer]
\begin{enumerate}
\item Traduisez en allemand :
      a) La rédaction publie un journal chaque trimestre.
      b) Le rédacteur en chef annonce une nouvelle rubrique.
      c) Résume l'article en deux phrases !
\item Classez ces mots dans les quatre champs (Rédaction, Genres, Objets, Style) :
      \all{die Kolumne, das Layout, prägnant, der Lead, die Reportage, sachlich, der Chefredakteur, die Bildunterschrift}.
\item Rédigez une \all{Ankündigung} (3 lignes, style télégraphique) pour un match de football de votre lycée : événement, date, lieu.
\end{enumerate}
\end{exercice}

\begin{corrige}[Exercice « À vous de jouer »]
1. a) \all{Die Redaktion veröffentlicht jedes Trimester ein Journal.}
   b) \all{Der Chefredakteur kündigt ein neues Ressort an.}
   c) \all{Fasse den Artikel in zwei Sätzen zusammen!}

2. Rédaction : \all{das Layout, der Chefredakteur} ; Genres : \all{die Kolumne, die Reportage} ; Objets : \all{der Lead, die Bildunterschrift} ; Style : \all{prägnant, sachlich}.

3. Modèle : \all{Fußballspiel: Klasse 1A gegen Tle C, Mittwoch, 16 Uhr, Sportplatz. Eintritt frei!}
\end{corrige}

\begin{aretenir}[L'essentiel du lexique]
\begin{itemize}
\item Toujours apprendre un nom avec son \textbf{article} et son \textbf{pluriel} : \all{die Schlagzeile, -n}.
\item \all{das Interview} s'écrit avec un seul « v » ; on \all{führt} un interview.
\item \all{ankündigen} et \all{zusammenfassen} sont séparables ; \all{veröffentlichen} et \all{berichten} sont inséparables (pas de \all{ge-} au Perfekt).
\item Une \all{Ankündigung} : courte, style télégraphique, événement + date + lieu.
\item Deux expressions idiomatiques avec \all{treffen} (fort : \all{traf, hat getroffen}) : \all{ins Schwarze treffen} et \all{den Nagel auf den Kopf treffen}.
\end{itemize}
\end{aretenir}

Ce lexique est maintenant votre boîte à outils. Dans la section suivante, nous réviserons la grammaire nécessaire — Perfekt, Präteritum et ordre des mots — pour employer ce vocabulaire dans des phrases correctes.

\coursec{Grammaire – Révision Perfekt, Präteritum, ordre des mots}

Dans la section précédente, nous avons enrichi notre vocabulaire de la presse et du journal scolaire : \all{die Zeitung}, \all{der Artikel}, \all{das Interview}, \all{die Anzeige}... Mais un bon article ne se contente pas de beaux mots : il exige une grammaire irréprochable. Or, pour raconter des événements — un match de football au lycée, une visite du proviseur, une fête scolaire à Yaoundé —, il faut maîtriser les temps du passé (\cle{Perfekt} et \cle{Präteritum}) et l'\cle{ordre des mots}, la règle d'or de la syntaxe allemande. C'est l'objet de cette révision guidée.

\courssub{Observation : deux façons de raconter le même événement}

Lisons d'abord ce court article paru dans le journal scolaire d'un lycée de Douala. Observez les verbes au passé : quelle forme revient le plus souvent ?

\begin{textebox}[Aus der Schülerzeitung „Die Brücke“, Douala]
\textbf{Großer Erfolg beim Sportfest}

Am Samstag fand das Sportfest unseres Gymnasiums statt. Über 300 Schüler nahmen teil. Der Direktor eröffnete die Veranstaltung um 9 Uhr mit einer kurzen Rede. Danach spielten die Fußballmannschaften der Klassen 1A und 1B gegeneinander. Das Spiel war sehr spannend: zuerst führte die Klasse 1B, aber in der zweiten Halbzeit drehte die Klasse 1A das Spiel und gewann 3:2. Die Zuschauer jubelten laut.

Eine Schülerin berichtete nach dem Spiel: „Wir haben sehr gut gespielt, und wir sind stolz auf unser Team!“ Auch der Sportlehrer war zufrieden: „Die Schüler haben den ganzen Tag fair gekämpft.“

Am Ende verteilte der Direktor die Medaillen, und alle gingen müde, aber glücklich nach Hause.

\emph{Glossaire :} das Sportfest, -e (la fête du sport) ; stattfinden (avoir lieu) ; teilnehmen (participer) ; die Veranstaltung, -en (l'événement) ; spannend (passionnant) ; führen (mener au score) ; die Halbzeit, -en (la mi-temps) ; gewinnen, gewann, hat gewonnen (gagner) ; der Zuschauer, - (le spectateur) ; jubeln (jubiler) ; stolz auf + Akkusativ (fier de) ; verteilen (distribuer) ; die Medaille, -n (la médaille).
\end{textebox}

\textbf{Questions d'observation :}
\begin{enumerate}
\item Soulignez tous les verbes conjugués au passé dans le texte. Combien de formes différentes voyez-vous ?
\item Où se trouvent les formes en \all{hat/haben + participe} ? Pourquoi, à votre avis ?
\item Dans la phrase \all{Am Samstag fand das Sportfest statt}, quelle est la place du verbe conjugué ?
\end{enumerate}

\textbf{Constat :} le récit journalistique écrit utilise massivement le \cle{Präteritum} (\all{fand}, \all{nahmen teil}, \all{eröffnete}, \all{spielten}, \all{war}, \all{gewann}, \all{verteilte}, \all{gingen}), tandis que les paroles rapportées au style direct (les citations orales) utilisent le \cle{Perfekt} (\all{haben gespielt}, \all{haben gekämpft}). C'est exactement la distinction que nous allons étudier.

\courssub{Le Perfekt : formation et emploi}

\begin{propriete}[Formation du Perfekt : auxiliaire + Partizip II]
Le Perfekt se forme avec l'auxiliaire \cle{haben} ou \cle{sein} conjugué au présent (en position 2) + le \cle{Partizip II} placé \textbf{à la fin} de la phrase.

\textbf{1. Le Partizip II régulier (verbes faibles) :} \all{ge- + radical + -t}
\begin{itemize}
\item \all{machen} $\rightarrow$ \all{gemacht} ; \all{spielen} $\rightarrow$ \all{gespielt} ; \all{loben} $\rightarrow$ \all{gelobt}
\item \all{organisieren} $\rightarrow$ \all{organisiert} (les verbes en \all{-ieren} ne prennent \textbf{pas} de \all{ge-} !)
\end{itemize}

\textbf{2. Le Partizip II irrégulier (verbes forts) :} \all{ge- + radical modifié + -en}
\begin{itemize}
\item \all{gehen} $\rightarrow$ \all{gegangen} ; \all{schreiben} $\rightarrow$ \all{geschrieben} ; \all{nehmen} $\rightarrow$ \all{genommen} ; \all{finden} $\rightarrow$ \all{gefunden}
\end{itemize}

\textbf{3. Verbes à particule :} séparables : \all{stattfinden} $\rightarrow$ \all{stattgefunden} (le \all{ge-} s'insère entre la particule et le radical) ; inséparables : \all{berichten} $\rightarrow$ \all{berichtet}, \all{verteilen} $\rightarrow$ \all{verteilt} (sans \all{ge-}).

\textbf{4. Choix de l'auxiliaire :}
\begin{itemize}
\item \cle{sein} : verbes de déplacement ou de changement d'état (\all{gehen, kommen, fahren, aufstehen, sterben}) + \all{sein} et \all{bleiben} $\rightarrow$ \all{Wir sind gegangen.}
\item \cle{haben} : tous les autres (la grande majorité) $\rightarrow$ \all{Wir haben organisiert.}
\end{itemize}
\end{propriete}

\begin{exemplebox}[Le Perfekt à l'oral]
\all{Der Lehrer hat die Klasse gelobt.} « Le professeur a félicité la classe. » (conversation, récit oral)

\all{Wir haben das Sportfest gut organisiert.} « Nous avons bien organisé la fête du sport. »

\all{Die Schüler sind müde nach Hause gegangen.} « Les élèves sont rentrés fatigués à la maison. »

\all{Hat der Direktor die Medaillen verteilt?} « Le directeur a-t-il distribué les médailles ? »
\end{exemplebox}

\courssub{Le Präteritum : la forme du récit écrit}

\begin{propriete}[Formation du Präteritum]
\textbf{1. Verbes faibles :} radical + \all{-te} (je/il), \all{-ten} (nous/ils), \all{-test} (tu), \all{-tet} (vous).
\begin{itemize}
\item \all{organisieren} $\rightarrow$ \all{ich organisierte, wir organisierten}
\item \all{spielen} $\rightarrow$ \all{er spielte} ; \all{berichten} $\rightarrow$ \all{sie berichtete}
\end{itemize}

\textbf{2. Verbes forts :} voyelle altérée, \textbf{sans terminaison} à la 1\iere{} et 3\ieme{} personne du singulier. Il faut apprendre les trois formes : \all{infinitif – Präteritum – Partizip II}.
\begin{itemize}
\item \all{gehen – ging – ist gegangen} ; \all{finden – fand – hat gefunden} ; \all{schreiben – schrieb – hat geschrieben} ; \all{gewinnen – gewann – hat gewonnen} ; \all{nehmen – nahm – hat genommen} ; \all{kommen – kam – ist gekommen}
\end{itemize}

\textbf{3. Verbes mixtes :} voyelle altérée (comme les forts) + terminaison \all{-te} (comme les faibles).
\begin{itemize}
\item \all{bringen – brachte – hat gebracht} ; \all{denken – dachte – hat gedacht} ; \all{wissen – wusste – hat gewusst}
\end{itemize}

\textbf{4. Les indispensables :} \all{sein – war – ist gewesen} ; \all{haben – hatte – hat gehabt} ; \all{werden – wurde – ist geworden} ; \all{es gab} (il y avait).
\end{propriete}

\begin{center}
\begin{tabularx}{\linewidth}{|X|X|X|X|}
\hline
\rowcolor{popblueL} Infinitif & Präteritum & Partizip II & Français \\
\hline
sein & war & ist gewesen & être \\
\hline
haben & hatte & hat gehabt & avoir \\
\hline
gehen & ging & ist gegangen & aller \\
\hline
kommen & kam & ist gekommen & venir \\
\hline
finden & fand & hat gefunden & trouver \\
\hline
schreiben & schrieb & hat geschrieben & écrire \\
\hline
gewinnen & gewann & hat gewonnen & gagner \\
\hline
bringen & brachte & hat gebracht & apporter \\
\hline
organisieren & organisierte & hat organisiert & organiser \\
\hline
berichten & berichtete & hat berichtet & rapporter \\
\hline
\end{tabularx}
\end{center}

\begin{methode}[Perfekt ou Präteritum ? La question à se poser]
\begin{enumerate}
\item Je me demande : « Est-ce un texte \textbf{écrit formel} (article de presse, récit, compte rendu) ? » $\rightarrow$ j'utilise le \cle{Präteritum} : \all{Der Lehrer lobte die Klasse.}
\item Est-ce une \textbf{conversation}, un dialogue, une citation orale, une lettre personnelle ? $\rightarrow$ j'utilise le \cle{Perfekt} : \all{Der Lehrer hat die Klasse gelobt.}
\item Cas particulier : \all{sein} et \all{haben} s'emploient souvent au Präteritum même à l'oral : \all{Ich war müde. Ich hatte keine Zeit.} (plus naturel que \all{ich bin gewesen}).
\item Dans un article, je cite les paroles des personnes au \textbf{style direct} avec le Perfekt, et je raconte les faits au Präteritum.
\end{enumerate}
\end{methode}

\begin{exemplebox}[Comparaison oral / écrit]
\all{Der Lehrer hat die Klasse gelobt.} « Le professeur a félicité la classe. » (Perfekt, oral)

\all{Der Lehrer lobte die Klasse.} « Le professeur félicita la classe. » (Präteritum, écrit journalistique)

\all{„Wir haben 3:2 gewonnen!“, sagte der Kapitän.} « "Nous avons gagné 3 à 2 !", dit le capitaine. » (Perfekt dans la citation, Präteritum dans la narration)
\end{exemplebox}

\begin{attention}[Erreur fréquente des francophones]
Ne traduis pas le passé composé français mot à mot ! Le français « a félicité » peut se traduire par \all{hat gelobt} (oral) \textbf{ou} \all{lobte} (écrit formel) selon le contexte. Dans un \textbf{article de presse formel}, n'utilise \textbf{pas} le Perfekt pour la narration : \all{Der Direktor hat die Veranstaltung eröffnet} est une faute de registre dans un article ; écris \all{Der Direktor eröffnete die Veranstaltung}. Réserve le Perfekt aux citations orales entre guillemets.
\end{attention}

\begin{savaistu}[Le Präteritum en Suisse alémanique]
En Suisse alémanique, le Präteritum est presque totalement absent à l'oral : les Suisses disent toujours \all{ich bin gegangen}, \all{ich habe gesehen}, même pour des événements lointains. Seuls \all{war} et \all{hatte} survivent parfois. À l'écrit, en revanche, la presse suisse utilise le Präteritum comme en Allemagne. Le sud de l'Allemagne et l'Autriche connaissent la même tendance : plus on descend vers le sud, plus le Perfekt domine !
\end{savaistu}

\courssub{L'ordre des mots : la règle V2 et la subordonnée}

\begin{propriete}[Phrase principale : le verbe conjugué en position 2 (V2)]
Dans une phrase principale allemande affirmative, le \textbf{verbe conjugué occupe toujours la deuxième place}. La première place peut être occupée par le sujet \textbf{ou} par un complément (temps, lieu...) ; dans ce cas, le sujet passe \textbf{après} le verbe (inversion).

\all{Der Direktor eröffnete um 9 Uhr die Veranstaltung.} (sujet – verbe – reste)

\all{Um 9 Uhr eröffnete der Direktor die Veranstaltung.} (complément – verbe – sujet – reste)

« À 9 heures, le directeur ouvrit la cérémonie. » — Attention : le verbe reste en position 2, jamais en position 3 !
\end{propriete}

\begin{propriete}[Phrase subordonnée : le verbe conjugué à la fin]
Après les conjonctions de subordination \all{dass} (que), \all{weil} (parce que), \all{wenn} (quand/si), \all{obwohl} (bien que), le verbe conjugué se place \textbf{en position finale}.

\all{Der Lehrer sagte, dass die Schüler gut gespielt haben.} « Le professeur a dit que les élèves avaient bien joué. »

\all{Wir freuen uns, weil unsere Mannschaft gewonnen hat.} « Nous nous réjouissons parce que notre équipe a gagné. »

\all{Wenn das Spiel beginnt, sind alle Zuschauer still.} « Quand le match commence, tous les spectateurs se taisent. »
\end{propriete}

\begin{propriete}[L'ordre des compléments : TEKAMOLO]
Quand plusieurs compléments se succèdent dans la phrase, ils respectent l'ordre \cle{TEKAMOLO} : \textbf{TE}mporel (quand ?) – \textbf{KA}usal (pourquoi ?) – \textbf{MO}dal (comment ?) – \textbf{LO}kal (où ?).

\all{Die Schüler spielten am Samstag wegen des Sportfestes sehr motiviert auf dem Schulhof.}

« Les élèves jouèrent samedi (TE), à cause de la fête du sport (KA), très motivés (MO), dans la cour de l'école (LO). »
\end{propriete}

\begin{popfigure}
\begin{tikzpicture}[node distance=1.2cm]
\node[draw=popblue, fill=popblueL, rounded corners, align=center, font=\small] (v2) {V2 – Position 2 : verbe conjugué};
\node[draw=poporange, fill=poporangeL, rounded corners, align=center, font=\small, below of=v2] (sub) {Subordonnée (dass, weil, wenn) – verbe à la fin};
\node[draw=popgreen, fill=popgreenL, rounded corners, align=center, font=\small, below of=sub] (tek) {TEKAMOLO – ordre des compléments};
\draw[->, thick, popdark] (v2) -- (sub);
\draw[->, thick, popdark] (sub) -- (tek);
\end{tikzpicture}
\legende{Les trois piliers de la syntaxe allemande : place du verbe et ordre des compléments.}
\end{popfigure}

\begin{center}
\begin{tabularx}{\linewidth}{|X|X|X|}
\hline
\rowcolor{popblueL} Français & Deutsch & Remarque \\
\hline
Le directeur a ouvert la fête à 9 h. & Der Direktor eröffnete das Fest um 9 Uhr. & Sujet – verbe (V2) \\
\hline
À 9 h, le directeur a ouvert la fête. & Um 9 Uhr eröffnete der Direktor das Fest. & Inversion : verbe toujours en 2\ieme{} position \\
\hline
Il a dit que les élèves avaient bien joué. & Er sagte, dass die Schüler gut gespielt haben. & Verbe final en subordonnée \\
\hline
Les élèves sont rentrés fatigués à la maison. & Die Schüler sind müde nach Hause gegangen. & Modal (müde) avant Lokal (nach Hause) \\
\hline
\end{tabularx}
\end{center}

\courssub{Cas accusatif et datif après les verbes de communication}

\begin{propriete}[Verbes de communication et leurs cas]
Certains verbes de communication régissent l'\textbf{accusatif} (la personne interrogée est un objet direct), d'autres le \textbf{datif} :
\begin{itemize}
\item \textbf{+ accusatif :} \all{interviewen} (interviewer), \all{befragen} (interroger), \all{fragen} (demander à) $\rightarrow$ \all{Der Journalist interviewte den Direktor.} « Le journaliste interviewa le directeur. »
\item \textbf{+ datif :} \all{antworten} (répondre à), \all{danken} (remercier) $\rightarrow$ \all{Der Direktor antwortete dem Journalisten.} « Le directeur répondit au journaliste. »
\end{itemize}
En français, « répondre à quelqu'un » a une préposition ; en allemand, \all{antworten} prend directement le datif, \textbf{sans préposition}.
\end{propriete}

\begin{attention}[Pièges de traduction]
\begin{itemize}
\item « Répondre au journaliste » = \all{dem Journalisten antworten} (datif), jamais \all{an den Journalisten antworten}.
\item « Interviewer quelqu'un » = \all{jemanden interviewen} (accusatif) : \all{Wir befragten den Sportlehrer.} « Nous avons interrogé le professeur de sport. »
\item Ne pas confondre \all{fragen} (demander) et \all{befragen} (interroger officiellement) : dans un interview de presse, on préfère \all{befragen} ou \all{interviewen}.
\item N'oublie pas la majuscule à tous les noms : \all{der Journalist, die Klasse, das Sportfest}.
\end{itemize}
\end{attention}

\begin{exoresolu}[Remettre la phrase dans le bon ordre et au bon temps]
\textbf{Énoncé :} Tu rédiges un article pour le journal scolaire. Mets les éléments suivants dans le bon ordre, au Präteritum : \all{um 10 Uhr / der Direktor / die Medaillen / verteilen / auf dem Schulhof}. Puis transforme la phrase en citation orale (Perfekt) introduite par : \all{Eine Schülerin sagte: „...“}.

\tcblower

\textbf{Solution détaillée :}

\textbf{Étape 1 – Article (Präteritum, registre écrit).} Le complément de temps \all{um 10 Uhr} peut occuper la première place ; le verbe conjugué \all{verteilte} reste en position 2, le sujet suit (inversion), le complément de lieu \all{auf dem Schulhof} vient après l'objet (TEKAMOLO : le lieu est en dernier) :

\all{Um 10 Uhr verteilte der Direktor die Medaillen auf dem Schulhof.}

« À 10 heures, le directeur distribua les médailles dans la cour de l'école. »

\textbf{Étape 2 – Citation orale (Perfekt).} Auxiliaire \all{haben} (verbe transitif, pas de déplacement) + Partizip II \all{verteilt} en fin de phrase :

\all{Eine Schülerin sagte: „Der Direktor hat die Medaillen auf dem Schulhof verteilt.“}

« Une élève dit : "Le directeur a distribué les médailles dans la cour de l'école." »

\textbf{Vérification :} verbe en position 2 dans les deux phrases principales ; participe à la fin au Perfekt ; majuscules aux noms (\all{Direktor}, \all{Medaillen}, \all{Schulhof}, \all{Schülerin}).
\end{exoresolu}

\begin{exercice}[À toi de jouer]
\begin{enumerate}
\item Conjugue au Präteritum : \all{sein, haben, kommen, schreiben, bringen, organisieren} (3\ieme{} personne du singulier).
\item Réécris ce début d'article au Präteritum : \all{Am Freitag hat ein Journalist unsere Schule besucht. Er hat den Direktor interviewt und ist dann in die Mensa gegangen.}
\item Mets dans le bon ordre : \all{gestern / die Schüler / wegen des Regens / im Klassenzimmer / gespielt / haben}.
\item Traduis : « Le journaliste a interrogé le capitaine, et le capitaine a répondu au journaliste. »
\end{enumerate}
\end{exercice}

\begin{corrige}[Exercice « À toi de jouer »]
\begin{enumerate}
\item \all{er war, er hatte, er kam, er schrieb, er brachte, er organisierte}.
\item \all{Am Freitag besuchte ein Journalist unsere Schule. Er interviewte den Direktor und ging dann in die Mensa.} (Article écrit $\rightarrow$ Präteritum partout ; \all{besuchen} est un verbe à particule inséparable : \all{besuchte}, \all{hat besucht}, sans \all{ge-}.)
\item \all{Gestern haben die Schüler wegen des Regens im Klassenzimmer gespielt.} (TE : gestern ; KA : wegen des Regens ; LO : im Klassenzimmer ; participe final.)
\item \all{Der Journalist befragte den Kapitän, und der Kapitän antwortete dem Journalisten.} (\all{befragen} + accusatif : \all{den Kapitän} ; \all{antworten} + datif : \all{dem Journalisten}.)
\end{enumerate}
\end{corrige}

\begin{aretenir}[L'essentiel de la section]
\begin{itemize}
\item \textbf{Perfekt} = \all{haben/sein} (position 2) + Partizip II (fin de phrase) : registre \textbf{oral}, conversations, citations.
\item \textbf{Präteritum} = forme simple (\all{organisierte, ging, brachte, war, hatte}) : registre \textbf{écrit}, articles de presse, récits.
\item Phrase principale : verbe conjugué en \textbf{position 2}, quoi qu'il y ait en première place.
\item Subordonnée (\all{dass, weil, wenn}) : verbe conjugué \textbf{à la fin}.
\item Compléments : ordre \textbf{TEKAMOLO}.
\item \all{interviewen/befragen} + accusatif ; \all{antworten} + datif.
\end{itemize}
\end{aretenir}

\coursec{Types de textes journalistiques – Article, Interview, Annonce}

Après avoir révisé les temps du récit (\emph{Perfekt} pour l'oral/le compte-rendu, \emph{Präteritum} pour l'écrit journalistique) et l'ordre des mots (place du verbe en position 2, structure du champ central), nous abordons maintenant la mise en pratique : comment structurer un texte journalistique en allemand. La grammaire n'est pas une fin en soi ; elle sert la \cle{fonction communicative} du texte. Un article informe, une interview révèle une parole, une annonce appelle à l'action. Chacun a ses codes rigoureux.

\courssub{L'article de presse : structure en entonnoir et registre objectif}

Commençons par l'observation d'un article type, tel qu'il pourrait paraître dans le journal du \emph{Lycée Bilingue de Yaoundé} ou dans une gazette scolaire allemande comme \emph{„Schulzeitung Aktuell“}.

\begin{textebox}[Exemple d'article scolaire : La fête de l'environnement]
\textbf{Umweltfest am Goethe-Gymnasium: Ein voller Erfolg}

\textbf{Yaoundé.} Am vergangenen Samstag feierte das Goethe-Gymnasium sein jährliches Umweltfest. Mehr als 500 Schüler, Eltern und Lehrer nahmen an der Veranstaltung teil, die vom Umweltclub der Schule organisiert worden war.

Schon am Morgen bauten die Mitglieder des Clubs Stände auf, an denen Recycling-Projekte vorgestellt wurden. „Wir wollen zeigen, dass Mülltrennung Spaß macht“, erklärte die Schülersprecherin Amina Njoya im Interview mit unserer Zeitung. Der Erlös des Festes, insgesamt 1.200.000 FCFA, soll für die Anschaffung neuer Solarpaneele für das Schulgebäude verwendet werden.

Das Highlight war ein Modenschau mit Kleidung aus recycelten Materialien. Die Jury, bestehend aus zwei Kunstlehrern und einem lokalen Designer, kürte drei Gewinner. Der Schulleiter, Herr Dr. Manga, lobte in seiner Abschlussrede das Engagement der Jugend: „Ihr seid die Zukunft dieses Landes.“

Das nächste Umweltfest ist für das kommende Schuljahr geplant. Interessierte können sich ab sofort beim Umweltclub melden.
\end{textebox}

\begin{definition}[Lead (der Einleitungsabsatz)]
Der \textbf{Lead} ist der erste Absatz eines Artikels. Er beantwortet die \cle{5 W-Fragen} (Wer? Was? Wann? Wo? Warum?) so vollständig und prägnant wie möglich. Er fasst die Kernnachricht zusammen, damit der Leser sofort das Wichtigste erfährt.
\end{definition}

\noindent
Analysons la structure de cet article, classique du \cle{journalisme informatif} (nachrichtenjournalistisch), souvent appelée \emph{Structure en pyramide inversée} (wichtigste Info zuerst).

\begin{popfigure}
\begin{tikzpicture}[
  level distance=1.6cm,
  sibling distance=3.2cm,
  every node/.style={draw=popblue, thick, rounded corners, fill=popblueL, font=\sffamily\small, align=center, inner sep=4pt},
  edge from parent/.style={draw=popblue, thick, ->, >={Stealth[length=3mm]}}
]
\node[align=center]{Artikel\\Struktur}
  child {node[align=center]{Titel\\\emph{(Headline)}}}
  child {node[align=center]{Lead\\\emph{(5 W-Fragen)}}}
  child {node[align=center]{Hauptteil\\Details, Fakten, Hintergründe}
    child {node[align=center]{Fakt 1\\+ Quelle/Zitat}}
    child {node[align=center]{Fakt 2\\+ Daten/Zahl}}
    child {node[align=center]{Zitat\\Persönliche Stimme}}
  }
  child {node[align=center]{Schluss\\Ausblick, Kontakt, Nächstes}};
\end{tikzpicture}
\legende{Structure type d'un article de presse : du général (Lead) au détail (Hauptteil) vers la conclusion (Schluss).}
\end{popfigure}

\begin{propriete}[Structure type : Titel $\rightarrow$ Lead $\rightarrow$ Hauptteil $\rightarrow$ Schluss]
\begin{itemize}
  \item \textbf{Titel (Überschrift) :} Nominal (souvent sans verbe), présent pour l'actualité immédiate, précis, accrocheur. Ex: \all{Umweltfest am Goethe-Gymnasium: Ein voller Erfolg}.
  \item \textbf{Lead :} 1 paragraphe. Répond aux 5 W. Verbe souvent au \emph{Präteritum} (écrit) ou \emph{Perfekt} (oral/reportage récent). Lieu (Dateline) en début de lead : \all{Yaoundé.}
  \item \textbf{Hauptteil (Corps) :} Développement par paragraphes distincts (un idée = un paragraphe). Ordre décroissant d'importance. Intégration de \textbf{Zitate} (citations) au style direct (\all{„...“}) ou indirect (\emph{Konjunktiv I} : \all{Sie erklärte, dass das Fest ein Erfolg sei.}). Utilisation de connecteurs logiques (\all{zunächst, außerdem, jedoch, schließlich}).
  \item \textbf{Schluss (Schlussabsatz) :} Souvent un \emph{Ausblick} (perspective future : \all{Das nächste Fest ist für... geplant}), une information de contact ou une citation finale forte.
\end{itemize}
\end{propriete}

\begin{center}
\begin{tabularx}{\linewidth}{|X|X|X|}
\hline
\rowcolor{popblueL} \textbf{Élément structurel} & \textbf{Deutsch} & \textbf{Français / Fonction} \\
\hline
Titre & \all{Titel / Schlagzeile} & Nominal, présent, informatif \\
Lieu/Date & \all{Ortszeile (z. B. Yaoundé.)} & Ancre spatio-temporelle \\
Lead & \all{Lead / Einleitung} & 5 W (Wer, Was, Wann, Wo, Warum) \\
Corps & \all{Hauptteil} & Détails, citations, données, causes/conséquences \\
Citation directe & \all{Wörtliche Rede / Zitat} & \all{„...“} + Verbum dicendi (\all{sagen, erklären, betonen}) \\
Citation indirecte & \all{Indirekte Rede (Konj. I)} & \all{Er sagte, er \textbf{komme} morgen.} (Subjonctif I) \\
Conclusion & \all{Schluss / Ausblick} & Futur, appel, info pratique \\
\hline
\end{tabularx}
\end{center}

\begin{attention}[Piège majeur : Le « ich » dans l'article objectif]
En français, on écrit parfois « J'ai assisté à la fête... » ou « Nous avons vu... ». En allemand journalistique \textbf{objectif} (\emph{sachlich}), le journaliste s'efface.
\begin{itemize}
  \item \textbf{Interdit :} \all{Ich sah viele Schüler. / Wir berichten heute...}
  \item \textbf{Correct (Passif) :} \all{Es wurden viele Schüler gesehen. / Es wird berichtet...}
  \item \textbf{Correct (Impersonnel « man ») :} \all{Man sah viele Schüler. / Man kann sagen...}
  \item \textbf{Correct (Sujet inanimé) :} \all{Das Fest bot viele Attraktionen. / Der Erlös beträgt...}
\end{itemize}
L'usage du \textbf{Passiv} (\emph{werden} + Partizip II) ou de \textbf{man} garantit la neutralité et le professionnalisme.
\end{attention}

\begin{exoresolu}[Analyse de structure : Identifier les 5 W dans le Lead]
\textbf{Énoncé :} Relevez les réponses aux 5 questions W dans le Lead de l'article ci-dessus (\emph{„Am vergangenen Samstag... teil“}).

\tcblower
\textbf{Solution détaillée :}
\begin{itemize}
  \item \textbf{Wer?} \all{das Goethe-Gymnasium} (Sujet / Acteur principal de la phrase).
  \item \textbf{Was?} \all{feierte ... sein jährliches Umweltfest} (Action / Événement principal).
  \item \textbf{Wann?} \all{Am vergangenen Samstag} (Temps).
  \item \textbf{Wo?} \all{am Goethe-Gymnasium} (Lieu — complété par la \emph{Ortszeile} \all{Yaoundé.}).
  \item \textbf{Warum? / Wozu?} \all{das vom Umweltclub der Schule organisiert worden war} (Contexte/Organisateur → but implicite : événement scolaire).
\end{itemize}
\emph{Remarque :} Le \emph{Warum} (la cause profonde) est souvent développé dans le \emph{Hauptteil} (ici : collecte de fonds pour panneaux solaires). Les participants (\all{Mehr als 500 Schüler...}) apparaissent dans la phrase suivante, toujours dans le Lead.
\end{exoresolu}

\courssub{L'interview : fidélité à la parole et mise en scène dialoguée}

L'interview (\emph{das Interview}, plur. \emph{die Interviews}) diffère radicalement de l'article : elle reproduit la \cle{parole vivante}. Ce n'est pas le journaliste qui résume, c'est l'interviewé qui s'exprime.

\begin{definition}[Interview (das Interview)]
Ein \textbf{Interview} ist ein dokumentiertes Gespräch zwischen einem Interviewer (Journalist) und einem Interviewpartner (Experte, Betroffener, Prominenter). Es wird im \textbf{Frage-Antwort-Format} wiedergegeben. Kennzeichen: \textbf{Einleitung} (Kontext, Vorstellung der Person), \textbf{wörtliche Rede} (exakte Zitate mit „...“), \textbf{Schlussdank}.
\end{definition}

\begin{textebox}[Exemple d'interview scolaire : La déléguée des élèves]
\textbf{„Wir wollen gehört werden“ – Interview mit der Schülersprecherin Amina Njoya}

\textbf{Frage:} Frau Njoya, Sie sind seit drei Monaten Schülersprecherin. Was hat sich bisher geändert?
\textbf{Antwort:} Wir haben einen festen Sprechstundentag eingeführt. Jeden Dienstag in der großen Pause kann jeder Schüler zu uns in den SV-Raum kommen. Das war vorher nicht so geregelt.

\textbf{Frage:} Welche Themen liegen Ihnen besonders am Herzen?
\textbf{Antwort:} Die Toiletten-Situation ist katastrophal. Und wir brauchen mehr Schattenplätze auf dem Schulhof. Ich habe dem Schulleiter schon einen Katalog mit 20 Vorschlägen übergeben.

\textbf{Frage:} Wie reagiert die Schulleitung?
\textbf{Antwort:} Herr Dr. Manga hört zu. Aber Entscheidungen dauern lange. Wir bleiben dran. Nächste Woche treffen wir den Elternbeirat.

\textbf{Frage:} Haben Sie eine Botschaft an die Mitschüler?
\textbf{Antwort:} Engagiert euch! Die SV (Schülervertretung) ist nur stark, wenn viele mitmachen. Danke für das Gespräch.
\end{textebox}

\begin{propriete}[Caractéristiques formelles de l'interview]
\begin{itemize}
  \item \textbf{Titre :} Souvent une citation choc ou une question : \all{„Wir wollen gehört werden“}.
  \item \textbf{Einleitung (optionnelle mais courante) :} Contexte, lieu, présentation de l'invité (nom, fonction). Ici intégrée dans la 1ère question.
  \item \textbf{Format Q/A :} \textbf{Frage:} (ou \all{F:}) / \textbf{Antwort:} (ou \all{A:}) en gras. Pas de guillemets pour les questions/réponses elles-mêmes (le deux-points suffit), mais guillemets \textbf{obligatoires} pour les mots cités \emph{dans} la réponse (\all{„Wir bleiben dran.“}).
  \item \textbf{Style direct (wörtliche Rede) :} Fidélité absolue. On ne corrige pas la syntaxe orale (anacoluthes, ellipses) sauf incompréhensible. Verbes au \emph{Präsens} ou \emph{Perfekt} (langage oral).
  \item \textbf{Schlussdank :} Politesse rituelle : \all{Danke für das Gespräch.} / \all{Danke für Ihre Zeit.}
\end{itemize}
\end{propriete}

\begin{center}
\begin{tabularx}{\linewidth}{|X|X|X|}
\hline
\rowcolor{poporangeL} \textbf{Français (Entretien)} & \textbf{Deutsch (Interview)} & \textbf{Remarque} \\
\hline
Question : / Réponse : & \textbf{Frage:} / \textbf{Antwort:} & Gras + deux-points, pas de guillemets \\
« Que prévoyez-vous ? » & \all{Frage: Was planen Sie?} & Verbe en position 2 (question W) \\
Réponse : Nous prévoyons une fête. & \all{Antwort: Wir planen ein Fest.} & Phrase complète, style direct \\
Il a dit : « Je viens. » & \all{Er sagte: „Ich komme.“} & Guillemets allemands „ “ pour citation \\
Merci pour l'entretien. & \all{Danke für das Interview.} / \all{Danke für das Gespräch.} & Formule de politesse finale \\
\hline
\end{tabularx}
\end{center}

\begin{exoresolu}[Traduction Thème : Mini-interview]
\textbf{Énoncé :} Traduisez en allemand ce court échange pour un journal scolaire.
\textit{Question : Quel est le thème de la fête cette année ?}
\textit{Réponse : Le thème est « Afrique et Environnement ». Nous voulons planter des arbres.}
\textit{Question : Quand la fête a-t-elle lieu ?}
\textit{Réponse : Elle a lieu samedi prochain à 10 heures. Merci.}

\tcblower
\textbf{Solution détaillée :}
\begin{itemize}
  \item \textbf{Frage:} Was ist das Thema der Feier in diesem Jahr? (\emph{Thème} = \all{das Thema}; \emph{cette année} = \all{in diesem Jahr} / \all{dieses Jahr}).
  \item \textbf{Antwort:} Das Thema ist „Afrika und Umwelt“. Wir wollen Bäume pflanzen. (\emph{Environnement} = \all{die Umwelt}; guillemets pour le titre du thème ; \emph{planter} = \all{pflanzen}, verbe fort \emph{pflanzt, pflanzte, hat gepflanzt}).
  \item \textbf{Frage:} Wann findet die Feier statt? (\emph{Avoir lieu} = \all{stattfinden}, verbe à particule séparable $\rightarrow$ \all{findet ... statt}).
  \item \textbf{Antwort:} Sie findet nächsten Samstag um 10 Uhr statt. Danke für das Gespräch. (\emph{Samedi prochain} = \all{nächsten Samstag} [accusatif de temps] ; \emph{à 10 heures} = \all{um 10 Uhr} ; formule de politesse standard).
\end{itemize}
\end{exoresolu}

\courssub{L'annonce (Ankündigung) : brièveté, futur et informations pratiques}

L'annonce (\emph{die Ankündigung}, \emph{die Bekanntmachung}, \emph{der Aushang}) est un texte \textbf{fonctionnel}, très court, à visée \textbf{appelative} (inciter à venir, à s'inscrire, à acheter). Elle répond aux 5 W de manière télégrammatique.

\begin{definition}[Ankündigung (Annonce / Avis)]
Eine \textbf{Ankündigung} ist ein kurzer, informationsdichter Text, der auf ein zukünftiges Ereignis hinweist. Sie beantwortet knapp die W-Fragen (Wer veranstaltet? Was? Wann? Wo? Warum/Wozu?). Häufig steht das Verb im \textbf{Futur I (werden + Infinitiv)} oder im \textbf{Präsens} (für feste Pläne). Wichtig: \textbf{Kontaktdaten} (E-Mail, Telefon, Raum, Ansprechpartner).
\end{definition}

\begin{textebox}[Exemple d'annonce : Tournoi de football]
\textbf{⚽ SCHULFUSSBALLTURNIER 2025 ⚽}

\textbf{Was:} Jährliches Fußballturnier der Klassen 10–13 (Mädchen \& Jungen)
\textbf{Wann:} Samstag, 15. März 2025, 09:00–16:00 Uhr
\textbf{Wo:} Sportplatz des Lycée Bilingue, Yaoundé
\textbf{Wer:} Organisiert von der SV (Schülervertretung) \& Sportlehrern
\textbf{Anmeldung:} Bis 08.03. per E-Mail an \texttt{sv@lycee-bilingue.cm} oder bei Herrn Kouam (Raum 102).
\textbf{Startgebühr:} 1.000 FCFA pro Team (inkl. Wasser \& Obst)
\textbf{Preise:} Pokale für 1.–3. Platz, Fair-Play-Pokal

\textbf{Kommt zahlreich und unterstützt eure Teams!}
\end{textebox}

\begin{propriete}[Caractéristiques linguistiques de l'Ankündigung]
\begin{itemize}
  \item \textbf{Structure visuelle :} Mots-clés en gras (\textbf{Was:}, \textbf{Wann:}, \textbf{Wo:}...) $\rightarrow$ lecture diagonale rapide.
  \item \textbf{Temps verbaux :}
        \begin{itemize}
          \item \textbf{Futur I (\all{werden} + Infinitiv) :} Pour l'annonce formelle, la promesse. \all{Das Turnier \textbf{wird stattfinden}.}
          \item \textbf{Präsens :} Pour les faits fixés, horaires. \all{Das Turnier \textbf{beginnt} um 09:00 Uhr.}
          \item \textbf{Infinitivgruppen / Nominalstil :} \all{Anmeldung bis zum 08.03.} (pas de verbe conjugué, gain de place).
        \end{itemize}
  \item \textbf{Registre :} \emph{Appellativ} (impératif ou infinitif à valeur impérative) : \all{Kommt zahlreich!} / \all{Bitte anmelden!}
  \item \textbf{Informations de contact :} Obligatoires (E-Mail, personne, lieu, date limite).
\end{itemize}
\end{propriete}

\begin{center}
\begin{tabularx}{\linewidth}{|X|X|X|}
\hline
\rowcolor{popgreenL} \textbf{W-Frage} & \textbf{Deutsch (Fragewort)} & \textbf{Exemple dans l'annonce} \\
\hline
Qui ? & \all{Wer?} & \all{Organisiert von der SV} \\
Quoi ? & \all{Was?} & \all{Fußballturnier der Klassen 10–13} \\
Quand ? & \all{Wann?} & \all{Samstag, 15. März 2025, 09:00–16:00 Uhr} \\
Où ? & \all{Wo?} & \all{Sportplatz des Lycée Bilingue} \\
Pourquoi / But ? & \all{Warum? / Wozu?} & \all{Preise: Pokale... Fair-Play} \\
Contact / Comment ? & \all{Wie? / Wo?} & \all{Anmeldung per E-Mail an... / bei Herrn Kouam} \\
\hline
\end{tabularx}
\end{center}

\courssub{Différence cruciale : Article (Bericht) vs. Commentaire (Kommentar)}

C'est une distinction fondamentale en examen (écrit comme oral). L'article \textbf{informe} (faits), le commentaire \textbf{opine} (opinion).

\begin{attention}[Article vs. Kommentar – Ne pas confondre !]
\begin{center}
\begin{tabularx}{\linewidth}{|X|X|X|}
\hline
\rowcolor{poppurpleL} \textbf{Kriterium} & \textbf{Artikel / Bericht (Reportage)} & \textbf{Kommentar / Glosse (Meinungstext)} \\
\hline
\textbf{Ziel} & Informieren (objektiv) & Meinung äußern, bewerten (subjektiv) \\
\textbf{Ich-Form} & \textbf{Nein} (Passiv, \all{man}, sujet inanimé) & \textbf{Ja} (\all{Ich finde..., Meiner Meinung nach...}) \\
\textbf{Verben} & \emph{Präteritum}, \emph{Perfekt}, \emph{Passiv} & \emph{Präsens}, \emph{Modalverben} (\all{sollen, müssen, können}) \\
\textbf{Adjektive} & Factuels (\all{groß, neu, teuer, 500}) & Évaluatifs (\all{wichtig, falsch, großartig, skandalös}) \\
\textbf{Structure} & Pyramide inversée (Lead $\rightarrow$ Details) & Thèse $\rightarrow$ Arguments $\rightarrow$ Conclusion \\
\textbf{Exemple} & \all{Das Fest fand statt. Es kamen 500 Leute.} & \all{Das Fest war ein großer Erfolg. Man hätte aber mehr Werbung machen sollen.} \\
\hline
\end{tabularx}
\end{center}
\textbf{Consigne d'examen :} Si le sujet demande „Schreiben Sie einen \textbf{Bericht}“, n'utilisez \textbf{jamais} « \all{Ich finde} » ou « \all{Meiner Meinung nach} ». Si le sujet demande „Schreiben Sie einen \textbf{Kommentar}“, l'opinion personnelle est \textbf{exigée}.
\end{attention}

\courssub{Mise en page, relecture et méthode de rédaction}

Un texte journalistique allemand soigne sa \cle{mise en page} (Layout) : colonnes (\emph{Spalten}), intertitres (\emph{Zwischenüberschriften}), images avec légende (\emph{Bild mit Bildunterschrift}), ligne d'auteur (\emph{Autorenzeile} : \all{Von Max Mustermann, 10a}).

\begin{methode}[Rédaction d'un article scolaire : 5 étapes]
\begin{enumerate}
  \item \textbf{Collecte \& Plan (Material \& Gliederung) :} Notez les faits (mots-clés), triez par importance. Répondez aux 5 W pour le Lead.
  \item \textbf{Rédigez le Lead en premier :} Une phrase (ou deux) max. Verbe au \emph{Präteritum}. Intégrez lieu/date.
  \item \textbf{Développez le Hauptteil :} Un paragraphe par idée nouvelle. Intégrez 1--2 citations (\emph{Zitate}) au style direct. Utilisez des connecteurs (\all{zunächst, darüber hinaus, jedoch, schließlich}).
  \item \textbf{Rédigez le Schluss (Ausblick) :} Perspective future ou info pratique (contact, prochaine date).
  \item \textbf{Relecture systématique (Korrekturlesen) :} Vérifiez : \textbf{Majuscules aux noms} (\all{der Schüler, die Umwelt}), \textbf{Ordre des mots} (Verbe position 2, verbe à la fin en subordonnée), \textbf{Cas} (Acc/Dat après prépositions/verbes), \textbf{Orthographe} (Umlaute \all{ä/ö/ü}, \all{ß} vs \all{ss}), \textbf{Guillemets} (\all{„ “}).
\end{enumerate}
\end{methode}

\begin{experience}[Atelier rédaction : « Une journée au lycée »]
\textbf{Consigne :} En binôme. Rédigez un \textbf{Lead} (3--4 lignes) pour un article sur « La journée portes ouvertes au lycée » (fictive ou réelle).
\begin{itemize}
  \item Utilisez le \emph{Präteritum}.
  \item Intégrez les 5 W.
  \item Pas de « ich ».
  \item Titre accrocheur.
\end{itemize}
\textbf{Exemple de production attendue :}
\all{Türen auf für die Zukunft: Tag der offenen Tür am Gymnasium Yaoundé.}
\all{Yaoundé. Am vergangenen Samstag öffnete das Gymnasium Yaoundé seine Türen für über 300 Besucher. Eltern und Grundschüler informierten sich über das bilinguale Angebot und die neuen naturwissenschaftlichen Labore.}
\end{experience}

\courssub{Exercice d'application guidé : De la prise de notes à l'article}

Cet exercice résolu synthétise la grammaire (temps, ordre des mots) et la structure journalistique.

\begin{exoresolu}[Production écrite : Écrire un Lead + 1 paragraphe Hauptteil]
\textbf{Énoncé :} À partir des notes suivantes, rédigez en allemand le \textbf{Titel}, le \textbf{Lead} et \textbf{un paragraphe du Hauptteil} (avec une citation directe) pour le journal scolaire.
\textit{Notes :}
\begin{itemize}
  \item Événement : Concert de la chorale du lycée (\emph{Schulchor}).
  \item Date/Lieu : Vendredi dernier, 20 octobre, Aula du lycée (\emph{Aula}).
  \item Participants : 80 chanteurs, chef de chœur M. Essomba (\emph{Chorleiter Herr Essomba}).
  \item Public : 400 personnes (élèves, parents, proviseur).
  \item Programme : Chants allemands, africains, gospel. Standing ovations.
  \item Citation chef de chœur : « Les élèves ont travaillé dur, le résultat est émouvant. »
  \item But : Collecte pour nouveaux pupitres (\emph{neue Notenpulte}).
\end{itemize}

\tcblower
\textbf{Solution détaillée (Modèle de correction) :}

\textbf{Chorkonzert in der Aula: Standing Ovations für 80 Stimmen}

\textbf{Yaoundé.} Am vergangenen Freitag, dem 20. Oktober, gab der Schulchor des Gymnasiums Yaoundé sein Jahreskonzert in der Aula. Unter der Leitung von Chorleiter Herrn Essomba sangen 80 Schülerinnen und Schüler vor rund 400 Zuhörern deutsche, afrikanische und gospellähnliche Lieder. Der Erlös des Abends soll der Anschaffung neuer Notenpulte dienen.

\emph{Remarques grammaticales :}
\begin{itemize}
  \item \textbf{Titel} : Nominal, deux points, accrocheur.
  \item \textbf{Lead} : \all{Yaoundé.} (Ortszeile). \all{gab ... sein Jahreskonzert} (Präteritum, verbe position 2). \all{Unter der Leitung ...} (Préposition \all{unter} + Datif). \all{sangen ... vor ... Zuhörern} (Präteritum verbe fort \emph{singen}, \emph{sang}, \emph{hat gesungen}). \all{sollen der Anschaffung ... dienen} (Passif substitutif \emph{sollen} + Infinitif + Datif \emph{der Anschaffung}).
  \item \textbf{Hauptteil (extrait) :}
\end{itemize}

\all{„Die Schüler haben hart gearbeitet, das Ergebnis ist bewegend“, sagte Herr Essomba im Gespräch mit unserer Zeitung. Er betonte, dass der Chor seit August zweimal wöchentlich geprobt habe. Der Schulleiter, Herr Dr. Manga, dankte den jungen Sängerinnen und Sängern für ihren Einsatz und kündigte an, dass die Schule die Finanzierung der Notenpulte unterstützen werde.}

\emph{Analyse du paragraphe :}
\begin{itemize}
  \item Citation directe : \all{„...“, sagte Herr Essomba.} (Verbum dicendi au Präteritum).
  \item Discours indirect : \all{Er betonte, dass der Chor ... \textbf{geprobt habe}.} (\textbf{Konjunktiv I} \emph{habe} pour la subordonnée \emph{dass} — niveau A2/B1, à reconnaître).
  \item Connecteurs : \all{und, dass, für}.
  \item Majuscules : \all{Schülerinnen, Schüler, Zuhörern, Erlös, Anschaffung, Notenpulte, Chorleiter, Schulleiter}.
\end{itemize}
\end{exoresolu}

\begin{savaistu}[Presse scolaire en Allemagne : « Schülerzeitung »]
En Allemagne, la \emph{Schülerzeitung} est une institution démocratique forte. Elle bénéficie d'une protection légale (\emph{Pressefreiheit}, Art. 5 GG). Des concours nationaux comme \emph{„Schülerzeitungswettbewerb der Länder«} récompensent les meilleures rédactions. Les thèmes vont de la vie scolaire (cantine, profs, numérique) aux grands débats sociétaux (climat, racisme, IA). Au Cameroun, les journaux scolaires (\emph{Le Lycéen}, \emph{La Voix du Lycée}) jouent un rôle similaire : former à la citoyenneté, à l'écriture et à l'esprit critique. Rédiger en allemand pour le journal du lycée, c'est s'inscrire dans cette tradition vivante !
\end{savaistu}

\coursec{Traduction – Thème (FR→DE) \& Version (DE→FR)}

Dans la section précédente, nous avons découvert les trois grands types de textes journalistiques — l'\all{Artikel}, l'\all{Interview} et l'\all{Anzeige} — et leurs structures propres. Mais au lycée, on ne se contente pas de rédiger : on traduit aussi. Traduire un article de journal scolaire, c'est mettre en pratique tout ce que nous avons révisé dans ce chapitre : le \cle{Perfekt}, le \cle{Präteritum}, l'\cle{ordre des mots} et le vocabulaire de la presse. Cette section vous donne une méthode rigoureuse pour le \cle{thème} (français $\rightarrow$ allemand) et la \cle{version} (allemand $\rightarrow$ français), deux exercices classiques à l'examen de Première et au Baccalauréat.

\courssub{Observation : quatre paires de phrases FR/DE}

Avant toute règle, observons. Lisez attentivement les quatre paires suivantes, extraites d'articles de journaux scolaires, et repérez ce qui change entre le français et l'allemand.

\begin{exemplebox}[Quatre paires à observer]
\begin{enumerate}
\item \textbf{FR :} Les élèves ont organisé une fête. \par \textbf{DE :} \all{Die Schüler organisierten ein Fest.} \par \emph{Observation : le passé composé français devient un Präteritum allemand (texte écrit).}
\item \textbf{FR :} Hier, la rédaction a publié un nouvel article. \par \textbf{DE :} \all{Gestern veröffentlichte die Redaktion einen neuen Artikel.} \par \emph{Observation : le verbe reste en deuxième position (V2), même après „Gestern“ ; le sujet passe après le verbe.}
\item \textbf{FR :} Le journal est lu par tous les élèves. \par \textbf{DE :} \all{Die Zeitung wird von allen Schülern gelesen.} \par \emph{Observation : le passif se construit avec „werden“ + participe II ; „par“ = „von“ + datif.}
\item \textbf{FR :} Il y a beaucoup de lecteurs au Cameroun. \par \textbf{DE :} \all{Es gibt viele Leser in Kamerun.} \par \emph{Observation : le „il“ expletif français devient „es“ — on ne l'oublie jamais.}
\end{enumerate}
\end{exemplebox}

\begin{attention}[Ce qu'il faut retenir de l'observation]
Dans chaque paire, trois choses changent presque toujours : le \textbf{temps du verbe} (passé composé $\rightarrow$ Präteritum dans un article écrit), la \textbf{place du verbe} (toujours V2 en allemand), et les \textbf{mots outils} (il y a = \all{es gibt} ; par = \all{von} + Dativ). La traduction n'est jamais mot à mot : c'est un changement de système.
\end{attention}

\courssub{La règle d'or du thème journalistique}

\begin{propriete}[Règle de traduction : le temps cible]
En \textbf{thème} (FR $\rightarrow$ DE), le choix du temps allemand dépend du \textbf{type de texte cible} :
\begin{itemize}
\item \textbf{Article, annonce, compte rendu écrit} $\rightarrow$ on privilégie le \textbf{Präteritum} : \all{Der Direktor gab den Termin bekannt.} « Le directeur a annoncé la date. »
\item \textbf{Dialogue, interview, lettre, récit oral} $\rightarrow$ on privilégie le \textbf{Perfekt} : \all{Der Direktor hat den Termin bekannt gegeben.}
\item Les verbes \all{sein}, \all{haben}, les verbes modaux et \all{werden} s'emploient au \textbf{Präteritum} même à l'oral : \all{ich war}, \all{ich hatte}, \all{ich konnte}.
\end{itemize}
\textbf{Exception :} dans un titre ou une accroche très vivante, le Präsens est possible : \all{Schüler gewinnen das Fußballturnier.} « Des élèves gagnent le tournoi de football. »
\end{propriete}

\begin{propriete}[Règle de traduction : l'ordre V2 et le registre]
Quel que soit le temps choisi, deux principes restent intangibles :
\begin{enumerate}
\item \textbf{Ordre V2 :} le verbe conjugué occupe toujours la deuxième position. Si une phrase française commence par un complément (« Hier, … », « Au Cameroun, … »), le sujet allemand passe \emph{après} le verbe : « Hier, la rédaction s'est réunie. » $\rightarrow$ \all{Gestern traf sich die Redaktion.}
\item \textbf{Registre :} un article de presse emploie un style sobre et informatif. On évite les tournures trop orales (\all{Also}, \all{echt}, \all{super}) et on préfère les verbes précis : \all{veröffentlichen} (publier), \all{berichten} (rapporter), \all{bekannt geben} (annoncer).
\end{enumerate}
\end{propriete}

\courssub{Tableau comparatif des formes verbales}

La correspondance entre les temps français et allemands n'est pas parfaite. Voici la carte de correspondance à connaître par cœur.

\begin{center}
\begin{tabularx}{\linewidth}{|X|X|X|}
\hline
\rowcolor{popblueL} Français & Deutsch & Exemple et remarque \\
\hline
Passé composé (oral, lettre, interview) & Perfekt & \all{Ich habe den Artikel gelesen.} « J'ai lu l'article. » \\
\hline
Passé composé / passé simple (article écrit) & Präteritum & \all{Die Schüler organisierten ein Fest.} « Les élèves ont organisé une fête. » \\
\hline
Imparfait (description, habitude) & Präteritum & \all{Es gab viele Leser.} « Il y avait beaucoup de lecteurs. » \\
\hline
Futur simple & Futur I & \all{Die Zeitung wird morgen erscheinen.} « Le journal paraîtra demain. » \\
\hline
Présent de vérité générale & Präsens & \all{Die Zeitung informiert die Schüler.} « Le journal informe les élèves. » \\
\hline
\end{tabularx}
\end{center}

\begin{popfigure}
\begin{tikzpicture}
\matrix (t) [matrix of nodes, nodes={anchor=center, minimum width=3cm}, column sep=0.5cm, row sep=0.3cm] {
FR & DE \\
Passé composé & Perfekt \\
Passé simple & Präteritum \\
Futur simple & Futur I \\
};
\draw (t-1-1.north west) rectangle (t-4-2.south east);
\draw[popblue, thick] (t-2-1.east) -- (t-2-2.west);
\draw[poporange, thick] (t-3-1.east) -- (t-3-2.west);
\draw[popgreen, thick] (t-4-1.east) -- (t-4-2.west);
\end{tikzpicture}
\legende{Correspondances des temps : le passé composé français a deux traductions possibles en allemand selon le registre (Perfekt à l'oral, Präteritum à l'écrit journalistique).}
\end{popfigure}

\begin{aretenir}[L'auxiliaire : sein ou haben ?]
Au \textbf{Perfekt}, le choix de l'auxiliaire est une source d'erreurs majeure :
\begin{itemize}
\item \textbf{sein} + participe II : verbes de \textbf{déplacement} ou de \textbf{changement d'état} : \all{gehen, kommen, fahren, aufstehen, sterben, passieren}. Ex. : \all{Die Delegation ist nach Yaoundé gefahren.} « La délégation est allée à Yaoundé. »
\item \textbf{haben} + participe II : tous les autres verbes, y compris les verbes \textbf{transitifs} et les verbes \textbf{réfléchis} : \all{Die Redaktion hat sich getroffen.} « La rédaction s'est réunie. »
\end{itemize}
Au \textbf{Präteritum}, pas d'auxiliaire : le problème disparaît ! C'est une raison supplémentaire de préférer le Präteritum dans les articles.
\end{aretenir}

\courssub{Pièges fréquents des francophones}

\begin{attention}[Faux amis et oublis classiques]
\begin{itemize}
\item \all{aktuell} ne signifie pas « actuel » au sens de « en fonction » mais « d'actualité, récent » : \all{die aktuellen Nachrichten} = « les nouvelles récentes ». Pour « le directeur actuel », dites \all{der jetzige Direktor}.
\item Ne jamais oublier le \all{es} expletif : « il y a » = \all{es gibt} ; « il pleut » = \all{es regnet}. Jamais *\all{Gibt viele Leser}.
\item \all{bekommen} = « recevoir », pas « devenir » (\all{werden}). \all{Ich habe einen Preis bekommen.} « J'ai reçu un prix. »
\item \all{die Redaktion} désigne à la fois la rédaction (l'équipe) et le bureau de rédaction ; elle est \textbf{féminine} : \all{Die Redaktion trifft sich.}
\item Verbe à particule : au Präteritum en V2, la particule va \textbf{en fin de phrase} : \all{Der Direktor gab den Termin bekannt.} (et non *\all{bekanntgab}).
\item Certains verbes régissent le \textbf{datif} : \all{helfen + Dativ} (\all{Der Direktor half der Redaktion.}), \all{danken + Dativ}, \all{gratulieren + Dativ}. Ne les mettez jamais à l'accusatif.
\end{itemize}
\end{attention}

\begin{exemplebox}[Exemples traduits commentés]
\begin{itemize}
\item « Le directeur a annoncé la date. » $\rightarrow$ \all{Der Direktor gab den Termin bekannt.} \emph{(Präteritum + particule en finale + accusatif „den Termin“)}
\item « Les journalistes ont interviewé le champion. » $\rightarrow$ \all{Die Journalisten interviewten den Champion.} \emph{(verbe faible : interviewte-n)}
\item « Un nouvel article paraîtra demain. » $\rightarrow$ \all{Ein neuer Artikel wird morgen erscheinen.} \emph{(Futur I : werden + infinitif en finale)}
\item « Le journal a été imprimé hier. » $\rightarrow$ \all{Die Zeitung wurde gestern gedruckt.} \emph{(passif au Präteritum : wurde + participe II)}
\end{itemize}
\end{exemplebox}

\courssub{Méthode : traduire en cinq étapes}

\begin{methode}[La méthode T.R.A.D.U.]
\begin{enumerate}
\item \textbf{T}ype de texte : article écrit ? $\rightarrow$ Präteritum. Dialogue ? $\rightarrow$ Perfekt.
\item \textbf{R}epérer le verbe et le conjuguer correctement (fort ou faible ?).
\item \textbf{A}ccorder l'ordre des mots : verbe en V2, particule et infinitif en finale.
\item \textbf{D}écliner : articles et adjectifs au bon cas (qui fait quoi ? Akkusativ ; à qui ? Dativ).
\item \textbf{U}ne relecture : majuscules aux noms, Umlaute, \all{es} expletif, auxiliaire sein/haben.
\end{enumerate}
\end{methode}

\begin{exoresolu}[Thème guidé : une phrase d'article]
Énoncé : traduisez en allemand, pour un article de journal scolaire : « Hier, les élèves de Première ont publié le premier numéro du journal. »
\tcblower
Solution détaillée :
\begin{enumerate}
\item \textbf{T}ype : article écrit $\rightarrow$ Präteritum.
\item \textbf{R} : „veröffentlichen“ (publier), verbe faible $\rightarrow$ \all{veröffentlichte} ; avec un pluriel : \all{veröffentlichten}.
\item \textbf{A} : la phrase commence par „Gestern“ $\rightarrow$ inversion du sujet : \all{Gestern veröffentlichten die Schüler der ersten Klasse …}
\item \textbf{D} : „le premier numéro“ : \all{die Nummer} est féminin $\rightarrow$ \all{die erste Nummer} (Akkusatif) ; „du journal“ = génitif : \all{der Schülerzeitung}.
\item \textbf{U} : relecture (majuscules, Umlaute, cas).
\end{enumerate}
Traduction finale : \all{Gestern veröffentlichten die Schüler der ersten Klasse die erste Nummer der Schülerzeitung.}
\end{exoresolu}

\courssub{Exercices de thème (FR $\rightarrow$ DE)}

\begin{exercice}[Thème : cinq phrases d'article]
Traduisez en allemand au temps approprié (texte écrit journalistique) :
\begin{enumerate}
\item Les élèves ont organisé une fête.
\item La rédaction a publié un article sur le football.
\item Le directeur a annoncé la date de la fête.
\item Il y avait beaucoup de visiteurs au marché de Yaoundé.
\item Le journal paraîtra demain.
\end{enumerate}
\end{exercice}

\begin{corrige}[Thème : cinq phrases]
\begin{enumerate}
\item \all{Die Schüler organisierten ein Fest.} \emph{(Präteritum, verbe faible)}
\item \all{Die Redaktion veröffentlichte einen Artikel über den Fußball.} \emph{(Akkusatif : „einen Artikel“)}
\item \all{Der Direktor gab den Termin der Feier bekannt.} \emph{(particule „bekannt“ en finale)}
\item \all{Es gab viele Besucher auf dem Markt in Yaoundé.} \emph{(„es gab“, ne pas oublier „es“)}
\item \all{Die Zeitung wird morgen erscheinen.} \emph{(Futur I : „wird“ + infinitif final)}
\end{enumerate}
\end{corrige}

\courssub{Exercices de version (DE $\rightarrow$ FR)}

En version, le travail est inverse : identifier le temps allemand, puis choisir le temps français naturel. Le Präteritum d'un article se rend généralement par le \textbf{passé composé} (ou le passé simple dans un récit soutenu) ; le Perfekt d'un dialogue se rend par le passé composé.

\begin{exercice}[Version : cinq phrases de presse]
Traduisez en français :
\begin{enumerate}
\item \all{Die Redaktion trifft sich morgen.}
\item \all{Die Schülerzeitung erschien gestern zum ersten Mal.}
\item \all{Ein Reporter hat den Bürgermeister von Douala interviewt.}
\item \all{Viele Leser schrieben Briefe an die Redaktion.}
\item \all{Der Artikel wurde von allen Schülern gelesen.}
\end{enumerate}
\end{exercice}

\begin{corrige}[Version : cinq phrases]
\begin{enumerate}
\item « La rédaction se réunit demain. » \emph{(Präsens à valeur de futur, comme en français)}
\item « Le journal des élèves a paru hier pour la première fois. » \emph{(Präteritum „erschien“ $\rightarrow$ passé composé)}
\item « Un reporter a interviewé le maire de Douala. » \emph{(Perfekt $\rightarrow$ passé composé)}
\item « Beaucoup de lecteurs ont écrit des lettres à la rédaction. » \emph{(„schrieben“, Präteritum de „schreiben“)}
\item « L'article a été lu par tous les élèves. » \emph{(passif : „wurde gelesen“ $\rightarrow$ « a été lu »)}
\end{enumerate}
\end{corrige}

\courssub{Traduction d'un texte court : version puis thème}

\begin{textebox}[Un article de journal scolaire]
\all{Gestern feierte die Schülerzeitung ‚Echos du Lycée‘ ihr erstes Jubiläum. Vor einem Jahr gründeten zehn Schüler der ersten Klasse die Zeitung. Zuerst hatten sie viele Probleme: Sie hatten keinen Computer und kein Geld für Papier. Aber der Direktor half der Redaktion, und eine Firma aus Douala gab Druckerpapier. Heute erscheint die Zeitung jeden Monat, und alle Schüler lesen sie gern. Die Redaktionsleiterin, Aminatou, sagte: „Wir wollen nächstes Jahr auch eine Seite auf Deutsch drucken.“ Die Schüler applaudierten laut. Es gab Musik, Kuchen und viele Fotos. Am Ende dankte der Direktor den jungen Journalisten für ihre Arbeit.}
\end{textebox}

\textbf{Glossaire :} \all{das Jubiläum, die Jubiläen} (l'anniversaire) ; \all{gründen} (fonder) ; \all{die Firma, die Firmen} (l'entreprise) ; \all{das Druckerpapier} (le papier d'impression) ; \all{die Redaktionsleiterin, -nen} (la rédactrice en chef) ; \all{drucken} (imprimer) ; \all{applaudieren} (applaudir) ; \all{danken + Dativ} (remercier).

\begin{exoresolu}[Version guidée du texte]
Énoncé : traduisez le texte ci-dessus en français.
\tcblower
Étapes : (1) le texte est un article $\rightarrow$ Präteritum dominant (\all{feierte, gründeten, half, gab, sagte, applaudierten, dankte}) $\rightarrow$ passé composé en français ; (2) un Präteritum de \all{haben} (\all{hatten}) dans un retour en arrière $\rightarrow$ passé composé ou imparfait ; (3) le discours direct reste au présent/futur.\par
Traduction : « Hier, le journal des élèves „Echos du Lycée“ a fêté son premier anniversaire. Il y a un an, dix élèves de Première ont fondé le journal. Au début, ils ont eu beaucoup de problèmes : ils n'avaient ni ordinateur ni argent pour le papier. Mais le directeur a aidé la rédaction, et une entreprise de Douala a donné du papier d'impression. Aujourd'hui, le journal paraît chaque mois, et tous les élèves le lisent avec plaisir. La rédactrice en chef, Aminatou, a déclaré : „Nous voulons l'année prochaine imprimer aussi une page en allemand.“ Les élèves ont applaudi fort. Il y a eu de la musique, du gâteau et beaucoup de photos. À la fin, le directeur a remercié les jeunes journalistes de leur travail. »
\end{exoresolu}

\begin{exercice}[Thème : texte court à traduire]
Traduisez ce texte en allemand (style article, Präteritum) : « La semaine dernière, notre lycée a organisé une journée de la presse. Les élèves ont écrit des articles et ont interviewé le directeur. La rédaction a imprimé cent exemplaires du journal. Tous les visiteurs ont lu le journal avec intérêt. Le directeur a félicité les jeunes journalistes. »
\end{exercice}

\begin{corrige}[Thème : texte court]
\all{Letzte Woche organisierte unser Gymnasium einen Pressetag. Die Schüler schrieben Artikel und interviewten den Direktor. Die Redaktion druckte hundert Exemplare der Zeitung. Alle Besucher lasen die Zeitung mit Interesse. Der Direktor gratulierte den jungen Journalisten.}\par
Points de vigilance : \all{schrieben} (Präteritum fort de \all{schreiben}), \all{lasen} (Präteritum fort de \all{lesen}), \all{gratulieren + Dativ} (\all{den jungen Journalisten}), \all{hundert Exemplare} (Akkusatif pluriel).
\end{corrige}

\begin{aretenir}[Bilan de la section]
Pour réussir thème et version de presse : (1) identifiez le type de texte pour choisir Präteritum ou Perfekt ; (2) placez le verbe en V2 et les particules en finale ; (3) méfiez-vous des faux amis (\all{aktuell}, \all{bekommen}) ; (4) n'oubliez jamais \all{es} expletif ; (5) relisez : majuscules, cas, auxiliaires. Vous êtes maintenant prêts pour la production guidée de la section 6.
\end{aretenir}

\coursec{Production guidée \& relecture}

Dans la section précédente, nous avons traduit des phrases et des textes courts de presse, du français vers l'allemand (thème) et de l'allemand vers le français (version). Traduire, c'est déjà produire : on choisit un temps, on place le verbe en deuxième position, on accorde les articles. Nous allons maintenant franchir une étape supplémentaire : \textbf{écrire nous-mêmes un article de journal scolaire complet}, de la première idée jusqu'à la relecture finale. C'est exactement la compétence demandée à l'examen : produire un texte journalistique clair, correct et bien présenté. Suivez les sept étapes de la méthode.

\courssub{Vue d'ensemble : les sept étapes de la production}

Rédiger un article n'est pas un acte magique : c'est un \textbf{processus en étapes}. Le schéma suivant résume le parcours complet, du choix du sujet à la révision finale.

\begin{popfigure}
\begin{tikzpicture}[node distance=0.7cm]
\node[draw=popblue, fill=popblueL, rounded corners, align=center, minimum width=5cm] (start) {1. Choix du sujet};
\node[draw=popblue, fill=popblueL, rounded corners, align=center, minimum width=5cm, below of=start] (lead) {2. Lead (5 W)};
\node[draw=poporange, fill=poporangeL, rounded corners, align=center, minimum width=5cm, below of=lead] (body) {3. Corps (3 paragraphes)};
\node[draw=poporange, fill=poporangeL, rounded corners, align=center, minimum width=5cm, below of=body] (ank) {4. Annonce (Ankündigung)};
\node[draw=popgreen, fill=popgreenL, rounded corners, align=center, minimum width=5cm, below of=ank] (layout) {5. Mise en page};
\node[draw=popgreen, fill=popgreenL, rounded corners, align=center, minimum width=5cm, below of=layout] (check) {6. Grille d'auto-correction};
\node[draw=poppurple, fill=poppurpleL, rounded corners, align=center, minimum width=5cm, below of=check] (pair) {7. Relecture en binôme};
\draw[->, thick, popdark] (start) -- (lead);
\draw[->, thick, popdark] (lead) -- (body);
\draw[->, thick, popdark] (body) -- (ank);
\draw[->, thick, popdark] (ank) -- (layout);
\draw[->, thick, popdark] (layout) -- (check);
\draw[->, thick, popdark] (check) -- (pair);
\end{tikzpicture}
\legende{Le processus de rédaction d'un article de journal scolaire en sept étapes.}
\end{popfigure}

\begin{aretenir}[Le principe d'or]
On ne rédige jamais un article « d'un coup ». On \textbf{planifie} (étapes 1 à 4), on \textbf{présente} (étape 5), puis on \textbf{corrige} (étapes 6 et 7). À l'examen, consacre environ un tiers du temps à la planification et à la relecture : c'est là que se gagnent les points.
\end{aretenir}

\courssub{Étape 1 : Choisir le sujet}

Le sujet doit être \textbf{concret, récent et intéressant pour les lecteurs} (les élèves de ton école). Trois sujets types, très fréquents dans les journaux scolaires camerounais :

\begin{itemize}
\item un \textbf{événement scolaire} : \all{das Schulfest} (la fête de l'école), \all{die Theateraufführung, -en} (la représentation théâtrale), la visite d'un député au lycée de Yaoundé ;
\item un \textbf{résultat sportif} : \all{das Fußballspiel, -e} (le match de football), la victoire de l'équipe du lycée contre le lycée voisin de Douala ;
\item un \textbf{projet écologique} : \all{das Umweltprojekt, -e} (le projet pour l'environnement), la plantation d'arbres dans la cour, le club de recyclage.
\end{itemize}

\begin{methode}[Choisir un bon sujet en trois questions]
\begin{enumerate}
\item \textbf{Qui est concerné ?} (les élèves, les professeurs, les parents)
\item \textbf{Qu'est-ce qui s'est passé de nouveau ?} (un fait précis, pas une opinion vague)
\item \textbf{Puis-je citer quelqu'un ?} (un témoin, un organisateur : indispensable pour l'interview de l'étape 3)
\end{enumerate}
Si tu réponds « oui » aux trois questions, ton sujet est bon.
\end{methode}

\courssub{Étape 2 : Rédiger le Lead (les 5 W) au Präteritum}

Le \cle{Lead} (en allemand aussi \all{der Lead} ou \all{der Vorspann}) est le premier paragraphe de l'article : en deux ou trois phrases, il répond aux \textbf{5 W}.

\begin{center}
\begin{tabularx}{\linewidth}{|l|X|X|}
\hline
\rowcolor{popblueL} \textbf{W} & \textbf{Question allemande} & \textbf{Exemple (Schulfest)} \\
\hline
\textbf{W}er ? & Qui ? & Die Schüler der Klasse 1ère A \\
\hline
\textbf{W}as ? & Quoi ? & organisierten ein Schulfest \\
\hline
\textbf{W}ann ? & Quand ? & am Samstag, dem 15. März \\
\hline
\textbf{W}o ? & Où ? & im Schulhof des Gymnasiums Yaoundé \\
\hline
\textbf{W}arum ? & Pourquoi ? & um Geld für die Bibliothek zu sammeln \\
\hline
\end{tabularx}
\end{center}

Le Lead se rédige de préférence au \textbf{Präteritum}, le temps du récit écrit journalistique.

\begin{exemplebox}[Un Lead complet]
\all{Die Schüler der Klasse 1ère A organisierten am Samstag, dem 15. März, ein Schulfest im Schulhof des Gymnasiums Yaoundé. Sie wollten Geld für die neue Bibliothek sammeln. Mehr als 300 Gäste kamen.}

« Les élèves de la classe de Première A ont organisé samedi 15 mars une fête scolaire dans la cour du lycée de Yaoundé. Ils voulaient collecter de l'argent pour la nouvelle bibliothèque. Plus de 300 invités sont venus. »
\end{exemplebox}

\begin{attention}[Präteritum ou Perfekt ?]
Compare : « Les élèves ont chanté. » se traduit \all{Die Schüler sangen.} (Präteritum, style écrit, article de presse) ou \all{Die Schüler haben gesungen.} (Perfekt, style oral, conversation, lettre à un ami). Dans un \textbf{article de journal}, privilégie le Präteritum ; dans les \textbf{citations} (paroles rapportées directement), les personnes interrogées parlent naturellement au Perfekt. Vérifie aussi les formes : verbes faibles en \all{-te} (\all{organisierte}), verbes forts avec changement de voyelle (\all{sang}, \all{kam}), et au Perfekt, le bon choix de l'auxiliaire : \all{haben} avec les verbes transitifs (\all{haben gesungen}), \all{sein} avec les verbes de mouvement ou de changement d'état (\all{sind gekommen}).
\end{attention}

\courssub{Étape 3 : Développer le corps de l'article (3 paragraphes)}

Le corps de l'article développe le Lead en trois paragraphes :

\begin{enumerate}
\item \textbf{Paragraphe 1 — les faits} : déroulement de l'événement, chiffres, détails concrets.
\item \textbf{Paragraphe 2 — les citations} : une mini-interview (une question, une ou deux réponses entre guillemets allemands „ … “).
\item \textbf{Paragraphe 3 — la suite} : conséquences, remerciements, perspectives.
\end{enumerate}

\begin{propriete}[Rapporter une citation : la place du verbe]
Deux constructions possibles :
\begin{itemize}
\item Citation d'abord, verbe de parole ensuite (ordre inversé) : \all{„Das Fest war ein großer Erfolg“, sagte die Schülerin Amina.} (« La fête a été un grand succès », a dit l'élève Amina.)
\item Verbe de parole d'abord, deux-points, puis la citation : \all{Die Schülerin Amina sagte: „Das Fest war ein großer Erfolg.“}
\end{itemize}
Dans les deux cas, la citation garde l'ordre V2 (verbe conjugué en deuxième position) ; elle s'ouvre avec „ et se ferme avec “.
\end{propriete}

\courssub{Étape 4 : Rédiger une Ankündigung (annonce) liée}

Après l'article, le journal scolaire publie souvent une \cle{Ankündigung} (annonce) : un texte très bref qui invite les lecteurs à un événement futur. Elle contient obligatoirement : \textbf{la date, le lieu, l'heure et la manière de s'inscrire}. Elle se rédige au présent ou au futur.

\begin{exemplebox}[Modèle d'annonce]
\all{Ankündigung: Das nächste Schulfest findet am 20. Juni um 15 Uhr im Schulhof statt. Alle Schüler sind herzlich eingeladen! Anmeldung bis zum 10. Juni bei Frau Mbarga (Raum 12).}

« Annonce : la prochaine fête scolaire aura lieu le 20 juin à 15 h dans la cour de l'école. Tous les élèves sont chaleureusement invités ! Inscription jusqu'au 10 juin auprès de Mme Mbarga (salle 12). »
\end{exemplebox}

Mots utiles : \all{findet ... statt} (a lieu ; attention, verbe à particule séparable : la particule \all{statt} va en fin de phrase), \all{die Anmeldung, -en} (l'inscription), \all{herzlich eingeladen} (chaleureusement invité), \all{der Eintritt ist frei} (l'entrée est libre/gratuite).

\courssub{Étape 5 : La mise en page}

Un article de journal se reconnaît à sa \textbf{forme} autant qu'à son contenu. Éléments obligatoires :

\begin{itemize}
\item \textbf{le titre} (\all{die Überschrift, -en}) : court, accrocheur, avec majuscules aux noms ;
\item \textbf{l'auteur et la date} : \all{von Marie Ngo Bell, 18. März 2025} ;
\item \textbf{les intertitres} (\all{die Zwischenüberschriften}) pour séparer les parties ;
\item des paragraphes courts et aérés.
\end{itemize}

\courssub{Modèle complet : lire avant d'écrire}

Avant de produire toi-même, observe ce modèle intégral. Repère le Lead, les trois paragraphes, la citation et l'annonce.

\begin{textebox}[Unser Schulfest — ein voller Erfolg]
\all{von Marie Ngo Bell, 18. März 2025}

\all{Die Schüler der Klasse 1ère A organisierten am Samstag, dem 15. März, ein Schulfest im Schulhof des Gymnasiums Yaoundé. Sie wollten Geld für die neue Bibliothek sammeln. Mehr als 300 Gäste kamen.}

\textbf{\all{Ein buntes Programm}}

\all{Am Nachmittag gab es Musik, Theater und traditionelle Tänze. Die Schüler verkauften auch Essen und Getränke: Puff-Puff, gegrillten Mais und frischen Saft. Um 17 Uhr spielte die Schulband, und alle Gäste tanzten.}

\textbf{\all{„Wir sind sehr stolz“}}

\all{„Das Fest war ein großer Erfolg“, sagte die Schülerin Amina. „Wir haben viel Geld gesammelt, und alle haben geholfen.“ Auch der Direktor, Herr Essono, war zufrieden: „Diese Schüler zeigen, dass Jugendliche viel erreichen können.“}

\textbf{\all{Und jetzt?}}

\all{Die Klasse 1ère A gab dem Bibliotheksprojekt 150 000 FCFA. Im Juni beginnt der Bau der neuen Bibliothek.}

\textbf{\all{Ankündigung:}} \all{Das nächste Schulfest findet am 20. Juni um 15 Uhr im Schulhof statt. Anmeldung bis zum 10. Juni bei Frau Mbarga (Raum 12). Der Eintritt ist frei!}
\end{textebox}

\textbf{Glossaire :} \all{der Schulhof, -höfe} (la cour de l'école) ; \all{sammeln} (collecter) ; \all{der Gast, Gäste} (l'invité) ; \all{gegrillter Mais} (le maïs grillé) ; \all{der Saft, Säfte} (le jus) ; \all{stolz} (fier) ; \all{erreichen} (atteindre, réussir) ; \all{der Bau, -ten} (la construction) ; \all{der Eintritt, -e} (l'entrée).

\textbf{Questions de compréhension :} 1. Wer organisierte das Fest? 2. Warum organisierten die Schüler das Fest? 3. Was sagte Amina? 4. Wann findet das nächste Fest statt?

\textbf{Réponses attendues :} 1. \all{Die Schüler der Klasse 1ère A.} 2. \all{Sie wollten Geld für die neue Bibliothek sammeln.} 3. \all{Das Fest war ein großer Erfolg; sie haben viel Geld gesammelt, und alle haben geholfen.} 4. \all{Am 20. Juni um 15 Uhr.}

\courssub{Étape 6 : L'auto-correction avec la grille KORREKTUR}

La relecture n'est pas une option : c'est là qu'on gagne des points. Utilise la grille mnémotechnique \cle{KORREKTUR} : chaque lettre correspond à un point de contrôle.

\begin{methode}[La checklist KORREKTUR]
\begin{itemize}
\item \textbf{K — Kasus} : les cas sont-ils corrects ? (\all{für die Bibliothek} : accusatif après \all{für})
\item \textbf{O — Orthographe} : Umlaute (ä, ö, ü), ß, pas de fautes de copie.
\item \textbf{R — Registre} : ton journalistique, pas de langage SMS ni de familiarités.
\item \textbf{R — Reihenfolge} : ordre des mots, verbe conjugué en position 2 (V2).
\item \textbf{E — Einheitlichkeit} : unité des temps (Präteritum pour le récit, présent/futur pour l'annonce).
\item \textbf{K — Kommas} : virgules devant \all{dass}, \all{weil}, \all{um ... zu}, et entre propositions.
\item \textbf{T — Tempus} : terminaisons du Präteritum (\all{-te}, \all{-test}, \all{-ten}) et participes passés corrects.
\item \textbf{U — Umfang} : longueur respectée (120 à 150 mots à l'examen).
\item \textbf{R — Reread} : relire une dernière fois à voix haute, lentement.
\end{itemize}
\end{methode}

\begin{center}
\begin{tabularx}{\linewidth}{|X|X|X|}
\hline
\rowcolor{poporangeL} \textbf{Erreur fréquente} & \textbf{Faux} & \textbf{Correct} \\
\hline
Majuscule oubliée & \all{die schüler sangen} & \all{Die Schüler sangen} \\
\hline
Verbe mal placé & \all{Am Samstag die Schüler organisierten...} & \all{Am Samstag organisierten die Schüler...} \\
\hline
Terminaison fausse & \all{sie singte} & \all{sie sang} (verbe fort) \\
\hline
Participe passé & \all{sie haben gesingt} & \all{sie haben gesungen} \\
\hline
Cas après \all{für} & \all{für der Bibliothek} & \all{für die Bibliothek} \\
\hline
Faux ami & \all{aktuel} pour « actuel » & \all{aktuell} = « d'actualité » ; « actuel » = \all{gegenwärtig} \\
\hline
\end{tabularx}
\end{center}

\begin{attention}[Les trois pièges préférés des francophones]
1. \textbf{La majuscule} : en allemand, TOUS les noms prennent une majuscule, pas seulement en début de phrase. 2. \textbf{L'ordre des mots} : si la phrase commence par un complément (\all{Am Samstag}), le verbe reste en deuxième position et le sujet suit : \all{Am Samstag kamen 300 Gäste.} 3. \textbf{Le temps} : ne mélange pas Präteritum et Perfekt dans le récit ; choisis le Präteritum et tiens-toi-y.
\end{attention}

\courssub{Étape 7 : L'échange en binôme et la révision finale}

Après l'auto-correction, échange ton article avec un camarade. Un œil neuf voit des erreurs que l'auteur ne voit plus.

\begin{experience}[Relecture croisée en binôme]
\begin{enumerate}
\item Échangez vos articles avec votre voisin.
\item Lisez l'article de l'autre avec la grille KORREKTUR : soulignez au crayon chaque problème et notez la lettre correspondante (K, O, R...) dans la marge.
\item Posez deux questions de lecteur : « Qu'est-ce qui n'est pas clair ? Quel détail manque ? »
\item Rendez l'article à son auteur, qui rédige la \textbf{version finale} propre.
\end{enumerate}
Durée conseillée : 10 minutes de relecture + 10 minutes de réécriture.
\end{experience}

\begin{exoresolu}[Corriger un extrait d'article]
Énoncé : Voici le Lead rédigé par un élève. Trouve et corrige les cinq erreurs.

\all{Am freitag die schüler der Klasse 1ère A hat organisiert ein Schulfest im Schulhof. Sie wollten Geld für der Bibliothek sammeln, weil die alte Bibliothek zu klein ist.}

\tcblower

Solution détaillée :
\begin{enumerate}
\item \all{freitag} → \all{Freitag} : majuscule obligatoire aux noms (jours de la semaine). [O]
\item \all{die schüler} → \all{die Schüler} : majuscule au nom. [O]
\item \all{Am Freitag die Schüler ... hat organisiert} → \all{Am Freitag organisierten die Schüler} : verbe conjugué en position 2 après le complément initial, et Präteritum (récit journalistique) au lieu du Perfekt. [R, T]
\item \all{für der Bibliothek} → \all{für die Bibliothek} : \all{für} exige l'accusatif ; \all{die Bibliothek} est féminin, accusatif \all{die}. [K]
\item La subordonnée avec \all{weil} est correcte (verbe \all{ist} en fin de proposition, virgule présente) : à vérifier, mais sans erreur ici. [K comme Kommas : vérifié]
\end{enumerate}
Version corrigée : \all{Am Freitag organisierten die Schüler der Klasse 1ère A ein Schulfest im Schulhof. Sie wollten Geld für die Bibliothek sammeln, weil die alte Bibliothek zu klein ist.}

« Vendredi, les élèves de la classe de Première A ont organisé une fête scolaire dans la cour. Ils voulaient collecter de l'argent pour la bibliothèque, car l'ancienne bibliothèque est trop petite. »
\end{exoresolu}

\begin{exercice}[À ton tour : rédige ton article]
Choisis l'un des trois sujets suivants et rédige un article complet de 120 à 150 mots en suivant les sept étapes :
\begin{enumerate}
\item Ton équipe de football a gagné un match contre un lycée voisin.
\item Ta classe a planté des arbres dans la cour de l'école.
\item Le club de théâtre a joué une pièce devant les parents.
\end{enumerate}
N'oublie pas : titre, auteur, date, Lead (5 W) au Präteritum, trois paragraphes avec une citation, une annonce finale, puis l'auto-correction KORREKTUR et la relecture en binôme.
\end{exercice}

\begin{corrige}[Exercice — éléments de correction]
Il n'existe pas une seule bonne réponse, mais voici un canevas attendu pour le sujet 2 :
\begin{itemize}
\item Titre : \all{Neue Bäume für unsere Schule} ; auteur et date.
\item Lead : \all{Die Schüler der Klasse 1ère A pflanzten am Dienstag zwanzig Bäume im Schulhof. Sie wollten die Umwelt schützen und Schatten für die Pause schaffen.}
\item Corps : déroulement (qui a creusé, arrosé) ; citation : \all{„Es war harte Arbeit, aber wir sind stolz“, sagte der Schüler Jean.} ; suite : un arrosage chaque semaine.
\item Annonce : \all{Das nächste Umweltprojekt beginnt am 5. Mai. Anmeldung bei Herrn Fotso.}
\item Grille KORREKTUR appliquée : majuscules aux noms, V2, Präteritum régulier (\all{pflanzte}, \all{sagte}), accusatif après \all{für}.
\end{itemize}
\end{corrige}

\begin{aretenir}[Astuce examen : trois Leads types à mémoriser]
Pour gagner du temps le jour de l'épreuve, apprends par cœur ces trois gabarits de Lead :
\begin{itemize}
\item \textbf{Événement} : \all{Am [Tag] organisierten die Schüler [was] im Schulhof / in der Aula.}
\item \textbf{Résultat sportif} : \all{Am [Tag] gewann die Fußballmannschaft unserer Schule das Spiel gegen [Gegner] mit [Ergebnis].}
\item \textbf{Annonce} : \all{Am [Tag] findet [Ereignis] um [Uhrzeit] in [Ort] statt. Alle sind eingeladen!}
\end{itemize}
Remplace simplement les éléments entre crochets par les informations du sujet.
\end{aretenir}

\begin{savaistu}[La production écrite à l'examen]
Aux épreuves d'allemand (examen de classe de Première, puis baccalauréat), on peut te demander la rédaction d'un article de journal scolaire de \textbf{120 à 150 mots} en un temps limité. Les correcteurs évaluent quatre critères : le respect du type de texte (titre, Lead, annonce), la correction grammaticale (temps, cas, ordre des mots), le vocabulaire et la présentation. Entraîne-toi régulièrement \textbf{avec un minuteur} : 5 minutes de planification, 20 minutes de rédaction, 5 minutes de relecture KORREKTUR. C'est exactement le rythme des sept étapes de cette leçon !
\end{savaistu}

\newpage

\coursec{Activité d'intégration}

\courssub{Objectifs de l'activité}

À l'issue de cette activité d'intégration, l'élève sera capable de :
\begin{itemize}
\item rédiger en allemand un \cle{article} de journal scolaire (environ 130 mots) sur un événement réel du lycée, au \cle{Perfekt} et au \cle{Präteritum} ;
\item rédiger une \cle{interview} fictive (3 à 4 échanges) avec des questions correctement formulées (ordre des mots, verbe en position 2) ;
\item rédiger une \cle{annonce} brève et attractive pour une activité à venir ;
\item traduire un texte de l'allemand vers le français (\cle{version}) et du français vers l'allemand (\cle{thème}) ;
\item relire et corriger sa production à l'aide d'une grille (orthographe, majuscules, ordre des mots, cas) ;
\item travailler en groupe, répartir les tâches et présenter le résultat à l'oral.
\end{itemize}

\courssub{Situation}

Le lycée bilingue de Yaoundé prépare la fin de l'année scolaire. La direction a décidé de publier un numéro spécial du journal scolaire \emph{„Die Löwenstimme“} (La Voix des Lions), entièrement \textbf{bilingue allemand-français}, qui sera distribué aux élèves, aux parents et aux partenaires allemands du lycée lors de la journée portes ouvertes. Le club d'allemand est chargé de la page 4, consacrée à la vie du lycée.

Voici la note de service du rédacteur en chef :

\begin{textebox}[Note de service du rédacteur en chef]
Liebe Mitglieder des Deutschklubs,\\
für die nächste Ausgabe unserer Schülerzeitung „Die Löwenstimme“ brauchen wir eine komplette Seite über das Schulleben. Eure Seite muss drei Texte enthalten:\\
1. einen Artikel (ca. 130 Wörter) über ein Ereignis an unserem Lyzeum,\\
2. ein Interview (3 bis 4 Fragen) mit einem Schüler oder einer Lehrerin,\\
3. eine Anzeige für eine nächste Aktivität.\\
Der Artikel muss auf Französisch übersetzt werden, die Anzeige auf Deutsch. Abgabe: Freitag, 17 Uhr. Viel Erfolg!\\
Euer Chefredakteur
\end{textebox}

\begin{definition}[Glossaire de la note de service]
der Deutschklub, -s : le club d'allemand ; die Ausgabe, -n : le numéro (d'un journal) ; die Schülerzeitung, -en : le journal scolaire ; das Ereignis, -se : l'événement ; das Lyzeum, Lyzeen : le lycée ; die Anzeige, -n : l'annonce ; übersetzen : traduire ; die Abgabe, -n : la remise (du travail).
\end{definition}

\courssub{Consignes}

Travail en groupes de 4 élèves. Durée : 2 heures (répartition conseillée : 20 min de préparation, 60 min de rédaction et traduction, 20 min de relecture, 20 min de présentation).

\begin{enumerate}
\item \textbf{Choix du sujet (10 min).} En groupe, choisissez un événement récent de votre lycée : la fête de la jeunesse, un match de football, la semaine de la culture, une visite d'étude, la cérémonie de remise des prix. Notez 5 mots-clés allemands sur cet événement.
\item \textbf{Rédaction de l'article (25 min).} Rédigez un article d'environ 130 mots en allemand. Il doit contenir : un \cle{titre} accrocheur, une \cle{introduction} (qui ? quoi ? quand ? où ?), un \cle{développement} (déroulement, réactions) et une \cle{conclusion}. Employez le Perfekt pour les événements récents et le Präteritum pour \all{sein}, \all{haben} et les verbes modaux.
\item \textbf{Rédaction de l'interview (15 min).} Imaginez une interview de 3 à 4 questions avec un élève ou un professeur du lycée. Attention à la place du verbe dans les questions (avec ou sans mot interrogatif).
\item \textbf{Rédaction de l'annonce en français puis en allemand (15 min).} Rédigez d'abord l'annonce en français (une prochaine activité : tournoi, sortie, concours), puis traduisez-la en allemand (\cle{thème}).
\item \textbf{Traduction de l'article (20 min).} Traduisez votre article allemand en français (\cle{version}). Restez fidèles au sens, mais rédigez un français naturel.
\item \textbf{Relecture croisée (15 min).} Échangez votre page avec un autre groupe. Corrigez à l'aide de la grille KORREKTUR : majuscules des noms, place du verbe (V2), terminaisons du Perfekt et du Präteritum, cas après les prépositions, orthographe (ä, ö, ü, ß).
\item \textbf{Mise en page et présentation (20 min).} Disposez les trois textes sur une feuille A3 (ou outil numérique) avec titres et rubriques. Chaque groupe présente sa page à la classe en 3 minutes, en allemand.
\end{enumerate}

\courssub{Ressources à mobiliser}

\begin{aretenir}[Boîte à outils de la page de journal]
\textbf{Grammaire :}
\begin{itemize}
\item \cle{Perfekt} : auxiliaire \all{haben} ou \all{sein} + participe II à la fin : \all{Wir haben ein Turnier organisiert.} / \all{Viele Gäste sind gekommen.}
\item \cle{Präteritum} : pour \all{sein} (\all{war}), \all{haben} (\all{hatte}) et les modaux (\all{konnte}, \all{musste}, \all{wollte}).
\item \cle{Ordre des mots} : verbe conjugué en position 2 ; dans les questions sans mot interrogatif, verbe en position 1 : \all{Hast du das Spiel gesehen ?}
\end{itemize}
\textbf{Lexique de la presse :} der Artikel, - ; die Schlagzeile, -n (le titre) ; das Interview, -s ; die Anzeige, -n ; der Redakteur, -e ; die Ausgabe, -n ; berichten (rapporter) ; stattfinden (avoir lieu, séparable : \all{findet statt}, \all{hat stattgefunden}).\\
\textbf{Structures utiles :} \all{Am Freitag hat ... stattgefunden.} ; \all{Wir haben Herrn N. gefragt: ...} ; \all{Kommt alle am Samstag um 15 Uhr in die Aula !} (impératif pour l'annonce).
\end{aretenir}

\begin{attention}[Pièges à éviter]
\begin{itemize}
\item Ne dites pas \all{*Das Fest hat gefunden*} : \all{stattfinden} est séparable et se conjugue \all{hat stattgefunden}.
\item Après un complément de temps en tête de phrase, le sujet passe après le verbe : \all{Am Montag \textbf{waren} wir im Stadion.} et non \all{*Am Montag wir waren...}.
\item Majuscule à tous les noms : \all{das Spiel, die Schule, der Preis}.
\item En version, ne traduisez pas mot à mot : \all{Es gab viele Gäste} se traduit naturellement par « il y avait beaucoup d'invités ».
\end{itemize}
\end{attention}

\courssub{Proposition de production et corrigé commenté}

\begin{exoresolu}[Modèle de page de journal — groupe de référence]
Énoncé : produire la page 4 de \emph{„Die Löwenstimme“} selon les consignes (article + interview + annonce, avec version de l'article et thème de l'annonce).
\tcblower
\textbf{1. Artikel (article) :}\\[2pt]
\textbf{Fußballfest am Lyzeum: Unsere Mannschaft gewinnt den Pokal}\\
Am Freitagnachmittag hat auf dem Sportplatz unseres Lyzeums das große Fußballturnier stattgefunden. Sechs Mannschaften aus Yaoundé haben teilgenommen. Das Wetter war schön, und viele Schüler und Lehrer sind gekommen, um die Spieler zu unterstützen.\\
Unser Team hat zuerst gegen das Lyzeum von Nkol-Eton gespielt und 2:1 gewonnen. Im Finale war das Spiel sehr schwer, aber unser Kapitän Arnaud hat in der letzten Minute ein Tor geschossen. Endstand: 1:0 für uns!\\
Der Direktor hat den Pokal übergeben und alle Spieler gratuliert. „Ich bin sehr stolz auf euch“, hat er gesagt. Es war ein unvergesslicher Tag für unser Lyzeum.\\[4pt]
\textit{Commentaire :} titre accrocheur (die Schlagzeile) ; introduction avec qui/quoi/quand/où ; Perfekt pour les actions (\all{hat stattgefunden}, \all{haben teilgenommen}, \all{hat geschossen}) ; Präteritum pour \all{war} ; verbe en position 2 après \all{Am Freitagnachmittag} et \all{Im Finale} ; participe II à la fin ; discours rapporté avec guillemets allemands.\\[6pt]
\textbf{Version (traduction française) :}\\
« Fête du football au lycée : notre équipe remporte la coupe »\\
Vendredi après-midi, le grand tournoi de football a eu lieu sur le terrain de sport de notre lycée. Six équipes de Yaoundé ont participé. Il faisait beau et de nombreux élèves et professeurs sont venus soutenir les joueurs. Notre équipe a d'abord joué contre le lycée de Nkol-Eton et a gagné 2 à 1. En finale, le match était très difficile, mais notre capitaine Arnaud a marqué un but à la dernière minute. Score final : 1 à 0 pour nous ! Le directeur a remis la coupe et a félicité tous les joueurs. « Je suis très fier de vous », a-t-il dit. C'était un jour inoubliable pour notre lycée.\\[6pt]
\textbf{2. Interview :}\\
\all{Die Löwenstimme: Arnaud, wie war das Finale für dich?}\\
\all{Arnaud: Es war sehr schwer, aber wir haben als Team gekämpft.}\\
\all{Die Löwenstimme: Wann hast du mit dem Fußball angefangen?}\\
\all{Arnaud: Mit zehn Jahren, in meinem Quartier in Yaoundé.}\\
\all{Die Löwenstimme: Was ist dein nächstes Ziel?}\\
\all{Arnaud: Wir wollen nächstes Jahr den Pokal wieder gewinnen!}\\[4pt]
\textit{Commentaire :} questions avec mot interrogatif (verbe en position 2) ; verbe séparable \all{anfangen} : \all{hast ... angefangen}.\\[6pt]
\textbf{3. Annonce (thème FR $\rightarrow$ DE) :}\\
Français : « Grande journée portes ouvertes ! Venez nombreux samedi à 9 h à la salle des fêtes du lycée. Au programme : exposition du club d'allemand, concert et match de football. Entrée libre ! »\\
Deutsch : \all{Großer Tag der offenen Tür! Kommt alle am Samstag um 9 Uhr in die Aula des Lyzeums. Programm: eine Ausstellung des Deutschklubs, ein Konzert und ein Fußballspiel. Der Eintritt ist frei!}\\[4pt]
\textit{Commentaire :} impératif pluriel \all{kommt} ; heure avec \all{um} ; \all{Tag der offenen Tür} = journée portes ouvertes (expression figée).
\end{exoresolu}

\courssub{Bilan de l'activité}

Cette activité d'intégration a permis de mobiliser l'ensemble des compétences du chapitre dans une tâche authentique et motivante :
\begin{itemize}
\item \textbf{Compréhension de l'écrit} : lecture et exploitation de la note de service en allemand ;
\item \textbf{Expression écrite} : production des trois genres journalistiques (article, interview, annonce) ;
\item \textbf{Grammaire réinvestie} : Perfekt et Präteritum dans un contexte réel, ordre des mots (V2, participe final), impératif ;
\item \textbf{Traduction} : version (article DE $\rightarrow$ FR) et thème (annonce FR $\rightarrow$ DE), avec attention à la fidélité et au naturel ;
\item \textbf{Relecture} : correction croisée avec la grille KORREKTUR, qui développe l'autonomie et l'auto-évaluation ;
\item \textbf{Expression orale} : présentation de la page devant la classe.
\end{itemize}

\begin{methode}[Grille d'évaluation simplifiée]
\begin{enumerate}
\item \textbf{Contenu} (4 pts) : les trois textes sont présents, complets et adaptés au genre.
\item \textbf{Grammaire} (6 pts) : Perfekt/Präteritum corrects, verbe en position 2, cas et accords.
\item \textbf{Lexique} (3 pts) : vocabulaire de la presse et de la vie scolaire bien employé.
\item \textbf{Traduction} (4 pts) : fidélité au sens, allemand et français naturels.
\item \textbf{Présentation} (3 pts) : mise en page soignée, titres, présentation orale claire.
\end{enumerate}
\end{methode}

\begin{savaistu}[Le journal scolaire en Allemagne]
En Allemagne, presque chaque école possède sa \all{Schülerzeitung}, rédigée entièrement par les élèves avec l'aide d'un professeur-conseil. Il existe même des concours nationaux qui récompensent les meilleurs journaux scolaires : un beau modèle pour \emph{„Die Löwenstimme“} de Yaoundé !
\end{savaistu}

\newpage

\coursec{Méthodes et exercices résolus}

\courssub{Méthode 1 : Lire et comprendre un texte de presse}

\begin{methode}[Comprendre un article de presse en 5 étapes]
\begin{enumerate}
\item \textbf{Lire le titre et le sous-titre} : ils annoncent le thème (\all{das Thema}) et souvent l'opinion du journaliste.
\item \textbf{Repérer les 5 W} : \all{Wer ?} (qui), \all{Was ?} (quoi), \all{Wo ?} (où), \all{Wann ?} (quand), \all{Warum ?} (pourquoi). Souligne-les dans le texte.
\item \textbf{Identifier la structure} : un article commence par l'essentiel (le \all{Lead}), puis donne les détails. La première phrase résume souvent tout l'article.
\item \textbf{Deviner les mots inconnus} par le contexte et la famille de mots : \all{der Schüler} → \all{die Schülerzeitung} (le journal des élèves).
\item \textbf{Répondre en phrases complètes}, en reprenant les mots de la question, sans recopier tout le texte.
\end{enumerate}
\textbf{Réflexe} : à l'examen, les questions de compréhension suivent l'ordre du texte. La réponse à la question 2 se trouve après celle de la question 1.
\end{methode}

\begin{exoresolu}[Compréhension de l'écrit — type examen]
\textbf{Texte :} \all{Seit diesem Schuljahr erscheint an unserem Gymnasium in Yaoundé eine neue Schülerzeitung. Sie heißt „Der Pausenfrosch“ und erscheint jeden Monat. Zwanzig Schülerinnen und Schüler schreiben Artikel über Sport, Musik und das Leben an der Schule. „Wir wollen über Themen sprechen, die uns interessieren“, sagt die Chefredakteurin Amina, 16 Jahre alt. Die Zeitung kostet 100 Franken. Das Geld finanziert den Druck der nächsten Ausgabe.}\par
\textbf{Questions :} 1. Wie heißt die Schülerzeitung ? 2. Wie oft erscheint sie ? 3. Wer schreibt die Artikel ? 4. Was macht die Zeitung mit dem Geld ?
\tcblower
\textbf{Raisonnement} : on repère les 5 W. Le nom du journal est entre guillemets allemands ; la fréquence est un complément de temps ; le sujet de \all{schreiben} répond à « qui » ; la dernière phrase répond à « que fait-on de l'argent ».\par
\textbf{Réponses rédigées :}\par
1. \all{Die Schülerzeitung heißt „Der Pausenfrosch“.}\par
2. \all{Sie erscheint jeden Monat.} (Attention : \all{jeden Monat} = accusatif de temps, car \all{der Monat} est masculin.)\par
3. \all{Zwanzig Schülerinnen und Schüler schreiben die Artikel.}\par
4. \all{Mit dem Geld finanziert die Zeitung den Druck der nächsten Ausgabe.} (Le verbe en 2\textsuperscript{e} position : le complément \all{Mit dem Geld} occupe la 1\textsuperscript{re} place, donc le sujet vient après le verbe.)\par
\textbf{Conclusion} : chaque réponse est une phrase complète avec verbe conjugué en 2\textsuperscript{e} position et majuscules aux noms.
\end{exoresolu}

\courssub{Lexique thématique : la presse et le journal scolaire}

\begin{definition}[Wortschatz — Medien und Schülerzeitung]
\begin{center}
\begin{tabularx}{\linewidth}{|X|X|}\hline
\rowcolor{popblueL} \textbf{Deutsch} & \textbf{Français} \\ \hline
die Zeitung, -en & le journal, le quotidien \\ \hline
die Zeitschrift, -en & le magazine, la revue \\ \hline
die Schülerzeitung, -en & le journal scolaire (des élèves) \\ \hline
der Artikel, - & l'article \\ \hline
die Überschrift, -en & le titre (d'un article) \\ \hline
der Bericht, -e & le reportage, le compte rendu \\ \hline
das Interview, -s & l'interview \\ \hline
die Anzeige, -n / die Ankündigung, -en & l'annonce \\ \hline
der Journalist, -en / die Journalistin, -nen & le / la journaliste \\ \hline
der Redakteur, -e / die Redakteurin, -nen & le / la rédacteur(-trice) \\ \hline
der Chefredakteur, -e / die Chefredakteurin, -nen & le / la rédacteur(-trice) en chef \\ \hline
die Ausgabe, -n & le numéro, l'édition \\ \hline
die Nachricht, -en & la nouvelle, l'information \\ \hline
die Presse (nur Singular) & la presse \\ \hline
der Druck, -e / drucken & l'impression / imprimer \\ \hline
erscheinen (erscheint, erschien, ist erschienen) & paraître \\ \hline
veröffentlichen (hat veröffentlicht) & publier \\ \hline
berichten (hat berichtet) & rapporter, rendre compte \\ \hline
\end{tabularx}
\end{center}
\textbf{Familles de mots} : \all{die Schule} → \all{der Schüler} → \all{die Schülerzeitung} ; \all{schreiben} → \all{die Überschrift} → \all{der Schriftsteller}.
\end{definition}

\begin{attention}[Faux amis du vocabulaire des médias]
\begin{itemize}
\item \all{die Zeitung} = le quotidien ; \all{die Zeitschrift} = le magazine. Ne traduis jamais « journal » par \all{das Journal} pour un quotidien scolaire.
\item \all{aktuelle Nachrichten} = les nouvelles récentes ; \all{aktuell} ne signifie pas toujours « actuellement » (qui se dit plutôt \all{zurzeit} ou \all{im Moment}).
\item \all{die Notiz} = la note écrite, pas « la notice ».
\item \all{die Presse} = la presse (les médias), jamais « la presse » (machine).
\end{itemize}
\end{attention}

\courssub{Grammaire : révision du Perfekt, du Präteritum et de l'ordre des mots}

\begin{propriete}[Le Perfekt : formation et auxiliaires]
\textbf{Formation} : auxiliaire \all{haben} ou \all{sein} conjugué en 2\textsuperscript{e} position + participe passé à la fin de la phrase.\par
\textbf{Choix de l'auxiliaire} :
\begin{itemize}
\item \all{haben} : par défaut (verbes transitifs, la plupart des verbes) : \all{Ich habe einen Artikel geschrieben.}
\item \all{sein} : verbes de mouvement ou de changement d'état (\all{gehen, kommen, fahren, aufstehen, sterben}) + \all{sein} et \all{bleiben} : \all{Wir sind zum Markt gegangen.}
\end{itemize}
\textbf{Participes passés} :
\begin{center}
\begin{tabular}{|l|l|l|}\hline
\rowcolor{popblueL} \textbf{Type de verbe} & \textbf{Formation} & \textbf{Exemple} \\ \hline
Verbe faible & ge + radical + t & \all{machen → gemacht} \\ \hline
Verbe fort & ge + (radical modifié) + en & \all{schreiben → geschrieben} \\ \hline
Verbe en -ieren & pas de ge- & \all{organisieren → organisiert} \\ \hline
Verbe à préfixe inséparable & pas de ge- & \all{veröffentlichen → veröffentlicht} \\ \hline
Verbe à particule séparable & ge entre particule et radical & \all{teilnehmen → teilgenommen} \\ \hline
\end{tabular}
\end{center}
\end{propriete}

\begin{propriete}[Le Präteritum : les formes indispensables]
Au niveau Première, on utilise surtout le Präteritum de \all{sein}, \all{haben}, \all{es geben} et des verbes modaux :
\begin{center}
\begin{tabular}{|l|l|l|l|}\hline
\rowcolor{popblueL} \textbf{Personne} & \textbf{sein} & \textbf{haben} & \textbf{können} \\ \hline
ich & war & hatte & konnte \\ \hline
du & warst & hattest & konntest \\ \hline
er/sie/es & war & hatte & konnte \\ \hline
wir & waren & hatten & konnten \\ \hline
ihr & wart & hattet & konntet \\ \hline
sie/Sie & waren & hatten & konnten \\ \hline
\end{tabular}
\end{center}
\all{Es gab ein Fest.} « Il y a eu une fête. » — \all{Der Direktor war zufrieden.} « Le directeur était satisfait. »
\end{propriete}

\begin{propriete}[L'ordre des mots : les trois positions du verbe]
\begin{enumerate}
\item \textbf{Phrase déclarative} : verbe conjugué en 2\textsuperscript{e} position. Si un complément ouvre la phrase, le sujet passe après le verbe (inversion) : \all{Gestern hat die Schule ein Fest organisiert.}
\item \textbf{Question} : sans mot interrogatif, verbe en 1\textsuperscript{re} position (\all{Spielst du Fußball ?}) ; avec mot interrogatif, verbe en 2\textsuperscript{e} position (\all{Warum spielst du Fußball ?}).
\item \textbf{Subordonnée} (\all{dass, weil, wenn}) : verbe conjugué à la fin : \all{Wir hoffen, dass es nächstes Jahr wieder einen Sporttag gibt.}
\end{enumerate}
\end{propriete}

\courssub{Méthode 2 : Rédiger un article de journal scolaire}

\begin{methode}[Écrire un article en 6 étapes]
\begin{enumerate}
\item \textbf{Trouver un titre} court et accrocheur (\all{die Überschrift}) : \all{Neue Mensa an unserer Schule !}
\item \textbf{Rédiger le Lead} : une première phrase qui répond à qui ? quoi ? où ? quand ?
\item \textbf{Développer} avec 2 ou 3 phrases de détails : causes, conséquences, citations.
\item \textbf{Choisir le temps} : \all{Perfekt} pour raconter un événement récent dans un journal scolaire (\all{Die Schule hat ein Fest organisiert.}) ; \all{Präteritum} pour \all{sein}, \all{haben} et les verbes modaux (\all{Es war ein großer Erfolg.}).
\item \textbf{Conclure} par une ouverture : projet, souhait, question aux lecteurs.
\item \textbf{Relire} : verbe en 2\textsuperscript{e} position, majuscules aux noms, participes passés corrects (\all{ge...t / ge...en}).
\end{enumerate}
\textbf{Structures à retenir} : \all{Am Montag hat ... stattgefunden.} / \all{Viele Schüler haben teilgenommen.} / \all{Wir hoffen, dass ...}
\end{methode}

\begin{exoresolu}[Rédaction d'un article — type examen]
\textbf{Sujet :} Votre lycée a organisé une journée sportive. Rédigez un article de 6 à 8 phrases pour le journal scolaire allemand.
\tcblower
\textbf{Plan} : titre → lead (quoi, quand, où) → détails (qui, quelles activités) → citation → bilan → ouverture.\par
\textbf{Article rédigé :}\par
\all{Sporttag an unserem Gymnasium}\par
\all{Am Freitag hat unser Gymnasium in Douala einen großen Sporttag organisiert. Alle Klassen haben an diesem Fest teilgenommen. Am Morgen haben die Schüler Fußball und Basketball gespielt. Am Nachmittag gab es einen Leichtathletik-Wettbewerb. „Der Sporttag ist super, wir haben viel Spaß gehabt“, sagt Marie aus der Klasse 1\textsuperscript{re} A. Der Direktor war sehr zufrieden. Am Ende hat die Klasse Tle C den ersten Preis gewonnen. Wir hoffen, dass es nächstes Jahr wieder einen Sporttag gibt.}\par
\textbf{Justifications grammaticales} : \all{hat organisiert / haben teilgenommen / hat gewonnen} = Perfekt des événements récents ; \all{war}, \all{gab es} = Präteritum de \all{sein} et \all{es geben} (usage naturel) ; \all{Am Freitag hat ...} = inversion du sujet après le complément de temps ; \all{dass es ... gibt} = verbe à la fin dans la subordonnée.
\end{exoresolu}

\courssub{Méthode 3 : Rédiger une interview et une annonce}

\begin{methode}[L'interview : questions et réponses]
\begin{enumerate}
\item \textbf{Préparer 4 à 6 questions} avec les mots interrogatifs : \all{Wie} (comment), \all{Warum} (pourquoi), \all{Seit wann} (depuis quand), \all{Was für} (quel genre de).
\item \textbf{Respecter la place du verbe} : dans une question sans mot interrogatif, le verbe est en 1\textsuperscript{re} position : \all{Spielst du gern Fußball ?} ; avec mot interrogatif, en 2\textsuperscript{e} : \all{Warum spielst du gern Fußball ?}
\item \textbf{Présenter l'invité} en une phrase : \all{Heute sprechen wir mit ...}
\item \textbf{Noter les réponses} en phrases complètes, au Perfekt pour le passé.
\end{enumerate}
\textbf{L'annonce} (\all{die Anzeige / die Ankündigung}) : titre en majuscules, information essentielle (quoi, où, quand, prix), phrase d'appel : \all{Kommt alle !} / \all{Wir freuen uns auf euch !}
\end{methode}

\begin{exoresolu}[Interview — type examen]
\textbf{Sujet :} Vous interviewez un élève champion de football de votre lycée pour le journal scolaire. Écrivez 4 questions et leurs réponses.
\tcblower
\textbf{Raisonnement} : alterner questions avec mot interrogatif (verbe en 2\textsuperscript{e} position) et questions fermées (verbe en 1\textsuperscript{re} position) ; réponses au présent ou au Perfekt.\par
\textbf{Interview rédigée :}\par
\all{Heute sprechen wir mit Jean, dem Kapitän unserer Fußballmannschaft.}\par
\all{— Seit wann spielst du Fußball ? — Ich spiele seit fünf Jahren Fußball.}\par
\all{— Trainierst du jeden Tag ? — Nein, ich trainiere nur am Dienstag und am Samstag.}\par
\all{— Was hat dir am letzten Spiel gefallen ? — Das letzte Spiel war schwer, aber wir haben 2:0 gewonnen. Das war fantastisch !}\par
\all{— Was sind deine Pläne für die Zukunft ? — Ich möchte Profifußballer werden.}\par
\textbf{Justifications} : \all{Seit wann spielst du} = mot interrogatif + verbe + sujet ; \all{Trainierst du} = question fermée, verbe en tête ; \all{haben gewonnen} = Perfekt de \all{gewinnen} (verbe fort : \all{gewinnen – gewann – hat gewonnen}) ; \all{Profifußballer werden} = \all{möchte} + infinitif à la fin ; \all{seit fünf Jahren} = datif pluriel après \all{seit}.
\end{exoresolu}

\begin{exemplebox}[Modèle d'annonce pour le journal scolaire]
\all{FUSSBALLTURNIER !}\par
\all{Am Samstag, den 15. Juni, findet auf dem Sportplatz unseres Gymnasiums ein großes Fußballturnier statt. Beginn: 9 Uhr. Eintritt: 200 Franken. Alle Klassen können eine Mannschaft anmelden. Kommt alle und unterstützt eure Klasse ! Wir freuen uns auf euch !}\par
« TOURNOI DE FOOTBALL ! Samedi 15 juin, un grand tournoi de football aura lieu sur le terrain de sport de notre lycée. Début : 9 h. Entrée : 200 francs. Toutes les classes peuvent inscrire une équipe. Venez tous soutenir votre classe ! Nous vous attendons avec joie ! »\par
\textbf{Points de langue} : \all{findet ... statt} = verbe à particule séparable (\all{stattfinden}) ; \all{am Samstag, den 15. Juni} = date à l'accusatif ; \all{Kommt alle !} = impératif pluriel.
\end{exemplebox}

\courssub{Méthode 4 : Traduire — le thème (français → allemand)}

\begin{methode}[Traduire du français vers l'allemand en 5 étapes]
\begin{enumerate}
\item \textbf{Lire toute la phrase} et identifier le temps (passé composé → Perfekt ; imparfait avec être/avoir → Präteritum).
\item \textbf{Trouver le sujet et le verbe}, puis placer le verbe conjugué en 2\textsuperscript{e} position.
\item \textbf{Choisir l'auxiliaire} : \all{haben} par défaut ; \all{sein} pour les verbes de mouvement ou de changement (\all{gehen, kommen, fahren, sterben, aufstehen}).
\item \textbf{Placer le participe passé à la fin} : \all{Ich habe einen Artikel geschrieben.}
\item \textbf{Vérifier} : majuscules aux noms, genre des articles, cas après prépositions (\all{mit dem}, \all{für die}, \all{in der Schule}).
\end{enumerate}
\textbf{Réflexe anti-calque} : le français dit « j'ai visité le marché avec mon ami » → l'allemand dit \all{Ich habe mit meinem Freund den Markt besucht} : participe en fin de phrase, jamais juste après l'auxiliaire.
\end{methode}

\begin{exoresolu}[Thème — type examen]
\textbf{Traduisez en allemand :}\par
1. Hier, les élèves ont publié un nouveau journal scolaire.\par
2. La rédactrice en chef a écrit un article sur le sport.\par
3. Nous sommes allés au marché avec notre professeur.
\tcblower
\textbf{1.} Passé composé → Perfekt ; \all{veröffentlichen} est un verbe faible à préfixe inséparable (participe sans \all{ge-}) → auxiliaire \all{haben}. « Hier » en tête → inversion du sujet.\par
\all{Gestern haben die Schüler eine neue Schülerzeitung veröffentlicht.}\par
\textbf{2.} \all{schreiben} est un verbe fort : \all{schreiben – schrieb – hat geschrieben} ; « sur le sport » = \all{über den Sport} (accusatif après \all{über} pour indiquer le thème).\par
\all{Die Chefredakteurin hat einen Artikel über den Sport geschrieben.}\par
\textbf{3.} \all{gehen} = verbe de mouvement → auxiliaire \all{sein} ; \all{mit} + datif : \all{mit unserem Lehrer} ; « au marché » (direction) = \all{zum Markt} (= \all{zu dem Markt}, datif).\par
\all{Wir sind mit unserem Lehrer zum Markt gegangen.}\par
\textbf{Conclusion} : dans les trois phrases, l'auxiliaire est en 2\textsuperscript{e} position et le participe passé à la fin — la règle d'or du Perfekt.
\end{exoresolu}

\courssub{Méthode 5 : Traduire — la version (allemand → français)}

\begin{methode}[Traduire de l'allemand vers le français en 4 étapes]
\begin{enumerate}
\item \textbf{Repérer le verbe conjugué} (2\textsuperscript{e} position) et son sujet : c'est le squelette de la phrase.
\item \textbf{Aller chercher le participe ou l'infinitif à la fin} avant de traduire : \all{Die Schüler haben gestern eine Zeitung gedruckt} se comprend « Les élèves + ont imprimé + hier + un journal ».
\item \textbf{Traduire les temps correctement} : Perfekt → passé composé ; Präteritum → passé composé ou imparfait selon le sens (\all{es war} = « c'était »).
\item \textbf{Reformuler en français naturel} : ne jamais garder l'ordre allemand mot à mot.
\end{enumerate}
\textbf{Attention aux faux amis} : \all{die Presse} = la presse (les médias) ; \all{die Notiz} = la note écrite ; \all{aktuelle Nachrichten} = les nouvelles récentes.
\end{methode}

\begin{exoresolu}[Version — type examen]
\textbf{Traduisez en français :}\par
\all{Letzte Woche hat unsere Schule einen Medientag organisiert. Ein Journalist aus Yaoundé hat den Schülern erklärt, wie man einen guten Artikel schreibt. Viele Schüler waren begeistert und wollen jetzt für die Schülerzeitung arbeiten.}
\tcblower
\textbf{Analyse phrase par phrase :}\par
1. \all{Letzte Woche hat ... organisiert} : verbe \all{hat organisiert} (Perfekt) → passé composé ; \all{Letzte Woche} = « la semaine dernière ».\par
→ « La semaine dernière, notre école a organisé une journée des médias. »\par
2. \all{hat den Schülern erklärt} : \all{den Schülern} est un datif pluriel (\all{jemandem etwas erklären} = expliquer quelque chose à quelqu'un) ; \all{wie man einen guten Artikel schreibt} = subordonnée, verbe à la fin → « comment on écrit un bon article ».\par
→ « Un journaliste de Yaoundé a expliqué aux élèves comment on écrit un bon article. »\par
3. \all{waren begeistert} : Präteritum de \all{sein} → « étaient enthousiastes » (imparfait, état) ; \all{wollen ... arbeiten} = « veulent travailler » ; \all{für die Schülerzeitung} = « pour le journal scolaire ».\par
→ « Beaucoup d'élèves étaient enthousiastes et veulent maintenant travailler pour le journal scolaire. »\par
\textbf{Conclusion} : on a repéré les verbes, respecté le datif (\all{den Schülern} = « aux élèves ») et reformulé en français naturel.
\end{exoresolu}

\courssub{Méthode 6 : Relire et corriger sa production}

\begin{methode}[La grille de relecture en 5 points]
\begin{enumerate}
\item \textbf{Verbe en 2\textsuperscript{e} position} dans chaque phrase principale ; verbe à la fin après \all{dass, weil, wenn}.
\item \textbf{Majuscule à tous les noms} : \all{die Schule, der Artikel, das Fest}.
\item \textbf{Participes passés} : verbes faibles \all{ge...t} (\all{gemacht}), verbes forts \all{ge...en} (\all{geschrieben}), verbes en \all{-ieren} sans \all{ge-} (\all{organisiert}).
\item \textbf{Cas} : accusatif après \all{für, durch, ohne, gegen, über} ; datif après \all{mit, von, zu, bei, nach, seit}.
\item \textbf{Orthographe} : \all{ß} ou \all{ss}, Umlaute \all{ä, ö, ü}, guillemets allemands „…“ pour les citations.
\end{enumerate}
\end{methode}

\begin{exoresolu}[Correction d'un texte fautif — type examen]
\textbf{Corrigez les 6 fautes de ce texte d'élève :}\par
\all{Gestern ich habe ein artikel für die schülerzeitung schreiben. Der artikel war über sport. Meine lehrerin hat gesagt, dass der Artikel sehr gut ist.}
\tcblower
\textbf{Corrections détaillées :}\par
1. \all{Gestern ich habe} → \all{Gestern habe ich} : le verbe conjugué doit être en 2\textsuperscript{e} position, le sujet après le verbe (inversion).\par
2. \all{ein artikel} → \all{einen Artikel} : majuscule au nom et accusatif masculin (\all{der Artikel} → \all{einen Artikel}).\par
3. \all{schülerzeitung} → \all{Schülerzeitung} : majuscule obligatoire à tous les noms.\par
4. \all{schreiben} → \all{geschrieben} : au Perfekt, le participe passé du verbe fort \all{schreiben} se place à la fin.\par
5. \all{Der artikel war über sport} → \all{Der Artikel war über den Sport} : majuscules + accusatif après \all{über}.\par
6. \all{Meine lehrerin} → \all{Meine Lehrerin} : majuscule au nom.\par
\textbf{Texte corrigé :} \all{Gestern habe ich einen Artikel für die Schülerzeitung geschrieben. Der Artikel war über den Sport. Meine Lehrerin hat gesagt, dass der Artikel sehr gut ist.}
\end{exoresolu}

\begin{attention}[Les 5 erreurs qui coûtent des points à l'examen]
\begin{enumerate}
\item \textbf{Oublier la majuscule aux noms} : écrire \all{die schule} au lieu de \all{die Schule}. En allemand, TOUS les noms prennent une majuscule, même au milieu de la phrase.
\item \textbf{Calquer l'ordre français} : \all{Ich habe geschrieben einen Artikel} est FAUX. Le participe passé va à la fin : \all{Ich habe einen Artikel geschrieben.}
\item \textbf{Mauvaise place du verbe} : après un complément en tête de phrase, le sujet passe après le verbe : \all{Gestern hat die Schule ein Fest organisiert} (et non \all{Gestern die Schule hat...}).
\item \textbf{Confondre les auxiliaires} : les verbes de mouvement prennent \all{sein} : \all{Wir sind zum Markt gegangen} (et non \all{haben gegangen}).
\item \textbf{Faux amis et calques de vocabulaire} : « actuellement » se dit \all{zurzeit} ou \all{im Moment} (\all{aktuelle Nachrichten} = nouvelles récentes) ; « un journal » se dit \all{eine Zeitung} (quotidien) ou \all{eine Zeitschrift} (magazine) ; ne traduis jamais « journal » par \all{Journal} pour un quotidien scolaire.
\end{enumerate}
\end{attention}

\begin{aretenir}[L'essentiel de la leçon]
\begin{itemize}
\item Un article de presse se lit avec les 5 W (\all{Wer, Was, Wo, Wann, Warum}) et se rédige avec un titre, un lead, des détails, une citation et une ouverture.
\item Événement récent raconté → \all{Perfekt} ; \all{sein}, \all{haben}, \all{es geben} → \all{Präteritum}.
\item Règle d'or : verbe conjugué en 2\textsuperscript{e} position, participe passé ou infinitif à la fin, verbe à la fin dans les subordonnées.
\item En traduction, on repère d'abord le verbe et son sujet, puis on reformule dans la langue cible sans calquer l'ordre des mots.
\item La relecture systématique (majuscules, cas, participes, place du verbe) fait gagner des points.
\end{itemize}
\end{aretenir}

\newpage

\coursec{Exercices}

\courssub{Je vérifie mes connaissances}

\begin{exercice}[QCM – Presse et journal scolaire]
Pour chaque question, entoure la bonne réponse.
\begin{enumerate}
\item Le titre d'un article de journal s'appelle :
\begin{enumerate}
\item der Artikel
\item die Überschrift
\item die Redaktion
\end{enumerate}
\item Le premier paragraphe d'un article, qui résume l'essentiel, s'appelle :
\begin{enumerate}
\item der Schluss
\item das Interview
\item der Lead (die Einleitung)
\end{enumerate}
\item Quel temps utilise-t-on de préférence dans un article de presse écrit ?
\begin{enumerate}
\item das Präteritum
\item das Futur II
\item der Imperativ
\end{enumerate}
\item \all{die Anzeige} signifie :
\begin{enumerate}
\item l'interview
\item l'annonce
\item la une
\end{enumerate}
\end{enumerate}
\end{exercice}

\begin{exercice}[Vrai ou faux – Justifie ta réponse]
Indique si les affirmations suivantes sont vraies (V) ou fausses (F) et justifie en une phrase.
\begin{enumerate}
\item Dans une phrase déclarative allemande, le verbe conjugué occupe toujours la deuxième place.
\item \all{Gestern habe ich den Direktor gesehen.} est une phrase au Präteritum.
\item Dans un journal scolaire, on peut publier un article, une interview et une annonce.
\item Le verbe \all{interviewen} est un verbe fort.
\item En allemand, tous les noms communs prennent une majuscule, même au milieu de la phrase.
\end{enumerate}
\end{exercice}

\begin{exercice}[Vocabulaire – Le monde de la presse]
Complète chaque phrase avec le mot qui convient, choisi dans la liste : \all{die Redaktion} – \all{der Journalist} – \all{die Schülerzeitung} – \all{die Überschrift} – \all{veröffentlichen}.
\begin{enumerate}
\item Unsere \ldots\ erscheint jeden Monat am Gymnasium.
\item Der \ldots\ schreibt einen Artikel über das Fußballspiel.
\item Die \ldots\ des Artikels ist zu lang : nur fünf Wörter, bitte !
\item Die \ldots\ trifft sich jeden Dienstag nach dem Unterricht.
\item Wir \ldots\ das Interview in der nächsten Ausgabe.
\end{enumerate}
\end{exercice}

\begin{exercice}[Conjugaison – Perfekt ou Präteritum ?]
Conjugue le verbe entre parenthèses au temps demandé.
\begin{enumerate}
\item Gestern \ldots\ (besuchen, Präteritum, er) der Trainer unsere Schule.
\item Wir \ldots\ (schreiben, Perfekt) einen Artikel über den Umwelttag.
\item Die Schüler \ldots\ (applaudieren, Präteritum) nach der Rede.
\item \ldots\ du \ldots\ (lesen, Perfekt) die neue Schülerzeitung ?
\item Der Direktor \ldots\ (eröffnen, Präteritum) die neue Bibliothek.
\end{enumerate}
\end{exercice}

\courssub{Je m'entraîne}

\begin{exercice}[Thème – Phrases à traduire (FR $\rightarrow$ DE)]
Traduis les phrases suivantes en allemand. Attention à l'ordre des mots et au choix du temps.
\begin{enumerate}
\item Hier, le journaliste a interviewé le directeur.
\item La rédaction publie un article chaque semaine.
\item Les élèves ont lu le journal scolaire avec intérêt.
\item Demain, nous écrirons une annonce pour le club de théâtre.
\item Le match a commencé à 15 heures et les spectateurs ont applaudi.
\end{enumerate}
\end{exercice}

\begin{exercice}[Transformation – Du Perfekt au Präteritum]
Transforme les phrases suivantes du Perfekt au Präteritum, puis indique dans quel contexte (oral ou écrit) chaque version est la plus naturelle.
\begin{enumerate}
\item Der Reporter hat den Bürgermeister interviewt.
\item Wir haben die Bibliothek besucht.
\item Die Schüler haben laut applaudiert.
\item Ich habe den Artikel gelesen.
\item Das Spiel hat um 16 Uhr begonnen.
\item Sie hat eine Frage gestellt.
\end{enumerate}
\end{exercice}

\begin{exercice}[Texte à trous – Un article de presse]
Complète l'article suivant avec les mots de la liste : \all{fand statt} – \all{die Schüler} – \all{der Direktor} – \all{ein Interview} – \all{die Redaktion} – \all{gestern}.
\begin{textebox}[Un article à compléter]
\all{Großer Erfolg am Gymnasium Yaoundé}

\all{\ldots\ (1) \ldots\ der Umwelttag an unserer Schule. \ldots\ (2) hat am Morgen die Veranstaltung eröffnet. \ldots\ (3) haben Bäume gepflanzt und den Schulhof geputzt. \ldots\ (4) der Schülerzeitung war auch da : sie hat mit dem Direktor \ldots\ (5) geführt. Am Nachmittag gab es ein Fußballspiel gegen das Gymnasium Biyem-Assi.}
\end{textebox}
\end{exercice}

\begin{exercice}[Appariement – Les genres journalistiques]
Relie chaque extrait au genre journalistique correspondant : \textbf{a)} l'article \quad \textbf{b)} l'interview \quad \textbf{c)} l'annonce.
\begin{enumerate}
\item \all{Frage : Wie oft trainieren Sie? – Antwort : Jeden Tag zwei Stunden.}
\item \all{Der Theaterclub sucht neue Mitglieder ! Treffen : Freitag, 14 Uhr, Raum 12.}
\item \all{Am Samstag fand das Finale des Schulfußballturniers statt. Über 200 Zuschauer waren gekommen.}
\item \all{Verkaufe : gebrauchtes Fahrrad, guter Zustand, Preis : 25 000 FCFA. Kontakt : Klasse Tle A.}
\item \all{Die neue Mensa wurde am Montag eingeweiht. Der Direktor dankte allen Helfern.}
\end{enumerate}
\end{exercice}

\courssub{Je me prépare au Bac}

\begin{exercice}[Compréhension de l'écrit et questions de langue]
Lis le texte suivant, puis réponds aux questions.

\begin{textebox}[Die Schülerzeitung „Der Palaver“]
\all{Seit einem Jahr gibt es am Gymnasium Nkol-Eton in Yaoundé eine Schülerzeitung. Sie heißt „Der Palaver“ und erscheint jeden Monat. Die Redaktion besteht aus zwölf Schülern der Klassen Seconde und Première. Jeden Dienstag treffen sie sich nach dem Unterricht und diskutieren über die Themen der nächsten Ausgabe : Sport, Umwelt, Musik und das Leben an der Schule.}

\all{Letzten Monat interviewte die Redaktion den neuen Sporttrainer. Er sagte : „Unsere Mannschaft trainiert hart, und wir wollen das Turnier gewinnen.“ Der Artikel war ein großer Erfolg : Viele Schüler lasen ihn und sprachen darüber.}

\all{„Die Zeitung hilft uns, besser Deutsch zu schreiben“, erklärt Amina, die Chefredakteurin. „Und wir lernen, im Team zu arbeiten.“ Die nächste Ausgabe erscheint im Mai – mit einem Bericht über den Umwelttag und einer Anzeige des Theaterclubs.}

\emph{(environ 150 mots)}
\end{textebox}

\textbf{A. Compréhension} – Réponds en français par des phrases complètes.
\begin{enumerate}
\item Depuis combien de temps le journal scolaire existe-t-il ?
\item Qui compose la rédaction et quand se réunit-elle ?
\item Quels sujets le journal aborde-t-il ?
\item Pourquoi l'article sur le sport a-t-il eu du succès ?
\item Selon Amina, quels sont les deux avantages du journal pour les élèves ?
\end{enumerate}

\textbf{B. Questions de langue}
\begin{enumerate}
\item Releve dans le texte trois verbes au Präteritum et donne leur infinitif.
\item Mets à la forme négative : \all{Die Zeitung erscheint jeden Monat.}
\item Mets à la question : \all{Die Redaktion interviewte den Trainer.}
\item Souligne dans le texte deux verbes à particule séparable et explique leur place dans la phrase.
\end{enumerate}
\end{exercice}

\begin{exercice}[Thème et version – Traduction]
\textbf{A. Thème} – Traduis en allemand le paragraphe suivant. (Attendu : Präteritum, ordre des mots V2, déclinaisons correctes.)

« Hier, le principal a inauguré la nouvelle bibliothèque. Les élèves ont applaudi. Ensuite, un journaliste a interviewé le principal. Il a dit : "Cette bibliothèque est un cadeau pour toute l'école." »

\textbf{B. Version} – Traduis en français le texte suivant.

\begin{textebox}[À traduire]
\all{Die Redaktion interviewte den Trainer. Er sagte : „Wir trainieren hart.“ Dann besuchten die Schüler das Stadion. Das Spiel begann um 15 Uhr, und am Ende applaudierten alle Zuschauer.}
\end{textebox}
\end{exercice}

\begin{exercice}[Production écrite – Article de journal scolaire]
Tu es membre de la rédaction du journal de ton lycée. Rédige en allemand un article d'environ \textbf{130 mots} sur le sujet suivant : « La journée de l'environnement au lycée ».

Ton article doit contenir :
\begin{itemize}
\item un titre (die Überschrift) ;
\item un lead (une ou deux phrases qui résument l'essentiel : qui ? quoi ? où ? quand ?) ;
\item trois paragraphes (le déroulement de la journée, les activités des élèves, une citation d'un participant) ;
\item une conclusion courte.
\end{itemize}

\textbf{Contraintes de langue :} utilise le Präteritum pour le récit, au moins deux verbes à particule séparable, une citation au discours direct avec les guillemets allemands „…“, et respecte l'ordre des mots (verbe en deuxième position). N'oublie pas la majuscule à tous les noms.
\end{exercice}

\newpage

\coursec{Corrigés des exercices}

\coursec{Corrigés des exercices — Leçon ALA20 : Traduction et production d'articles de journaux scolaires et de textes de presse}

\courssub{Exercice 1 — QCM : Presse et journal scolaire}
\begin{corrige}[Exercice 1 — QCM : Presse et journal scolaire]
\begin{enumerate}
\item \textbf{b)} \all{die Überschrift}. Le titre d'un article se dit \all{die Überschrift} ; \all{der Artikel} désigne l'article lui-même et \all{die Redaktion} l'équipe de rédaction.
\item \textbf{c)} \all{der Lead (die Einleitung)}. Le premier paragraphe qui résume l'essentiel (qui ? quoi ? où ? quand ?) est le lead ; \all{der Schluss} est la conclusion.
\item \textbf{a)} \all{das Präteritum}. C'est le temps du récit écrit (presse, romans, rapports) ; le Perfekt domine à l'oral.
\item \textbf{b)} l'annonce. \all{die Anzeige} = l'annonce (publicité, petite annonce) ; l'interview se dit \all{das Interview}.
\end{enumerate}
\end{corrige}

\courssub{Exercice 2 — Vrai ou faux}
\begin{corrige}[Exercice 2 — Vrai ou faux]
\begin{enumerate}
\item \textbf{Vrai.} Dans une phrase déclarative allemande, le verbe conjugué occupe toujours la deuxième position (règle V2), même si un complément ouvre la phrase : \all{Gestern \textbf{habe} ich den Direktor gesehen.}
\item \textbf{Faux.} Cette phrase est au \cle{Perfekt} (auxiliaire \all{habe} + participe passé \all{gesehen}) ; au Präteritum, on écrirait \all{Gestern sah ich den Direktor.}
\item \textbf{Vrai.} Un journal scolaire (\all{die Schülerzeitung}) publie des articles (\all{der Artikel}), des interviews (\all{das Interview}) et des annonces (\all{die Anzeige}).
\item \textbf{Faux.} \all{interviewen} est un verbe \cle{faible} (régulier) : \all{interviewte}, \all{hat interviewt}. Les verbes en \all{-ieren} et les verbes récents empruntés sont toujours faibles.
\item \textbf{Vrai.} En allemand, tous les noms communs prennent une majuscule, où qu'ils soient dans la phrase : \all{Die Schüler lesen die Zeitung.}
\end{enumerate}
\end{corrige}

\courssub{Exercice 3 — Vocabulaire : Le monde de la presse}
\begin{corrige}[Exercice 3 — Vocabulaire : Le monde de la presse]
\begin{enumerate}
\item Unsere \textbf{Schülerzeitung} erscheint jeden Monat am Gymnasium. « Notre journal scolaire paraît chaque mois au lycée. »
\item Der \textbf{Journalist} schreibt einen Artikel über das Fußballspiel. « Le journaliste écrit un article sur le match de football. »
\item Die \textbf{Überschrift} des Artikels ist zu lang: nur fünf Wörter, bitte ! « Le titre de l'article est trop long : seulement cinq mots, s'il vous plaît ! »
\item Die \textbf{Redaktion} trifft sich jeden Dienstag nach dem Unterricht. « La rédaction se réunit chaque mardi après les cours. » (Attention : \all{sich treffen} est réfléchi ; ici \all{die Redaktion} est singulier, donc \all{trifft sich}.)
\item Wir \textbf{veröffentlichen} das Interview in der nächsten Ausgabe. « Nous publions l'interview dans le prochain numéro. »
\end{enumerate}
\end{corrige}

\courssub{Exercice 4 — Conjugaison : Perfekt ou Präteritum ?}
\begin{corrige}[Exercice 4 — Conjugaison : Perfekt ou Präteritum ?]
\begin{enumerate}
\item Gestern \textbf{besuchte} der Trainer unsere Schule. (Präteritum faible : \all{besuch-te} ; le verbe reste en position 2 après \all{gestern}.)
\item Wir \textbf{haben geschrieben}. (Perfekt du verbe fort \all{schreiben} : \all{schreiben – schrieb – hat geschrieben}.)
\item Die Schüler \textbf{applaudierten}. (Präteritum faible en \all{-ierten} ; les verbes en \all{-ieren} sont toujours faibles.)
\item \textbf{Hast} du die neue Schülerzeitung \textbf{gelesen} ? (Perfekt de \all{lesen} : \all{las – hat gelesen} ; à la question, l'auxiliaire est en position 1 et le participe à la fin.)
\item Der Direktor \textbf{eröffnete} die neue Bibliothek. (Präteritum faible ; \all{eröffnen} est un verbe à particule \emph{inséparable} en \all{er-}, donc jamais séparé.)
\end{enumerate}
\end{corrige}

\courssub{Exercice 5 — Thème : phrases à traduire (FR $\rightarrow$ DE)}
\begin{corrige}[Exercice 5 — Thème : phrases à traduire (FR $\rightarrow$ DE)]
\begin{enumerate}
\item \all{Gestern hat der Journalist den Direktor interviewt.} — Perfekt (récit oral) ; \all{gestern} en tête $\Rightarrow$ inversion sujet-verbe ; \all{den Direktor} à l'accusatif masculin.
\item \all{Die Redaktion veröffentlicht jede Woche einen Artikel.} — Présent (habitude) ; \all{jede Woche} : accusatif de temps, article féminin.
\item \all{Die Schüler haben die Schülerzeitung mit Interesse gelesen.} — Perfekt de \all{lesen} (\all{hat gelesen}) ; participe en fin de phrase.
\item \all{Morgen werden wir eine Anzeige für den Theaterclub schreiben.} — Futur I : \all{werden} conjugué en position 2 + infinitif à la fin ; \all{für} + accusatif (\all{den Theaterclub}).
\item \all{Das Spiel begann um 15 Uhr, und die Zuschauer applaudierten.} — Präteritum (récit écrit) : \all{beginnen – begann}, verbe fort ; \all{applaudieren} est faible : \all{applaudierte(n)} ; après \all{und}, chaque proposition garde l'ordre V2.
\end{enumerate}
\textbf{Points de vigilance :} majuscule aux noms (\all{Journalist, Direktor, Spiel, Zuschauer}) ; verbe conjugué toujours en deuxième position ; participe passé et infinitif renvoyés en fin de phrase.
\end{corrige}

\courssub{Exercice 6 — Transformation : du Perfekt au Präteritum}
\begin{corrige}[Exercice 6 — Transformation : du Perfekt au Präteritum]
\begin{enumerate}
\item \all{Der Reporter \textbf{interviewte} den Bürgermeister.} — verbe faible ; Präteritum : contexte écrit (article) ; Perfekt : oral.
\item \all{Wir \textbf{besuchten} die Bibliothek.} — faible ; écrit : Präteritum, oral : Perfekt.
\item \all{Die Schüler \textbf{applaudierten} laut.} — faible (verbe en \all{-ieren}) ; écrit : Präteritum.
\item \all{Ich \textbf{las} den Artikel.} — verbe fort \all{lesen – las – gelesen} ; écrit : Präteritum, oral : Perfekt.
\item \all{Das Spiel \textbf{begann} um 16 Uhr.} — fort \all{beginnen – begann – begonnen} ; écrit : Präteritum.
\item \all{Sie \textbf{stellte} eine Frage.} — faible \all{stellen – stellte – gestellt} ; écrit : Präteritum.
\end{enumerate}
\textbf{Règle :} le Präteritum est le temps du récit \emph{écrit} (presse, roman, rapport) ; le Perfekt domine dans la conversation \emph{orale}. Exception : \all{sein}, \all{haben} et les modaux s'emploient au Präteritum même à l'oral.
\end{corrige}

\courssub{Exercice 7 — Texte à trous : un article de presse}
\begin{corrige}[Exercice 7 — Texte à trous : un article de presse]
\begin{enumerate}
\item \textbf{Gestern fand} der Umwelttag an unserer Schule \textbf{statt}. — \all{stattfinden} est un verbe à particule séparable : \all{fand} en position 2, \all{statt} à la fin. « Hier a eu lieu la journée de l'environnement dans notre lycée. »
\item \textbf{Der Direktor} hat am Morgen die Veranstaltung eröffnet. — sujet au nominatif. « Le directeur a ouvert la manifestation le matin. »
\item \textbf{Die Schüler} haben Bäume gepflanzt und den Schulhof geputzt. — sujet pluriel. « Les élèves ont planté des arbres et nettoyé la cour. »
\item \textbf{Die Redaktion} der Schülerzeitung war auch da. — génitif : \all{der Schülerzeitung}. « La rédaction du journal scolaire était là aussi. »
\item sie hat mit dem Direktor \textbf{ein Interview} geführt. — \all{ein Interview führen} = mener une interview ; accusatif neutre. « elle a mené une interview avec le directeur. »
\end{enumerate}
\end{corrige}

\courssub{Exercice 8 — Appariement : les genres journalistiques}
\begin{corrige}[Exercice 8 — Appariement : les genres journalistiques]
\begin{enumerate}
\item \textbf{b) l'interview} : la structure question-réponse (\all{Frage … Antwort …}) est la marque de l'interview.
\item \textbf{c) l'annonce} : appel avec informations pratiques (jour, heure, lieu) ; \all{sucht neue Mitglieder} = cherche de nouveaux membres.
\item \textbf{a) l'article} : récit factuel d'un événement au Präteritum/Plusquamperfekt (\all{fand … statt}, \all{waren gekommen}).
\item \textbf{c) l'annonce} : petite annonce de vente (\all{Verkaufe … Preis … Kontakt}).
\item \textbf{a) l'article} : compte rendu d'événement (\all{wurde eingeweiht} = a été inaugurée, passif au Präteritum).
\end{enumerate}
\end{corrige}

\courssub{Exercice 9 — Compréhension de l'écrit et questions de langue}
\begin{corrige}[Exercice 9 — Compréhension de l'écrit et questions de langue]
\textbf{A. Compréhension}
\begin{enumerate}
\item Le journal scolaire existe depuis un an (\all{seit einem Jahr}).
\item La rédaction est composée de douze élèves des classes de Seconde et de Première ; elle se réunit chaque mardi après les cours.
\item Le journal aborde le sport, l'environnement, la musique et la vie au lycée.
\item L'article sur le sport a eu du succès parce que beaucoup d'élèves l'ont lu et en ont parlé.
\item Selon Amina, le journal aide les élèves à mieux écrire l'allemand et leur apprend à travailler en équipe.
\end{enumerate}
\textbf{B. Questions de langue}
\begin{enumerate}
\item Trois verbes au Präteritum : \all{interviewte} (\all{interviewen}), \all{sagte} (\all{sagen}), \all{lasen} (\all{lesen}). On accepte aussi \all{sprachen} (\all{sprechen}), \all{war} (\all{sein}).
\item \all{Die Zeitung erscheint \textbf{nicht} jeden Monat.} — négation de la phrase entière : \all{nicht} se place devant le complément nié (ici devant \all{jeden Monat}).
\item \all{\textbf{Interviewte} die Redaktion den Trainer ?} — question sans mot interrogatif : verbe conjugué en position 1.
\item Deux constructions à élément séparable en fin de proposition : \all{treffen … sich} (\all{sich treffen}, pronom réfléchi) et \all{sprachen … darüber} (\all{darüber sprechen}, adverbe prépositionnel). Règle : au Präsens/Präteritum, le verbe conjugué est en position 2 et l'élément séparable (particule, pronom réfléchi, adverbe prépositionnel) est renvoyé à la fin de la proposition. Note : \all{erscheint} (\all{erscheinen}) n'est \emph{pas} un verbe à particule séparable (préfixe \all{er-} inséparable).
\end{enumerate}
\textbf{Barème indicatif (sur 20) :} A. Compréhension : 5 $\times$ 2 = 10 pts ; B. Langue : 2,5 pts par question = 10 pts. Toute réponse de compréhension doit être une phrase complète en français.
\end{corrige}

\courssub{Exercice 10 — Thème et version : traduction}
\begin{corrige}[Exercice 10 — Thème et version : traduction]
\textbf{A. Thème — proposition de traduction :}

\all{Gestern eröffnete der Direktor die neue Bibliothek. Die Schüler applaudierten. Dann interviewte ein Journalist den Direktor. Er sagte: „Diese Bibliothek ist ein Geschenk für die ganze Schule.“}

\textbf{Points de vigilance :}
\begin{itemize}
\item \all{eröffnen} (inaugurer) : particule \all{er-} inséparable, Präteritum \all{eröffnete}.
\item « Hier » en tête $\Rightarrow$ inversion : \all{Gestern eröffnete der Direktor …} (verbe en position 2).
\item « le principal » : sujet au nominatif \all{der Direktor} ; « a interviewé le principal » : accusatif \all{den Direktor}.
\item « ensuite » : \all{dann} en tête $\Rightarrow$ inversion \all{Dann interviewte ein Journalist …}.
\item Citation au discours direct avec les guillemets allemands „…“ ; \all{für} + accusatif : \all{für die ganze Schule}.
\end{itemize}

\textbf{B. Version — proposition de traduction :}

« La rédaction a interviewé l'entraîneur. Il a dit : "Nous nous entraînons dur." Ensuite, les élèves ont visité le stade. Le match a commencé à 15 heures, et à la fin, tous les spectateurs ont applaudi. »

\textbf{Points de vigilance :} \all{der Trainer} = l'entraîneur (et non « le formateur ») ; \all{hart trainieren} = s'entraîner dur ; \all{am Ende} = à la fin ; \all{die Zuschauer} = les spectateurs.

\textbf{Barème indicatif (sur 20) :} thème 10 pts (2 pts par phrase : temps 0,5, ordre des mots 0,5, cas et vocabulaire 1) ; version 10 pts (fidélité 6, qualité du français 4).
\end{corrige}

\courssub{Exercice 11 — Production écrite : article de journal scolaire}
\begin{corrige}[Exercice 11 — Production écrite : article de journal scolaire]
\textbf{Proposition de rédaction modèle (environ 135 mots) :}

\begin{textebox}[Modèle d'article]
\all{\textbf{Der Umwelttag am Gymnasium: Ein großer Erfolg}}

\all{Am Freitag fand am Gymnasium in Yaoundé der erste Umwelttag statt. Über 300 Schüler und Lehrer nahmen an dieser Veranstaltung teil.}

\all{Der Direktor eröffnete den Tag um 8 Uhr mit einer kurzen Rede. Dann begannen die Aktivitäten: Die Klassen pflanzten zwanzig Bäume im Schulhof und sammelten den Müll rund um die Schule ein.}

\all{Die Schüler arbeiteten den ganzen Vormittag mit großem Eifer. Einige malten Plakate, andere bauten Blumenkästen aus altem Holz. Am Nachmittag gab es einen Vortrag über das Recycling.}

\all{„Ich habe heute viel gelernt“, sagte Amina aus der Klasse Première A. „Unsere Schule ist jetzt sauberer und schöner.“}

\all{Der Umwelttag war ein voller Erfolg. Die Redaktion hofft: Nächstes Jahr gibt es ihn wieder!}
\end{textebox}

\textbf{Traduction du modèle :} « Vendredi a eu lieu au lycée de Yaoundé la première journée de l'environnement. Plus de 300 élèves et professeurs ont participé à cette manifestation. Le directeur a ouvert la journée à 8 heures par un court discours. Puis les activités ont commencé : les classes ont planté vingt arbres dans la cour et ont ramassé les déchets autour de l'école. Les élèves ont travaillé toute la matinée avec beaucoup d'ardeur. Certains ont peint des affiches, d'autres ont construit des jardinières avec du vieux bois. L'après-midi, il y a eu une conférence sur le recyclage. "J'ai beaucoup appris aujourd'hui", dit Amina de la classe de Première A. "Notre école est maintenant plus propre et plus belle." La journée de l'environnement a été un plein succès. La rédaction espère : l'année prochaine, elle aura lieu de nouveau ! »

\textbf{Contraintes vérifiées dans le modèle :}
\begin{itemize}
\item Präteritum du récit : \all{fand … statt, eröffnete, begannen, pflanzten, arbeiteten, sagte, war}.
\item Deux verbes à particule séparable (au minimum) : \all{stattfinden} (\all{fand … statt}), \all{einsammeln} (\all{sammelten … ein}), \all{teilnehmen} (\all{nahmen … teil}).
\item Citation au discours direct avec guillemets allemands : \all{„Ich habe heute viel gelernt“, sagte Amina …}.
\item Ordre des mots V2 respecté, y compris après \all{am Freitag}, \all{dann}, \all{am Nachmittag}.
\end{itemize}

\textbf{Barème indicatif (sur 20) :} respect du sujet et du genre (titre, lead, 3 paragraphes, conclusion) : 5 pts ; correction grammaticale (Präteritum, ordre des mots, cas) : 7 pts ; lexique de la presse et de l'environnement : 4 pts ; contraintes imposées (séparables, citation, guillemets) : 2 pts ; présentation et orthographe (majuscules, Umlaute) : 2 pts.

\textbf{Points de vigilance pour la relecture :} vérifier la majuscule à chaque nom ; le verbe conjugué en position 2 dans chaque proposition principale ; la particule séparable en fin de proposition ; le Präteritum des verbes forts (\all{finden – fand}, \all{beginnen – begann}, \all{nehmen – nahm}) ; les guillemets allemands „…“ et non les guillemets français.
\end{corrige}

\newpage

\coursec{Fiche bilan}

\begin{popfigure}
\begin{center}
\begin{tikzpicture}[
  mindmap,
  grow cyclic,
  every node/.style={concept, align=center, font=\small},
  root concept/.append style={concept color=popblue, font=\bfseries\small, text=white},
  level 1 concept/.append style={level distance=3.4cm, sibling angle=51.4, font=\footnotesize},
  level 2 concept/.append style={level distance=2.1cm, font=\scriptsize},
]
\node[root concept]{Journal scolaire \& presse}
  child[concept color=poporange]{ node {Lexique}
    child { node {die Zeitung} }
    child { node {der Artikel} }
    child { node {die Schlagzeile} }
  }
  child[concept color=popgreen]{ node {Perfekt}
    child { node {haben/sein + Partizip II} }
  }
  child[concept color=poppurple]{ node {Präteritum}
    child { node {war, hatte, kam} }
  }
  child[concept color=poppink]{ node {Ordre des mots}
    child { node {verbe en 2\textsuperscript{e} position} }
    child { node {TeKaMoLo} }
  }
  child[concept color=popgold]{ node {Types de textes}
    child { node {Artikel} }
    child { node {Interview} }
    child { node {Anzeige} }
  }
  child[concept color=popblue!60]{ node {Traduction}
    child { node {thème FR $\rightarrow$ DE} }
    child { node {version DE $\rightarrow$ FR} }
  }
  child[concept color=popdark]{ node {Relecture}
    child { node {majuscules, cas, verbe} }
  };
\end{tikzpicture}
\end{center}
\legende{Carte mentale de la leçon : le journal scolaire et les textes de presse}
\end{popfigure}

\begin{aretenir}[L'essentiel de la leçon]
\textbf{1. Lexique clé de la presse (à connaître par cœur)}
\begin{itemize}
  \item \all{die Zeitung, -en} : le journal ; \all{die Schülerzeitung, -en} : le journal scolaire
  \item \all{der Artikel, -} : l'article ; \all{die Schlagzeile, -n} : le titre (gros titre)
  \item \all{der Journalist, -en / die Journalistin, -nen} : le/la journaliste
  \item \all{das Interview, -s} : l'interview ; \all{die Anzeige, -n} : l'annonce
  \item \all{die Redaktion, -en} : la rédaction ; \all{der Leser, - / die Leserin, -nen} : le/la lecteur(-trice)
  \item \all{berichten (über + Akkusativ)} : rendre compte de ; \all{veröffentlichen} : publier
  \item \all{erscheinen (ist erschienen)} : paraître ; \all{sich ereignen (hat sich ereignet)} : se produire (événement)
\end{itemize}

\textbf{2. Grammaire à maîtriser}
\begin{center}
\begin{tabularx}{\linewidth}{|l|X|X|}
\hline
\rowcolor{popblueL} \textbf{Point} & \textbf{Règle} & \textbf{Exemple} \\
\hline
Perfekt & auxiliaire \all{haben} ou \all{sein} (mouvement/changement d'état) + Partizip II en fin de phrase & \all{Wir haben einen Artikel geschrieben.} / \all{Er ist nach Yaoundé gefahren.} \\
\hline
Präteritum & surtout \all{sein} (\all{war}), \all{haben} (\all{hatte}), les modaux et les verbes forts dans les récits de presse & \all{Das Fest war ein großer Erfolg.} \\
\hline
Ordre des mots & verbe conjugué en 2\textsuperscript{e} position ; le sujet se place après le verbe si un autre élément ouvre la phrase & \all{Gestern \textbf{hat} die Redaktion die Zeitung veröffentlicht.} \\
\hline
TeKaMoLo & ordre des compléments : temps -- cause -- manière -- lieu & \all{Sie arbeitet heute (Te) fleißig (Mo) in der Schule (Lo).} \\
\hline
Subordonnées & verbe conjugué à la fin (\all{weil, dass, als}) & \all{..., weil das Fest gut \textbf{war}.} \\
\hline
\end{tabularx}
\end{center}

\textbf{3. Les trois types de textes}
\begin{itemize}
  \item \cle{Article} : titre accrocheur (\all{die Schlagzeile}), répond aux questions \all{Wer? Was? Wann? Wo? Warum?}, rédigé au passé (Perfekt/Präteritum), ton objectif.
  \item \cle{Interview} : alternance question/réponse (\all{Frage -- Antwort}), verbes introducteurs \all{fragen, antworten, sagen, meinen}, inversion du sujet après le verbe introducteur : \all{„Haben Sie viele Leser?“, fragt der Journalist.}
  \item \cle{Annonce} : courte et informative : quoi, quand, où, contact ; impératif possible : \all{Kommt alle!}
\end{itemize}

\textbf{4. Méthodes express}
\begin{itemize}
  \item \textbf{Thème (FR $\rightarrow$ DE)} : repérer le temps $\rightarrow$ placer le verbe conjugué en 2\textsuperscript{e} position $\rightarrow$ mettre le Partizip II en fin de phrase $\rightarrow$ vérifier majuscules et cas.
  \item \textbf{Version (DE $\rightarrow$ FR)} : identifier sujet/verbe/compléments $\rightarrow$ traduire le sens, pas mot à mot $\rightarrow$ reformuler en français naturel.
  \item \textbf{Relecture} : verbe conjugué à la bonne place ? noms avec majuscule ? auxiliaire correct (\all{haben}/\all{sein}) ? article et cas corrects ?
\end{itemize}
\end{aretenir}

\begin{attention}[Pièges fréquents des francophones]
\begin{itemize}
  \item Ne dites pas \all{*Ich habe nach Yaoundé gefahren} : les verbes de mouvement prennent \all{sein} au Perfekt : \all{Ich \textbf{bin} nach Yaoundé gefahren.}
  \item Après un complément en tête de phrase, le verbe reste en 2\textsuperscript{e} position : \all{Gestern \textbf{hat} die Redaktion...} et non \all{*Gestern die Redaktion hat...}.
  \item \all{die Zeitung} signifie « le journal » (le support), pas « le temps » : attention au faux ami avec \all{die Zeit} (le temps).
  \item \all{aktualisieren} n'est pas « actualiser » au sens journalistique courant ; pour « faire le point sur l'actualité », on dit \all{über die aktuellen Ereignisse berichten}.
  \item Tous les noms prennent la majuscule : \all{das Fest, der Erfolg, die Schule}.
\end{itemize}
\end{attention}

\begin{aretenir}[Checklist avant le Bac]
\begin{itemize}
  \item $\square$ Je connais le lexique de la presse : \all{die Zeitung, der Artikel, die Schlagzeile, das Interview, die Anzeige}.
  \item $\square$ Je forme le Perfekt sans erreur : \all{haben}/\all{sein} + Partizip II en fin de phrase.
  \item $\square$ Je connais le Präteritum de \all{sein} (\all{war}), \all{haben} (\all{hatte}) et des modaux (\all{musste, konnte, wollte}).
  \item $\square$ Je place toujours le verbe conjugué en 2\textsuperscript{e} position, même après \all{Gestern}, \all{Heute}, \all{Am Montag}.
  \item $\square$ Je mets le verbe conjugué à la fin dans les subordonnées avec \all{weil, dass, als}.
  \item $\square$ Je respecte l'ordre TeKaMoLo pour les compléments.
  \item $\square$ Je sais structurer un article : titre, introduction (Wer? Was? Wann? Wo?), développement, conclusion.
  \item $\square$ Je sais rédiger une interview avec des verbes introducteurs variés et l'inversion du sujet.
  \item $\square$ Je mets une majuscule à tous les noms allemands.
  \item $\square$ Je traduis le sens et non mot à mot, puis je reformule dans la langue cible.
  \item $\square$ Je relis ma production : cas, articles, orthographe, Umlaute (\all{ä, ö, ü, ß}).
  \item $\square$ Je gère mon temps : je laisse 5 minutes pour la relecture finale.
\end{itemize}
\end{aretenir}

\findecours

\end{document}
